% !TeX program = pdflatex
\documentclass[11pt,a4paper]{article}
\usepackage[T1]{fontenc}
\usepackage[utf8]{inputenc}
\usepackage[polish]{babel}
\usepackage{lmodern}
\usepackage{amsmath,amssymb,amsthm,mathtools}
\usepackage{microtype}
\usepackage{xcolor}
\usepackage{geometry}
\usepackage{enumitem}
\usepackage{fancyhdr}
\usepackage[unicode]{hyperref}
\hypersetup{colorlinks=true,linkcolor=blue!45!black,urlcolor=blue!45!black}
\geometry{margin=2.5cm}
\setlength{\parindent}{0pt}
\setlength{\parskip}{4pt}
\setlength{\emergencystretch}{2em}
\setlist[enumerate]{topsep=4pt,itemsep=2pt}
\setlist[itemize]{topsep=4pt,itemsep=2pt}
\numberwithin{equation}{section}
\theoremstyle{plain}
\newtheorem{twierdzenie}{Twierdzenie}[section]
\newtheorem{lemat}[twierdzenie]{Lemat}
\newtheorem{stwierdzenie}[twierdzenie]{Stwierdzenie}
\newtheorem{wniosek}[twierdzenie]{Wniosek}
\theoremstyle{definition}
\newtheorem{definicja}[twierdzenie]{Definicja}
\newtheorem{przyklad}[twierdzenie]{Przykład}
\theoremstyle{remark}
\newtheorem{uwaga}[twierdzenie]{Uwaga}
\renewcommand{\proofname}{Dowód}
\pagestyle{fancy}
\fancyhf{}
\fancyhead[L]{\small Rozszerzenia pojęcia całki}
\fancyhead[R]{\small\ifnum\value{section}>0 Sekcja \thesection\else Wprowadzenie\fi}
\fancyfoot[C]{\thepage}
\renewcommand{\headrulewidth}{0.4pt}
\newcommand{\R}{\mathbb R}
\newcommand{\N}{\mathbb N}
\newcommand{\C}{\mathbb C}
\newcommand{\Q}{\mathbb Q}
\newcommand{\La}{\lambda}
\newcommand{\ind}[1]{\mathbf 1_{#1}}
\newcommand{\osc}{\operatorname{osc}}
\newcommand{\Var}{\operatorname{Var}}
\newcommand{\mesh}{\operatorname{mesh}}
\newcommand{\supp}{\operatorname{supp}}
\newcommand{\AC}{\mathrm{AC}}
\newcommand{\BV}{\mathrm{BV}}
\newcommand{\HK}{\mathrm{HK}}
\newcommand{\RS}{\mathrm{RS}}
\newcommand{\KS}{\mathrm{KS}}
\newcommand{\Dn}{\underline D}
\newcommand{\Du}{\overline D}
\newcommand{\pv}{\mathrm{p.w.}}
\DeclareMathOperator{\esssup}{ess\,sup}
\title{Rozszerzenia pojęcia całki\\\large Definicje, kryteria, dowody i przykłady}
\author{}
\date{}

% BEGIN GLOBAL CODE LISTINGS (generated from latex/pl/code-listings.tex)
\usepackage{xcolor}
\usepackage{listings}
\definecolor{codebg}{RGB}{248,249,251}
\definecolor{codeborder}{RGB}{180,191,203}
\definecolor{codecomment}{RGB}{61,113,79}
\definecolor{codekeyword}{RGB}{41,78,132}
\lstdefinelanguage{Julia}{
  morekeywords={abstract,baremodule,begin,break,catch,const,continue,do,else,elseif,
    end,export,false,finally,for,function,global,if,import,in,isa,let,local,macro,
    module,mutable,primitive,quote,return,struct,true,try,type,using,where,while,
    missing,nothing,println,print},
  sensitive=true,morecomment=[l]{\#},morecomment=[n]{\#=}{=\#},
  morestring=[b]",
}
\lstdefinestyle{skrypt}{
  backgroundcolor=\color{codebg},frame=single,rulecolor=\color{codeborder},
  framesep=5pt,frameround=tttt,framerule=.5pt,
  basicstyle=\ttfamily\footnotesize,keywordstyle=\color{codekeyword}\bfseries,
  commentstyle=\color{codecomment}\itshape,stringstyle=\color{purple!65!black},
  numbers=left,numberstyle=\tiny\color{gray!75!black},numbersep=8pt,
  stepnumber=1,breaklines=true,breakatwhitespace=false,
  columns=fullflexible,keepspaces=true,showstringspaces=false,tabsize=4,
  linewidth=.90\linewidth,captionpos=b,aboveskip=11pt,belowskip=12pt,
  xleftmargin=12pt,xrightmargin=2pt,
}
\lstdefinestyle{pythonpde}{style=skrypt,language=Python}
\lstdefinestyle{python}{style=skrypt,language=Python}
\lstdefinestyle{julia}{style=skrypt,language=Julia}
\lstset{style=skrypt}
\renewcommand{\lstlistingname}{Kod}
\renewcommand{\lstlistlistingname}{Spis kodów}
% END GLOBAL CODE LISTINGS
\begin{document}
\maketitle
\begin{abstract}
Przedstawiamy całki Darboux, Riemanna i Riemanna--Stieltjesa, całki Henstocka--Kurzweila, Perrona, Kurzweila--Stieltjesa i McShane'a, wąską i szeroką całkę Denjoya oraz ich związki z teorią miary. Omawiamy ponadto całkę Lebesgue'a, konstrukcję Daniella--Stone'a, twierdzenie Riesza--Markowa, całkę Radona i miarę Haara. Tekst obejmuje \emph{część I}: kończy się wprowadzeniem do całki Mawhina. Uzupełniono dowody, rachunki wzorcowe, kryteria zbieżności i porównania klas całkowalności; fragmenty nieuwzględnione w skanie oznaczono jako autorskie rozszerzenia. Dalsze tematy wymienione w pierwotnym spisie treści pozostają zapowiedziane, lecz nie są tutaj opracowywane.
\end{abstract}
\tableofcontents

\section*{Wstęp i konwencje}
\addcontentsline{toc}{section}{Wstęp i konwencje}
W sekcjach dotyczących jednej zmiennej $[a,b]\subset\R$ jest ustalonym przedziałem zwartym, $a<b$, a $\lambda$ oznacza miarę Lebesgue'a. Zakładamy podstawy teorii miary, algebry liniowej oraz analizy funkcjonalnej. Podstawy grup topologicznych i lokalnej zwartości wprowadzamy przed konstrukcją miary Haara.

\textbf{Co porównujemy.} Całkowalność zależy od funkcji, dziedziny, miary lub funkcji całkującej oraz przestrzeni wartości. Zmiana któregokolwiek z tych elementów może zmienić odpowiedź. O ile nie zaznaczamy inaczej, całki przyjmują wartości skończone. Dla nieujemnych funkcji mierzalnych dopuszczamy również $+\infty$. O ile nie zaznaczamy inaczej, funkcje są rzeczywiste. Dla funkcji zespolonej $f=u+iv$ przyjmujemy, że jest całkowalna w danej teorii, jeżeli całkowalne są obie funkcje rzeczywiste $u,v$; wówczas jej całka jest sumą całki z $u$ oraz $i$ razy całki z $v$. Symbolem $\chi_A$ oznaczamy funkcję charakterystyczną zbioru $A$, równą $1$ na $A$ i $0$ poza $A$.

Dla przestrzeni z miarą $(X,\mathcal A,\mu)$ przez $\mathcal L^1(\mu)$ oznaczamy zbiór konkretnych funkcji $\mathcal A$-mierzalnych $f:X\to\R$ o skończonych wartościach i $\int |f|\,\mathrm d\mu<\infty$, a przez $L^1(\mu)$ jego przestrzeń ilorazową względem równości prawie wszędzie. Jest to istotne przy porównywaniu Riemanna i Lebesgue'a: $\chi_{\Q\cap[0,1]}$ nie jest całkowalna w sensie Riemanna (zob. przykład funkcji Dirichleta w sekcji o całce Darboux oraz twierdzenie o równoważności całek Darboux i Riemanna), chociaż reprezentuje w $L^1(\lambda)$ tę samą klasę co funkcja zerowa: zbiór $\Q\cap[0,1]$ jest przeliczalny, a każdy jego singleton ma miarę Lebesgue'a równą zeru, więc z przeliczalnej podaddytywności miary cały zbiór ma miarę zero.

\textbf{Terminologia.} \emph{Integrand} to funkcja podcałkowa; \emph{integrator} to funkcja całkująca (w teorii Stieltjesa). \emph{Antiderivative} oznacza funkcję pierwotną, zaś \emph{primitive} w teorii całki HK oznacza funkcję pierwotną w sensie tej teorii, tj. funkcję postaci $F(x)=C+\int_a^x f(t)\,\mathrm dt$ dla funkcji $f$ całkowalnej w sensie HK i stałej $C\in\R$. Jej ciągłość oraz równość $F'=f$ prawie wszędzie są własnościami tej całki, a nie warunkami definiującymi funkcję pierwotną. \emph{Gauge} oznacza cechowanie, \emph{tagged partition} --- podział znaczony, \emph{refinement} --- zagęszczenie. \emph{Bounded variation} to wahanie ograniczone, a \emph{regulated function} --- funkcja regulowana (mająca skończone granice jednostronne). \emph{Finite-valued} oznacza funkcję o skończonych wartościach, niekoniecznie ograniczoną.

Podziałem przedziału $[a,b]$ nazywamy zbiór $P=\{a=x_0<x_1<\dots<x_n=b\}$. Przyjmujemy $n\ge1$; jego podprzedziałami są domknięte przedziały $[x_{i-1},x_i]$ dla $1\le i\le n$. Ich wnętrza są parami rozłączne, a sąsiednie podprzedziały mają wspólny koniec. Średnicą podziału nazywamy $\mesh P=\max_{1\le i\le n}(x_i-x_{i-1})$. Po wyborze znaczników $t_i\in[x_{i-1},x_i]$ powstaje podział znaczony $\mathcal P=\{(t_i,[x_{i-1},x_i])\}_{i=1}^n$, o sumie
\[
S(f,\mathcal P)=\sum_{i=1}^n f(t_i)(x_i-x_{i-1}).
\]
W sumach Stieltjesa długość przedziału zastępujemy przyrostem funkcji całkującej.


\section{Całka Darboux}
% Import: Analiza matematyczna I, rozdział 29 i rozdział 19 (zachowane etykiety).
\begin{definicja}
\label{ma1:def:riemann-partition}
Niech $a<b$. Podziałem przedziału $[a,b]$ nazywamy skończony ciąg $P=(x_0,\ldots,x_n)$, gdzie $n\ge1$ oraz $a=x_0<x_1<\cdots<x_n=b$. Oznaczamy $\Delta x_i=x_i-x_{i-1}$ oraz
\[
 |P|=\max_{1\le i\le n}\Delta x_i.
\]
Liczbę $|P|$ (oznaczaną też przez $\mesh P$) nazywamy średnicą podziału. Podział $Q$ jest rozdrobnieniem $P$, jeśli zawiera wszystkie jego punkty. Suma zbiorów punktów dwóch podziałów, uporządkowana rosnąco, jest ich wspólnym rozdrobnieniem.
\end{definicja}
\begin{definicja}
\label{ma1:def:darboux-sums}
Niech $f:[a,b]\to\R$ będzie ograniczona. Dla $I_i=[x_{i-1},x_i]$ definiujemy
\[
 m_i=\inf_{I_i}f,\qquad M_i=\sup_{I_i}f,\qquad
 L(f,P)=\sum_{i=1}^n m_i\Delta x_i,\qquad
 U(f,P)=\sum_{i=1}^n M_i\Delta x_i.
\]
Są to odpowiednio suma dolna i suma górna Darboux. Ich różnica wynosi $\sum_i\omega_i\Delta x_i$, gdzie $\omega_i=M_i-m_i$ jest oscylacją $f$ na $I_i$.
\end{definicja}
Kresy istnieją dzięki ograniczoności $f$ i zupełności $\R$. Nie muszą być osiągane. Na przykład funkcja $f(x)=x$ dla $0<x<1$, z wartościami $f(0)=f(1)=1/2$, ma na $[0,1]$ kres dolny $0$ i górny $1$: wartości $f(1/n)$ i $f(1-1/n)$ dla $n\ge2$ dążą odpowiednio do $0$ i $1$, lecz wszystkie wartości funkcji należą do $(0,1)$. Jeżeli $f$ jest nieujemna, sumy można interpretować jako pola prostokątów zawartych pod wykresem i pokrywających obszar pod wykresem. Dla funkcji zmieniających znak są to sumy ze znakiem.

\begin{center}
\begin{tikzpicture}[x=5.2cm,y=3.2cm,>=Stealth]
 \draw[->] (-.07,0)--(1.12,0) node[right] {$x$};
 \draw[->] (0,-.07)--(0,1.15) node[above] {$y$};
 \foreach \i in {0,1,2,3}{
   \pgfmathsetmacro{\l}{\i/4}
   \pgfmathsetmacro{\r}{(\i+1)/4}
   \pgfmathsetmacro{\lo}{\l*\l}
   \pgfmathsetmacro{\hi}{\r*\r}
   \fill[blue!18] (\l,0) rectangle (\r,\lo);
   \fill[orange!30] (\l,\lo) rectangle (\r,\hi);
   \draw[orange!80!black] (\l,0) rectangle (\r,\hi);
   \draw[blue!65!black] (\l,\lo)--(\r,\lo);
 }
 \draw[very thick,black,domain=0:1,samples=100] plot (\x,{\x*\x});
 \foreach \i in {1,2,3,4}{
   \pgfmathsetmacro{\t}{\i/4}
   \draw (\t,-.018)--(\t,.018);
   \node[below] at (\t,-.02) {$x_{\i}$};
 }
 \node[below left] at (0,0) {$x_0$};
 \node[above left] at (.88,.83) {$f(x)=x^2$};
\end{tikzpicture}

\small Rysunek: suma dolna (niebieska) i suma górna (niebieska wraz z pomarańczową częścią) dla $f(x)=x^2$ i czterech równych podprzedziałów $[0,1]$. Pomarańczowa część przedstawia różnicę sum.
\end{center}

\begin{lemat}
\label{ma1:lem:darboux-refinement}
Jeżeli $Q$ jest rozdrobnieniem $P$, to
\[
 L(f,P)\le L(f,Q)\le U(f,Q)\le U(f,P).
\]
Ponadto dla dowolnych podziałów $P,R$ zachodzi $L(f,P)\le U(f,R)$.
\end{lemat}
\begin{proof}
Niech $f$ będzie ograniczona na $[a,b]$. Każdy podprzedział $P$ jest sumą kolejnych podprzedziałów $Q$. Kres dolny na każdym mniejszym przedziale jest nie mniejszy niż kres dolny na większym, a kres górny jest nie większy niż kres górny na większym. Mnożymy te nierówności przez dodatnie długości mniejszych przedziałów i sumujemy; suma ich długości jest długością większego przedziału. Otrzymujemy zewnętrzne nierówności. Środkowa wynika z $m_i\le M_i$. Dla dowolnych $P,R$ bierzemy wspólne rozdrobnienie $Q$; wtedy $L(f,P)\le L(f,Q)\le U(f,Q)\le U(f,R)$.
\end{proof}
\begin{definicja}
\label{ma1:def:riemann-integral}
Dla ograniczonej funkcji $f:[a,b]\to\R$ definiujemy całkę dolną i górną, biorąc kresy po wszystkich podziałach $[a,b]$:
\[
 \underline{\mathcal I}(f)=\sup_P L(f,P),\qquad
 \overline{\mathcal I}(f)=\inf_P U(f,P).
\]
Jeżeli te liczby są równe, mówimy, że $f$ jest całkowalna w sensie Darboux na $[a,b]$, a wspólną wartość oznaczamy przez $\int_a^b f(x)\,\dd x$.
\end{definicja}
Jeśli $m\le f\le M$, to wszystkie sumy leżą w $[m(b-a),M(b-a)]$. Oba kresy są więc skończone. Z lematu \ref{ma1:lem:darboux-refinement}, najpierw biorąc supremum po $P$, a następnie infimum po $R$, dostajemy $\underline{\mathcal I}(f)\le\overline{\mathcal I}(f)$.

\begin{twierdzenie}[Kryterium Darboux]
\label{ma1:thm:darboux-criterion}
Funkcja ograniczona $f$ jest całkowalna w sensie Darboux na $[a,b]$ wtedy i tylko wtedy, gdy dla każdego $\varepsilon>0$ istnieje podział $P$ taki, że $U(f,P)-L(f,P)<\varepsilon$.
\end{twierdzenie}
\begin{proof}
Jeżeli obie całki są równe $A$, z definicji kresów wybieramy $P,R$ takie, że $L(f,P)>A-\varepsilon/2$ i $U(f,R)<A+\varepsilon/2$. Dla wspólnego rozdrobnienia $Q$ mamy $U(f,Q)-L(f,Q)\le U(f,R)-L(f,P)<\varepsilon$. Odwrotnie, dla każdego podziału $P$ mamy
\[
 0\le\overline{\mathcal I}(f)-\underline{\mathcal I}(f)
 \le U(f,P)-L(f,P).
\]
Jeżeli prawa strona może być mniejsza od dowolnego $\varepsilon>0$, różnica całek jest zerowa.
\end{proof}

\subsection{Kryteria całkowalności i przykłady}
\begin{twierdzenie}
\label{ma1:thm:riemann-continuous}
Każda funkcja ciągła $f:[a,b]\to\R$ jest całkowalna.
\end{twierdzenie}
\begin{proof}
Funkcja jest ograniczona z twierdzenia Weierstrassa i jednostajnie ciągła z twierdzenia Heinego--Cantora. Dla $\varepsilon>0$ istnieje $\delta>0$ takie, że $|f(s)-f(t)|<\varepsilon/(2(b-a))$ dla $|s-t|<\delta$. Dla podziału o $|P|<\delta$ oscylacja na każdym $I_i$ jest nie większa niż $\varepsilon/(2(b-a))$ (możemy wziąć punkty osiągania maksimum i minimum). Stąd $U(f,P)-L(f,P)\le\varepsilon/2<\varepsilon$. Stosujemy kryterium Darboux.
\end{proof}
\begin{twierdzenie}
\label{ma1:thm:riemann-monotone}
Każda funkcja monotoniczna $f:[a,b]\to\R$ jest całkowalna.
\end{twierdzenie}
\begin{proof}
Jeśli $f$ jest niemalejąca, to $m_i=f(x_{i-1})$, $M_i=f(x_i)$, a wartości funkcji leżą pomiędzy $f(a)$ i $f(b)$. Zatem
\[
 U(f,P)-L(f,P)\le |P|\sum_i(f(x_i)-f(x_{i-1}))
 =|P|(f(b)-f(a)).
\]
Podziały o dowolnie małej średnicy dają kryterium Darboux. Dla funkcji nierosnącej mamy $m_i=f(x_i)$ i $M_i=f(x_{i-1})$; ten sam rachunek daje oszacowanie $|P|(f(a)-f(b))$.
\end{proof}
\begin{twierdzenie}
\label{ma1:thm:riemann-finite-discontinuities}
Funkcja ograniczona na $[a,b]$ i ciągła poza skończonym zbiorem punktów jest całkowalna.
\end{twierdzenie}
\begin{proof}
Niech $|f|\le B$ dla $B>0$, a zbiór możliwych nieciągłości ma $k$ punktów. Gdy $k=0$, stosujemy twierdzenie \ref{ma1:thm:riemann-continuous}. Gdy $k>0$, oznaczamy te punkty przez $d_1,\ldots,d_k$. Wybieramy $\rho>0$ mniejsze od $(b-a)/2$ oraz $\varepsilon/(8Bk)$, a gdy $k\ge2$, także od połowy najmniejszej odległości między różnymi $d_j$. Dla każdego $j$ bierzemy
\[
 J_j=[\max(a,d_j-\rho),\min(b,d_j+\rho)].
\]
Każdy $J_j$ ma dodatnią długość, zawiera $d_j$ we wnętrzu względem $[a,b]$ i nie przecina żadnego innego $J_i$. Suma długości jest nie większa niż $2k\rho<\varepsilon/(4B)$. Włączamy do podziału $a,b$ oraz wszystkie końce $J_j$. Na nich łączny wkład do różnicy sum jest mniejszy od $2B\varepsilon/(4B)=\varepsilon/2$.

Pozostałe odcinki o dodatniej długości pomiędzy tymi przedziałami są domknięte i nie zawierają punktów nieciągłości: końce wewnętrzne wybranych przedziałów nie są tymi punktami. Na każdym z tych odcinków funkcja jest ciągła. Jest ich skończenie wiele, powiedzmy $\ell$. Jeżeli $\ell>0$, wybieramy na każdym podział z różnicą sum mniejszą od $\varepsilon/(2\ell)$ na mocy twierdzenia \ref{ma1:thm:riemann-continuous} i kryterium Darboux. Łączymy wszystkie podziały. Wkłady na pozostałych odcinkach sumują się do mniej niż $\varepsilon/2$; gdy $\ell=0$, ich wkład wynosi zero. Cała różnica sum jest mniejsza od $\varepsilon$.
\end{proof}
\begin{przyklad}
\label{ma1:ex:riemann-dirichlet}
Funkcja $\chi_{\Q\cap[a,b]}$ nie jest całkowalna, gdy $a<b$. Każdy podprzedział dodatniej długości zawiera liczbę wymierną oraz niewymierną, więc jego kresy wynoszą $0$ i $1$. Dla przypomnienia: między $u<v$ wybieramy liczbę całkowitą $n$ z $n(v-u)>1$, a następnie całkowite $m$ z $nu<m<nv$; wtedy $m/n\in(u,v)$. Wybierając wymierną w $(u-\sqrt2,v-\sqrt2)$ i dodając $\sqrt2$, otrzymujemy liczbę niewymierną w $(u,v)$. Stąd dla każdego podziału $L=0$, $U=b-a$, co przeczy kryterium Darboux. 
\end{przyklad}

\begin{przyklad}[Funkcja Thomae'a]
\label{ma1:ex:thomae}
Definiujemy $T:\R\to[0,1]$ przez
\[
 T(x)=
 \begin{cases}
 1/q,& x=p/q,\ p\in\Z,\ q\in\N,\ q\ge1,\ \gcd(|p|,q)=1,\\
 0,&x\notin\Q.
 \end{cases}
\]
Ułamek nieskracalny o dodatnim mianowniku jest jednoznaczny; w szczególności $T(0)=1$. Dla każdego $a\in\R$ istnieje granica przekłuta $\lim_{x\to a}T(x)=0$. Zatem $T$ jest ciągła dokładnie w punktach niewymiernych.
\end{przyklad}
\begin{proof}
Ustalmy $a\in\R$ i $\varepsilon>0$. Wybierzmy całkowite $N\ge1$ z $1/N<\varepsilon$. Zbiór
$E_N=\{p/q\in[a-1,a+1]:p\in\Z,\ 1\le q\le N,\ \gcd(|p|,q)=1\}$
jest skończony: dla każdego spośród skończenie wielu mianowników $q$ licznik należy do ograniczonego przedziału $[q(a-1),q(a+1)]$. Przyjmijmy $\delta$ jako minimum liczby $1$ i wszystkich dodatnich odległości $|r-a|$ dla $r\in E_N\setminus\{a\}$. Jeśli ten ostatni zbiór jest pusty, przyjmujemy $\delta=1$. Mamy $\delta>0$.

Dla $0<|x-a|<\delta$ punkt $x$ należy do $[a-1,a+1]$, ale nie do $E_N$. Jeśli $x$ jest niewymierny, to $T(x)=0$; jeśli jest wymierny, jego nieskracalny mianownik jest większy od $N$, więc $0<T(x)<1/N<\varepsilon$. Dowodzi to granicy przekłutej. Ponieważ każdy punkt $\R$ jest punktem skupienia dziedziny, ciągłość jest równoważna $T(a)=0$, czyli $a\notin\Q$.
\end{proof}
\begin{center}
\begin{tikzpicture}[x=9cm,y=2.7cm,>=Stealth]
 \draw[->] (-.03,0)--(1.08,0) node[right] {$x$};
 \draw[->] (0,-.08)--(0,1.15) node[above] {$y$};
 \fill[blue!65!black] ({0/1},{1/1}) circle (.9pt);
 \fill[blue!65!black] ({1/1},{1/1}) circle (.9pt);
 \fill[blue!65!black] ({1/2},{1/2}) circle (.9pt);
 \fill[blue!65!black] ({1/3},{1/3}) circle (.9pt);
 \fill[blue!65!black] ({2/3},{1/3}) circle (.9pt);
 \fill[blue!65!black] ({1/4},{1/4}) circle (.9pt);
 \fill[blue!65!black] ({3/4},{1/4}) circle (.9pt);
 \fill[blue!65!black] ({1/5},{1/5}) circle (.9pt);
 \fill[blue!65!black] ({2/5},{1/5}) circle (.9pt);
 \fill[blue!65!black] ({3/5},{1/5}) circle (.9pt);
 \fill[blue!65!black] ({4/5},{1/5}) circle (.9pt);
 \fill[blue!65!black] ({1/6},{1/6}) circle (.9pt);
 \fill[blue!65!black] ({5/6},{1/6}) circle (.9pt);
 \fill[blue!65!black] ({1/7},{1/7}) circle (.9pt);
 \fill[blue!65!black] ({2/7},{1/7}) circle (.9pt);
 \fill[blue!65!black] ({3/7},{1/7}) circle (.9pt);
 \fill[blue!65!black] ({4/7},{1/7}) circle (.9pt);
 \fill[blue!65!black] ({5/7},{1/7}) circle (.9pt);
 \fill[blue!65!black] ({6/7},{1/7}) circle (.9pt);
 \fill[blue!65!black] ({1/8},{1/8}) circle (.9pt);
 \fill[blue!65!black] ({3/8},{1/8}) circle (.9pt);
 \fill[blue!65!black] ({5/8},{1/8}) circle (.9pt);
 \fill[blue!65!black] ({7/8},{1/8}) circle (.9pt);
 \fill[blue!65!black] ({1/9},{1/9}) circle (.9pt);
 \fill[blue!65!black] ({2/9},{1/9}) circle (.9pt);
 \fill[blue!65!black] ({4/9},{1/9}) circle (.9pt);
 \fill[blue!65!black] ({5/9},{1/9}) circle (.9pt);
 \fill[blue!65!black] ({7/9},{1/9}) circle (.9pt);
 \fill[blue!65!black] ({8/9},{1/9}) circle (.9pt);
 \fill[blue!65!black] ({1/10},{1/10}) circle (.9pt);
 \fill[blue!65!black] ({3/10},{1/10}) circle (.9pt);
 \fill[blue!65!black] ({7/10},{1/10}) circle (.9pt);
 \fill[blue!65!black] ({9/10},{1/10}) circle (.9pt);
 \fill[blue!65!black] ({1/11},{1/11}) circle (.9pt);
 \fill[blue!65!black] ({2/11},{1/11}) circle (.9pt);
 \fill[blue!65!black] ({3/11},{1/11}) circle (.9pt);
 \fill[blue!65!black] ({4/11},{1/11}) circle (.9pt);
 \fill[blue!65!black] ({5/11},{1/11}) circle (.9pt);
 \fill[blue!65!black] ({6/11},{1/11}) circle (.9pt);
 \fill[blue!65!black] ({7/11},{1/11}) circle (.9pt);
 \fill[blue!65!black] ({8/11},{1/11}) circle (.9pt);
 \fill[blue!65!black] ({9/11},{1/11}) circle (.9pt);
 \fill[blue!65!black] ({10/11},{1/11}) circle (.9pt);
 \node[below] at (0,0) {$0$};
 \node[below] at (1,0) {$1$};
 \node[left] at (0,1) {$1$};
 \node[blue!65!black,above] at (.5,.5) {$1/2$};
\end{tikzpicture}

\small Rysunek: wartości funkcji Thomae'a dla nieskracalnych mianowników nie większych niż $11$; pominięto wartości o większych mianownikach i wartości zerowe w punktach niewymiernych.
\end{center}

\begin{przyklad}[Całka Darboux funkcji Thomae'a]
Rozważamy ograniczenie funkcji z przykładu \ref{ma1:ex:thomae} do $[0,1]$.
Dla dowolnego podziału $P$ mamy $L(T,P)=0$, gdyż każdy podprzedział zawiera liczbę niewymierną, a $T\ge0$. Ustalmy $\varepsilon>0$ i wybierzmy $N$ z $1/N<\varepsilon/2$. Niech $E_N$ będzie zbiorem liczb wymiernych w $[0,1]$ o mianowniku nie większym niż $N$ w zapisie nieskracalnym; jest skończony z poprzedniego dowodu. Oznaczmy $k=|E_N|$; mamy $k\ge1$. Wybieramy podział $P$ o średnicy mniejszej niż $\varepsilon/(4k)$. Każdy punkt $E_N$ należy do co najwyżej dwóch domkniętych podprzedziałów, dlatego podprzedziałów przecinających $E_N$ jest co najwyżej $2k$, a suma ich długości nie przekracza $2k\mesh P$. Na nich $T\le1$, zaś na pozostałych $T\le1/N$. Otrzymujemy
\[
0\le U(T,P)-L(T,P)=U(T,P)
\le 2k\mesh P+\frac1N<\varepsilon.
\]
Z kryterium Darboux wynika całkowalność w sensie Darboux. Ponieważ każda suma dolna wynosi zero, wspólna wartość obu całek Darboux jest równa zeru.
\end{przyklad}
\begin{przyklad}[Obliczenie z definicji]
Dla $f(x)=x^2$ na $[0,1]$ funkcja jest niemalejąca, zatem na $[(i-1)/n,i/n]$ infimum wynosi $((i-1)/n)^2$, a supremum $(i/n)^2$. Korzystamy ze wzoru $\sum_{i=1}^n i^2=n(n+1)(2n+1)/6$: prawa strona dla $n=0$ wynosi zero, a jej przyrost od $n-1$ do $n$ to
\[
\frac{n}{6}\bigl((n+1)(2n+1)-(n-1)(2n-1)\bigr)=n^2,
\]
co dowodzi wzoru przez sumowanie przyrostów. Dla $P_n=\{i/n:0\le i\le n\}$ mamy
\[
U(f,P_n)=\frac1{n^3}\sum_{i=1}^n i^2=\frac{(n+1)(2n+1)}{6n^2},\qquad
L(f,P_n)=\frac1{n^3}\sum_{i=0}^{n-1}i^2=\frac{(n-1)(2n-1)}{6n^2}.
\]
Dzieląc liczniki przez $n^2$, otrzymujemy, że obie sumy dążą do $1/3$. Dla każdego $n$ suma dolna nie przekracza całki dolnej, ta nie przekracza całki górnej, a całka górna nie przekracza sumy górnej. Z twierdzenia o trzech ciągach (lub bezpośrednio z tych nierówności) obie całki Darboux wynoszą $1/3$, czyli $\int_0^1x^2\,\mathrm dx=1/3$.
\end{przyklad}

\section{Całka Riemanna}
\label{ma1:ch:riemann-integral}
% Import: Analiza matematyczna I, rozdział 29; Teoria miary i całki, rozdział 08.
\begin{definicja}
\label{ma1:def:riemann-sum}
Dla podziału $P$ i punktów wyróżnionych $\xi_i\in[x_{i-1},x_i]$ definiujemy sumę Riemanna
\[
 S(f,P,\xi)=\sum_{i=1}^n f(\xi_i)\Delta x_i.
\]
Granica sum Riemanna przy $|P|\to0$ wynosi $A$, jeżeli dla każdego $\varepsilon>0$ istnieje $\delta>0$ takie, że dla każdego podziału $P$ o $|P|<\delta$ i każdego wyboru $\xi_i$ mamy $|S(f,P,\xi)-A|<\varepsilon$.
\end{definicja}
Dla funkcji ograniczonej zawsze $L(f,P)\le S(f,P,\xi)\le U(f,P)$, ponieważ $m_i\le f(\xi_i)\le M_i$.

\begin{stwierdzenie}
Każda funkcja całkowalna w sensie Riemanna jest ograniczona.
\end{stwierdzenie}
\begin{proof}
Gdyby $f$ była nieograniczona, dla dowolnego ustalonego podziału musiałaby być nieograniczona na co najmniej jednej komórce. Zmieniając znacznik tylko na niej, uzyskalibyśmy sumy o dowolnie dużej różnicy, wbrew definicji całki.
\end{proof}

\begin{twierdzenie}
\label{ma1:thm:riemann-darboux-equivalence}
Dla ograniczonej $f$ całkowalność z definicji \ref{ma1:def:riemann-integral} jest równoważna istnieniu skończonej granicy sum Riemanna przy $|P|\to0$. W obu przypadkach otrzymujemy tę samą wartość.
\end{twierdzenie}
\begin{proof}
Załóżmy całkowalność i oznaczmy całkę przez $A$. Ustalmy $\varepsilon>0$ i podział $P_0$ z $U(f,P_0)-L(f,P_0)<\varepsilon/2$. Niech $k$ oznacza liczbę jego punktów wewnętrznych, a $B>0$ będzie ograniczeniem $|f|$.

Dla dowolnego podziału $P$ weźmy wspólne rozdrobnienie $Q$ podziałów $P,P_0$. Podprzedział $P$ zmienia się w tym rozdrobnieniu tylko wtedy, gdy zawiera w swoim wnętrzu punkt $P_0$. Takich podprzedziałów jest co najwyżej $k$, ponieważ ich wnętrza są parami rozłączne i w każdym znajduje się przynajmniej jeden punkt wewnętrzny $P_0$. Ich łączna długość jest nie większa niż $k|P|$. Punkty $P_0$ pokrywające się z końcami podprzedziałów $P$ nie powodują ich rozdrobnienia. Na każdym z nich zarówno stary wkład do sumy górnej, jak i suma nowych wkładów należą do przedziału $[-B\Delta x_i,B\Delta x_i]$. Różnica jest więc nie większa niż $2B\Delta x_i$. Analogicznie dla sumy dolnej. Z lematu \ref{ma1:lem:darboux-refinement} wynika
\[
 \begin{split}
 U(f,P)&\le U(f,Q)+2Bk|P|\le U(f,P_0)+2Bk|P|,\\
 L(f,P)&\ge L(f,Q)-2Bk|P|\ge L(f,P_0)-2Bk|P|.
 \end{split}
\]
Ponieważ $L(f,P_0)\le A\le U(f,P_0)$ oraz $L(f,P)\le S(f,P,\xi)\le U(f,P)$, mamy
\[
 -\bigl(U(f,P_0)-L(f,P_0)\bigr)-2Bk|P|
 \le S(f,P,\xi)-A
 \le U(f,P_0)-L(f,P_0)+2Bk|P|.
\]
Każda suma Riemanna dla $P$ spełnia zatem
\[
 |S(f,P,\xi)-A|<\varepsilon/2+2Bk|P|.
\]
Wybieramy $\delta=\varepsilon/(4B(k+1))$. Wtedy prawa strona jest mniejsza od $\varepsilon$ dla $|P|<\delta$, niezależnie od wyboru punktów wyróżnionych.

Odwrotnie, załóżmy warunek graniczny dla sum Riemanna z granicą $A$. Dla $\varepsilon>0$ wybieramy $\delta$ odpowiadające dokładności $\varepsilon/4$ i podział $P$ o średnicy mniejszej od $\delta$ (na przykład podział na $n$ równych części dla $n>(b-a)/\delta$). Z definicji kresów dla każdego $i$ można wybrać $\xi_i^+,\xi_i^-\in I_i$ tak, aby
\[
 M_i-f(\xi_i^+)<\frac{\varepsilon}{4(b-a)},\qquad
 f(\xi_i^-)-m_i<\frac{\varepsilon}{4(b-a)}.
\]
Po pomnożeniu przez długości i zsumowaniu otrzymujemy
$U(f,P)<S(f,P,\xi^+)+\varepsilon/4<A+\varepsilon/2$
oraz $L(f,P)>S(f,P,\xi^-)-\varepsilon/4>A-\varepsilon/2$.
Zatem $U(f,P)-L(f,P)<\varepsilon$. Kryterium Darboux daje całkowalność. Jeżeli $J$ oznacza uzyskaną całkę, to $L(f,P)\le J\le U(f,P)$, więc $|J-A|<\varepsilon/2$. Ponieważ $\varepsilon>0$ było dowolne, mamy $J=A$.
\end{proof}
W zapisie ze wstępu $S(f,P,\xi)=S(f,\mathcal P)$, gdzie $\mathcal P$ jest podziałem znaczonym, a $\xi_i$ są jego znacznikami. Istotna jest dowolność podziałów i punktów wyróżnionych. Zbieżność sum dla jednego szczególnego ciągu podziałów i jednego wyboru punktów nie wystarcza.

\subsection{Kryterium Lebesgue'a}
\begin{twierdzenie}[Kryterium Lebesgue'a całkowalności Riemanna]
\label{thm:riemann-criterion}
Niech $f:[a,b]\to\mathbb R$ będzie ograniczona, gdzie
$a<b$. Wówczas $f$ jest całkowalna w sensie Riemanna
tedy i tylko wtedy, gdy zbiór punktów nieciągłości $f$
ma miarę Lebesgue'a zero.
\end{twierdzenie}
\begin{proof}
Dla $x\in[a,b]$ oznaczmy przez $\omega_f(x)$
oscylację $f$ w punkcie, czyli
\[
\omega_f(x):=\inf_{r>0}
\left(\sup_{y\in[a,b]\cap(x-r,x+r)}f(y)
-\inf_{y\in[a,b]\cap(x-r,x+r)}f(y)\right).
\]
Funkcja $f$ jest ciągła w $x$ wtedy i tylko wtedy,
gdy $\omega_f(x)=0$. Dla $\varepsilon>0$ zbiór
$D_\varepsilon=\{x:\omega_f(x)\ge\varepsilon\}$
jest domknięty: jeśli $\omega_f(x)<\varepsilon$,
to na pewnym otoczeniu punktu $x$ oscylacja $f$
jest mniejsza niż $\varepsilon$, a zatem także
oscylacje w punktach pobliskich są mniejsze od
$\varepsilon$. Zatem $D_\varepsilon$ jest zwarty.

Załóżmy najpierw całkowalność Riemanna i
ustalmy $\varepsilon>0$. Weźmy dowolny podział
$P=(a=x_0<\cdots<x_N=b)$. Jeśli
$x\in D_\varepsilon$ nie jest punktem podziału,
to oscylacja $f$ na domkniętym przedziale
$[x_{j-1},x_j]$ zawierającym $x$ jest co najmniej
$\varepsilon$. Dlatego $D_\varepsilon$ mieści się
w sumie skończenie wielu punktów podziału i tych
przedziałów, których oscylacja jest nie mniejsza
niż $\varepsilon$. Suma długości tych przedziałów
jest nie większa niż
$\varepsilon^{-1}(U(f,P)-L(f,P))$.
Ponieważ dla funkcji całkowalnej w sensie Riemanna
różnicę tę można uczynić dowolnie małą,
$\lambda_1(D_\varepsilon)=0$.
Zbiór nieciągłości jest sumą
$\bigcup_{m\ge1}D_{1/m}$, więc jest zerowy.

Odwrotnie, niech zbiór nieciągłości $D$ ma miarę
zero, a $|f|\le M$. Ustalmy $\varepsilon>0$.
Zbiór $D_\varepsilon$ jest zwarty i zerowy, więc
istnieje zbiór otwarty $O$ zawierający $D_\varepsilon$
o dowolnie małej mierze Lebesgue'a (regularność zewnętrzna miary Lebesgue'a). Wybieramy przy tym
$\lambda_1(O)<b-a$, dzięki czemu zwarty zbiór
$K=[a,b]\setminus O$ jest niepusty. Dla każdego
$x\in K$ mamy $\omega_f(x)<\varepsilon$,
zatem istnieje promień $r_x>0$, na którego kuli
$B(x,r_x)\cap[a,b]$ oscylacja $f$ jest
mniejsza niż $\varepsilon$. Zwartość $K$ pozwala
wybrać skończenie wiele kul $B(x_j,r_{x_j}/2)$
pokrywających $K$. Połóżmy
$\delta=\frac12\min_j r_{x_j}>0$.
Jeżeli domknięty przedział $J\subseteq[a,b]$
o długości mniejszej niż $\delta$ przecina $K$,
wybieramy $y\in J\cap K$ i indeks $j$ taki, że
$y\in B(x_j,r_{x_j}/2)$. Dla każdego $z\in J$
mamy
$|z-x_j|\le|z-y|+|y-x_j|<\delta+r_{x_j}/2
\le r_{x_j}$.
Zatem $J\subseteq B(x_j,r_{x_j})$, gdzie
oscylacja $f$ jest mniejsza niż $\varepsilon$.

Weźmy podział $P$ o średnicy mniejszej niż $\delta$.
Na przedziałach przecinających $K$ oscylacja $f$
jest mniejsza niż $\varepsilon$. Pozostałe
przedziały (z wyjątkiem ewentualnych końców)
mieszczą się w $O$, więc suma ich długości
nie przekracza $\lambda_1(O)$. Otrzymujemy
\[
U(f,P)-L(f,P)
\le\varepsilon(b-a)+2M\lambda_1(O).
\]
Wybierając najpierw $\varepsilon$, a potem
$\lambda_1(O)$ dostatecznie małe, czynimy prawą
stronę dowolnie małą. Kryterium Darboux dowodzi
całkowalności Riemanna.
\end{proof}

\subsection{Własności całki oznaczonej}
\begin{twierdzenie}
\label{ma1:thm:riemann-algebra}
Jeżeli $f,g$ są całkowalne na $[a,b]$ oraz $\alpha,\beta\in\R$, to $\alpha f+\beta g$, $fg$ oraz $|f|$ są całkowalne. Ponadto
\[
 \int_a^b(\alpha f+\beta g)\,\dd x
 =\alpha\int_a^b f\,\dd x+\beta\int_a^b g\,\dd x.
\]
\end{twierdzenie}
\begin{proof}
Dla każdego podziału z punktami wyróżnionymi suma Riemanna funkcji $\alpha f+\beta g$ jest sumą odpowiednich sum pomnożonych przez skalary. Jeżeli $A,C$ oznaczają całki $f,g$, to dla zadanej $\varepsilon>0$ wybieramy wspólną średnicę graniczną dla dokładności $\varepsilon/(1+|\alpha|+|\beta|)$ obu sum. Nierówność trójkąta daje błąd mniejszy od $\varepsilon$ względem $\alpha A+\beta C$. Twierdzenie \ref{ma1:thm:riemann-darboux-equivalence} dowodzi całkowalności i wzoru.

Niech $|f|,|g|\le B$ dla $B>0$. Dla dowolnych $s,t$ w jednym podprzedziale
\[
 |f(s)g(s)-f(t)g(t)|\le B|g(s)-g(t)|+B|f(s)-f(t)|.
\]
Oscylacja ograniczonej funkcji $h$ na niepustym zbiorze $E$ jest równa $\sup_{s,t\in E}|h(s)-h(t)|$. Istotnie, dla $m=\inf_E h$ i $M=\sup_E h$ zachodzi $|h(s)-h(t)|\le M-m$. Dla dowolnego $\eta>0$ wybieramy $s,t\in E$ z $h(s)>M-\eta/2$ i $h(t)<m+\eta/2$. Wówczas $|h(s)-h(t)|\ge h(s)-h(t)>M-m-\eta$. Dowolność $\eta$ daje przeciwną nierówność. Zatem oscylacja iloczynu na podprzedziale jest nie większa niż $B$ razy suma oscylacji czynników. Wybieramy wspólne rozdrobnienie podziałów, dla których różnice sum $f$ i $g$ są mniejsze od $\varepsilon/(2B)$. Na wspólnym rozdrobnieniu $Q$ różnice nadal spełniają te oszacowania, zatem
\[
 U(fg,Q)-L(fg,Q)
 \le B\bigl(U(f,Q)-L(f,Q)+U(g,Q)-L(g,Q)\bigr)
 <\varepsilon.
\]
Iloczyn jest ograniczony przez $B^2$, więc stosujemy kryterium Darboux. Dla modułu korzystamy z $||f(s)|-|f(t)||\le |f(s)-f(t)|$, co daje nie większą oscylację niż dla $f$. Kryterium Darboux kończy dowód.
\end{proof}

\begin{twierdzenie}
\label{ma1:thm:riemann-order}
Dla całkowalnych $f,g$ na $[a,b]$ z $f\le g$ wynika $\int_a^b f\,\dd x\le\int_a^b g\,\dd x$. Jeżeli $m\le f\le M$, to
\[
 m(b-a)\le\int_a^b f\,\dd x\le M(b-a),\qquad
 \left|\int_a^b f\,\dd x\right|\le\int_a^b |f|\,\dd x.
\]
\end{twierdzenie}
\begin{proof}
Dla każdego podziału $S(f,P,\xi)\le S(g,P,\xi)$, więc przejście do granicy zachowuje nieostrą nierówność. Całka funkcji stałej $c$ wynosi $c(b-a)$, gdyż każda jej suma Riemanna jest równa tej liczbie. Porównujemy $f$ z funkcjami stałymi $m,M$. Nierówności $-|f|\le f\le |f|$ wraz z liniowością i całkowalnością modułu dają ostatnią tezę.
\end{proof}

\begin{twierdzenie}[Addytywność względem przedziału]
\label{ma1:thm:riemann-additivity}
Niech $a<c<b$ i niech $f$ będzie ograniczona na $[a,b]$. Jest ona całkowalna na $[a,b]$ wtedy i tylko wtedy, gdy jej ograniczenia do $[a,c]$ i $[c,b]$ są całkowalne. Wówczas
\[
 \int_a^b f\,\dd x=\int_a^c f\,\dd x+\int_c^b f\,\dd x.
\]
\end{twierdzenie}
\begin{proof}
Jeżeli $f$ jest całkowalna na całym przedziale, dla $\varepsilon>0$ wybieramy podział z różnicą sum mniejszą od $\varepsilon$ i dodajemy do niego punkt $c$. Różnica nie wzrasta. Jest ona sumą dwóch nieujemnych różnic dla podziałów ograniczonych do obu przedziałów, więc każda z nich jest mniejsza od $\varepsilon$. Kryterium Darboux daje całkowalność ograniczeń.

Odwrotnie, dla obu ograniczeń wybieramy podziały z różnicą sum mniejszą od $\varepsilon/2$ i łączymy je. Różnica sum na całym przedziale jest sumą obu różnic, więc jest mniejsza od $\varepsilon$. Aby wyznaczyć wartość całki, oznaczmy całki na obu częściach przez $A_1,A_2$. Dla połączonego podziału $P$ liczba $A_1+A_2$ leży pomiędzy $L(f,P)$ i $U(f,P)$, podobnie jak całka na całym przedziale. Moduł różnicy tych dwóch liczb jest mniejszy od $\varepsilon$. Dowolność $\varepsilon$ daje wzór.
\end{proof}
Przyjmujemy $\int_a^a f\,\dd x=0$ oraz $\int_b^a f\,\dd x=-\int_a^b f\,\dd x$ dla $a<b$. Wówczas addytywność zachodzi dla dowolnego porządku trzech końców należących do przedziału całkowalności. Istotnie, dla trzech różnych punktów $r<s<t$ oznaczmy $A=\int_r^s f\,\dd x$ i $C=\int_s^t f\,\dd x$. Wówczas całki w kierunku rosnącym mają wartości $A,C,A+C$, a w kierunku malejącym ich przeciwności. Przypisując punktom $r,s,t$ odpowiednio liczby $0,A,A+C$, każdą całkę między tymi punktami otrzymujemy jako różnicę liczby przypisanej końcowi i początkowi. Takie różnice spełniają addytywność w dowolnej kolejności. Jeżeli dwa punkty są równe, teza wynika z całki zerowej lub zmiany orientacji. Z twierdzenia \ref{ma1:thm:riemann-additivity} wynika też całkowalność na każdym domkniętym podprzedziale.

\begin{wniosek}
\label{ma1:cor:riemann-finite-changes}
Zmiana wartości funkcji całkowalnej w skończonej liczbie punktów nie zmienia jej całkowalności ani wartości całki.
\end{wniosek}
\begin{proof}
Niech $g$ różni się od $f$ jedynie w $k$ punktach. Różnica $h=g-f$ jest ograniczona, powiedzmy $|h|\le B$, oraz zerowa poza tymi punktami. Dla dowolnego podziału co najwyżej $2k$ jego podprzedziałów zawiera któryś z tych punktów (punkt może być wspólnym końcem dwóch). Zatem $|S(h,P,\xi)|\le 2kB|P|$. Sumy dążą do zera przy $|P|\to0$, więc $h$ jest całkowalna i jej całka wynosi zero. Liniowość dla $g=f+h$ daje tezę.
\end{proof}

\subsection{Twierdzenie o wartości średniej i przykłady}
\begin{twierdzenie}[O wartości średniej dla całek]
\label{ma1:thm:riemann-mean-value}
Niech $f$ będzie ciągła na $[a,b]$ i niech $g$ będzie całkowalna oraz nieujemna. Istnieje $\xi\in[a,b]$ takie, że
\[
 \int_a^b fg\,\dd x=f(\xi)\int_a^b g\,\dd x.
\]
W szczególności $\int_a^b f\,\dd x=f(\xi)(b-a)$ dla pewnego $\xi\in[a,b]$.
\end{twierdzenie}
\begin{proof}
Niech $m,M$ będą minimum i maksimum $f$, istniejące z twierdzenia Weierstrassa. Z $g\ge0$ wynika $mg\le fg\le Mg$. Całkowalność iloczynu i porównanie całek dają
$m\int g\le\int fg\le M\int g$, gdzie wszystkie całki są po $[a,b]$.
Jeżeli $\int g=0$, to również $\int fg=0$, więc dowolne $\xi$ spełnia wzór. Jeżeli $\int g>0$, dzielimy przez tę liczbę i otrzymujemy iloraz leżący w $[m,M]$. Twierdzenie o wartości pośredniej daje punkt $\xi$ o tej wartości $f$. Wybór $g=1$ daje szczególny przypadek.
\end{proof}
\begin{przyklad}
\label{ma1:ex:riemann-square}
Obliczymy $\int_0^1 x^2\,\dd x$ bez użycia funkcji pierwotnej. Dla podziału na $n$ równych części i prawych końców jako punktów wyróżnionych otrzymujemy
\[
 S_n=\frac1{n^3}\sum_{i=1}^n i^2
 =\frac{n(n+1)(2n+1)}{6n^3}.
\]
Tożsamość dla sumy kwadratów wynika z indukcji: dla $n=1$ obie strony wynoszą $1$, a krok indukcyjny ma postać
\[
 \frac{n(n+1)(2n+1)}6+(n+1)^2
 =\frac{(n+1)(2n^2+7n+6)}6
 =\frac{(n+1)(n+2)(2n+3)}6.
\] Ponieważ $|P_n|=1/n\to0$ i $x^2$ jest ciągła, twierdzenie \ref{ma1:thm:riemann-darboux-equivalence} daje całkę równą $\lim S_n=1/3$.
\end{przyklad}
\begin{przyklad}
\label{ma1:ex:riemann-step}
Dla $a<c<b$ funkcja $\chi_{[c,b]}$ na $[a,b]$ jest całkowalna (ma tylko jedną nieciągłość). Na $[a,c]$ różni się od zera wyłącznie w $c$, więc jej całka na tym przedziale jest zerowa. Na $[c,b]$ jest stała równa $1$, więc całka wynosi $b-c$. Addytywność daje $\int_a^b\chi_{[c,b]}(x)\,\dd x=b-c$.
\end{przyklad}

\subsection{Podstawowe twierdzenie rachunku całkowego}
\begin{twierdzenie}
\label{ma1:thm:riemann-fundamental}
Niech $f$ będzie całkowalna na $[a,b]$ i $x_0\in[a,b]$. Funkcja
\[
 F(x)=\int_{x_0}^x f(t)\,\dd t,\qquad x\in[a,b],
\]
jest ciągła na $[a,b]$. Jeżeli $f$ jest ciągła w $x\in(a,b)$, to $F$ jest różniczkowalna w $x$ i $F'(x)=f(x)$. Przy ciągłości jednostronnej $f$ na końcu analogiczna równość zachodzi dla pochodnej jednostronnej.
\end{twierdzenie}
\begin{proof}
Niech $|f|\le B$ na $[a,b]$. Addytywność i oszacowanie całki dają dla $u<v$
\[
 |F(v)-F(u)|=\left|\int_u^v f(t)\,\dd t\right|\le B(v-u).
\]
A zatem $F$ jest Lipschitzowska (dla $B=0$ jest stała), w szczególności ciągła.

Ustalmy punkt $x$ ciągłości $f$ i przyrost $h\ne0$ taki, że $x+h\in[a,b]$. Addytywność, liniowość i całka funkcji stałej (z uwzględnieniem orientacji dla $h<0$) dają
\[
 \frac{F(x+h)-F(x)}h-f(x)
 =\frac1h\int_x^{x+h}(f(t)-f(x))\,\dd t.
\]
Dla $\varepsilon>0$ z ciągłości $f$ wybieramy $\delta>0$, aby $|f(t)-f(x)|<\varepsilon/2$, gdy $|t-x|<\delta$ i $t\in[a,b]$. Jeżeli $0<|h|<\delta$, każdy $t$ pomiędzy $x$ i $x+h$ spełnia ten warunek. Stosując oszacowanie całki na przedziale z końcami w kolejności rosnącej, otrzymujemy
\[
 \left|\frac{F(x+h)-F(x)}h-f(x)\right|
 \le\frac{\varepsilon}{2|h|}|h|=\frac{\varepsilon}{2}<\varepsilon.
\]
Z definicji pochodnej dostajemy $F'(x)=f(x)$. Przy końcu przedziału używamy jedynie dopuszczalnych przyrostów i ciągłości z odpowiedniej strony.
\end{proof}
\begin{wniosek}
\label{ma1:cor:continuous-primitive-existence}
Każda funkcja ciągła $f$ na niepustym otwartym przedziale $I$ ma funkcję pierwotną. Dla $x_0\in I$ dokładnie jedna z nich spełnia warunek $F(x_0)=0$ i jest dana wzorem
\[
 F(x)=\int_{x_0}^x f(t)\,\dd t.
\]
\end{wniosek}
\begin{proof}
Dla każdego $x\in I$ domknięty odcinek o końcach $x_0,x$ zawiera się w $I$. Całka jest określona przez twierdzenie \ref{ma1:thm:riemann-continuous} i konwencję orientacji. Aby sprawdzić pochodną w ustalonym $x$, wybieramy domknięty przedział zawarty w $I$, zawierający $x_0$ oraz $x$ we wnętrzu. Na tym przedziale stosujemy twierdzenie \ref{ma1:thm:riemann-fundamental}; otrzymujemy $F'(x)=f(x)$. Całka od $x_0$ do $x_0$ wynosi zero. Jednoznaczność wynika z twierdzenia Lagrange'a: różnica dwóch funkcji pierwotnych ma pochodną zerową, jest więc stała na $I$, a warunek normalizacji wymusza, że stała wynosi zero.
\end{proof}


\begin{twierdzenie}[Wzór Newtona--Leibniza]
\label{ma1:thm:newton-leibniz}
Niech $f$ będzie całkowalna na $[a,b]$. Jeżeli $G$ jest ciągła na $[a,b]$, różniczkowalna na $(a,b)$ i $G'=f$ w jego wnętrzu, to
\[
 \int_a^b f(x)\,\dd x=G(b)-G(a).
\]
Nie wymagamy istnienia pochodnych na końcach.
\end{twierdzenie}
\begin{proof}
Na każdym podprzedziale $[x_{i-1},x_i]$ funkcja $G$ jest ciągła i różniczkowalna w jego wnętrzu, więc twierdzenie Lagrange'a daje $\xi_i\in(x_{i-1},x_i)$ takie, że $G(x_i)-G(x_{i-1})=G'(\xi_i)\Delta x_i=f(\xi_i)\Delta x_i$. Sumując, otrzymujemy dokładnie
\[
 \sum_{i=1}^n\bigl(G(x_i)-G(x_{i-1})\bigr)
 =G(b)-G(a)=S(f,P,\xi).
\] Wybieramy podziały na $n$ równych części i odpowiadające im punkty z twierdzenia Lagrange'a. Ich średnice dążą do zera, a całkowalność $f$ i twierdzenie \ref{ma1:thm:riemann-darboux-equivalence} dają granicę sum równą całce. Lewa strona równości dla sum jest niezależna od podziału, więc otrzymujemy tezę.
\end{proof}
Dla funkcji ciągłej można także odjąć od $G$ funkcję $F$ z twierdzenia \ref{ma1:thm:riemann-fundamental}. Różnica ma zerową pochodną we wnętrzu i jest ciągła na końcach, więc jest stała z twierdzenia Lagrange'a. Powyższy dowód obejmuje jednak również całkowalne pochodne nieciągłe.

\begin{przyklad}
\label{ma1:ex:riemann-newton-leibniz}
Z wzorów na funkcje pierwotne i wzoru Newtona--Leibniza otrzymujemy
\[
 \int_0^1 x^2\,\dd x=\left[\frac{x^3}{3}\right]_0^1=\frac13,\qquad
 \int_1^2\frac{\dd x}{x}=[\log x]_1^2=\log2.
\]
Zapis $[G(x)]_a^b$ oznacza $G(b)-G(a)$, a nie iloraz. Obie funkcje podcałkowe są ciągłe na wskazanych przedziałach, więc są całkowalne. Na przedziale przecinającym zero nie wolno stosować drugiego wzoru, gdyż $1/x$ nie jest tam określona i ograniczona. Nie jest to przypadek klasycznej całki Riemanna na przedziale domkniętym.
\end{przyklad}


\begin{wniosek}
Klasa $\mathcal R[a,b]$ funkcji całkowalnych w sensie Riemanna jest kratą wektorową, a granica jednostajna funkcji z tej klasy także do niej należy i jej całka jest granicą całek.
\end{wniosek}
\begin{proof}
Twierdzenie \ref{ma1:thm:riemann-algebra} daje domkniętość na kombinacje liniowe i moduł, a wzory $\max(f,g)=(f+g+|f-g|)/2$ oraz $\min(f,g)=(f+g-|f-g|)/2$ dają domkniętość na maksimum i minimum punktowe.
Jeżeli $f_n\to f$ jednostajnie, to $f$ jest ograniczona, ponieważ dla pewnego $n$ mamy $\|f-f_n\|_\infty<1$. Dla każdego podziału $P$ zachodzi
\[
 U(f,P)-L(f,P)\le U(f_n,P)-L(f_n,P)+2(b-a)\|f-f_n\|_\infty.
\]
Istotnie, na każdym podprzedziale kres górny może wzrosnąć, a kres dolny zmaleć co najwyżej o $\|f-f_n\|_\infty$. Dla zadanej dokładności $\varepsilon>0$ wybieramy najpierw $n$ z $2(b-a)\|f-f_n\|_\infty<\varepsilon/2$, a potem podział z różnicą sum $f_n$ mniejszą od $\varepsilon/2$. Kryterium Darboux daje całkowalność $f$. Z liniowości i twierdzenia \ref{ma1:thm:riemann-order} otrzymujemy
\[
 \left|\int_a^b f_n\,\mathrm dx-\int_a^b f\,\mathrm dx\right|
 \le (b-a)\|f_n-f\|_\infty\longrightarrow0.
\]
\end{proof}

\begin{uwaga}[Słabości całki Riemanna]
Zbieżność punktowa nie zachowuje całkowalności: jeżeli $\Q\cap[0,1]=\{q_1,q_2,\dots\}$, to $\ind{\{q_1,\dots,q_n\}}\to\ind{\Q\cap[0,1]}$. Całka Riemanna nie obejmuje funkcji nieograniczonych ani wszystkich pochodnych (istnieją różniczkowalne funkcje o nieograniczonej pochodnej). Seminorma $\int_a^b|f|$ po utożsamieniu funkcji równych p.w. jest normą, lecz otrzymana przestrzeń nie jest zupełna.
\end{uwaga}
\begin{przyklad}[Dlaczego brakuje zupełności]
Na $[0,1]$ niech $f_n(x)=\min\{n,x^{-1/2}\}$ dla $x>0$ i $f_n(0)=n$. Każde $f_n$ jest ciągłe, a $f_n\uparrow f$ poza zerem, gdzie $f(x)=x^{-1/2}$. Zbieżność zmajoryzowana (lub bezpośrednie całkowanie) daje $\int_0^1|f_n-f|\,\mathrm dx\to0$, zatem $(f_n)$ jest ciągiem Cauchy'ego w normie $L^1$. Nie istnieje jednak ograniczona funkcja całkowalna w sensie Riemanna, która byłaby równa $f$ prawie wszędzie.
\end{przyklad}

\section{Całka Riemanna--Stieltjesa}
\begin{definicja}
Niech $f,\alpha:[a,b]\to\R$. Piszemy $f\in\mathcal R(\alpha)$, jeśli istnieje $I\in\R$ takie, że dla każdego $\varepsilon>0$ istnieje $\delta>0$ i dla każdego podziału znaczonego $\mathcal P$ o $\mesh P<\delta$ zachodzi
\[
\left|\sum_{i=1}^n f(t_i)(\alpha(x_i)-\alpha(x_{i-1}))-I\right|<\varepsilon.
\]
Wtedy $I=\int_a^b f\,\mathrm d\alpha$ jest całką Riemanna--Stieltjesa. Przy $\alpha(x)=x$ otrzymujemy całkę Riemanna.
\end{definicja}
\begin{uwaga}[Rozszerzenie]
Dla niemalejącej $\alpha$ i ograniczonej $f$ można definiować sumy dolne i górne Stieltjesa z przyrostami $\alpha$. Jeżeli $\alpha$ ma skok, nawet ciągła w sensie jednostronnym funkcja $f$ może wymagać zgodności znaczników w otoczeniu tego skoku; zwykła całka Riemanna--Stieltjesa może wówczas nie istnieć.
\end{uwaga}
Dla ograniczonej funkcji $f$ i niepustego zbioru $J\subset[a,b]$ oznaczamy $\osc_J f=\sup_J f-\inf_J f$. Lokalna oscylacja to $\osc_x f=\inf_{\rho>0}\osc_{(x-\rho,x+\rho)\cap[a,b]}f$. Jest równa zeru dokładnie w punktach ciągłości względnej $f$: jeśli oscylacja na otoczeniu jest mała, to w szczególności wszystkie wartości w nim są bliskie $f(x)$; odwrotnie ciągłość daje małą różnicę każdej pary wartości. Dla $r>0$ zbiór $D_r=\{x:\osc_x f\ge r\}$ jest domknięty, ponieważ gdy $\osc_x f<r$, odpowiednio małe otoczenia punktów pobliskich zawierają się w otoczeniu $x$ o oscylacji mniejszej od $r$. W szczególności $D(f)=\bigcup_{n\ge1}D_{1/n}$ jest zbiorem borelowskim.

\begin{twierdzenie}[Narzędzia]\label{it:thm:rs-tools}
Niech $V_a^b(\alpha)=\sup_P\sum_i|\alpha(x_i)-\alpha(x_{i-1})|$.
\begin{enumerate}[label=(\alph*)]
\item Jeśli $f$ jest ciągła i $\alpha$ ma wahanie ograniczone, to $f\in\mathcal R(\alpha)$ i $\bigl|\int_a^b f\,\mathrm d\alpha\bigr|\le\|f\|_\infty V_a^b(\alpha)$.
\item Jeśli $f,\alpha$ mają wahanie ograniczone i nie mają wspólnych punktów nieciągłości, to obie całki Stieltjesa istnieją oraz
$\int_a^b f\,\mathrm d\alpha+\int_a^b\alpha\,\mathrm df=f(b)\alpha(b)-f(a)\alpha(a)$.
\item Jeśli $\alpha\in C^1[a,b]$ i $f\in\mathcal R[a,b]$, to $\int_a^b f\,\mathrm d\alpha=\int_a^b f(x)\alpha'(x)\,\mathrm dx$.
\item Całka jest liniowa w każdej zmiennej, jeśli odpowiednie całki istnieją. Dla niemalejącej $\alpha$ jest dodatnia i monotoniczna względem $f$.
\end{enumerate}
\end{twierdzenie}
\begin{proof}
Dla (a) rozbijamy każdą komórkę jednego podziału według punktów drugiego. Różnica sum przy zagęszczeniu jest ograniczona przez $\omega_f(\mesh P)V_a^b(\alpha)$, gdzie $\omega_f(h)=\sup_{|s-t|\le h}|f(s)-f(t)|$. Porównując dwa podziały z ich wspólnym zagęszczeniem i korzystając z $\omega_f(h)\to0$, dostajemy kryterium Cauchy'ego. Oszacowanie całki wynika z oszacowania każdej sumy.

Dla (b) udowodnimy najpierw istnienie całki. Oznaczmy $V(x)=\Var_{[a,x]}\alpha$ oraz $T=V(b)$. Z addytywności wariacji mamy $V(v)-V(u)=\Var_{[u,v]}\alpha$. Funkcja $V$ jest ciągła w każdym punkcie ciągłości $\alpha$. Istotnie, dla $x<b$ wybieramy podział $[x,b]$, którego suma wariacyjna różni się od $\Var_{[x,b]}\alpha$ o mniej niż $\eta>0$. Dla $y>x$ przed pierwszym punktem tego podziału, zastąpienie jego początku $x$ przez $y$ zmniejsza sumę o co najwyżej $|\alpha(y)-\alpha(x)|$. Z addytywności wariacji wynika
\[
 V(y)-V(x)\le\eta+|\alpha(y)-\alpha(x)|.
\]
Przechodząc do granicy przy $y\downarrow x$, a następnie $\eta\downarrow0$, uzyskujemy ciągłość prawostronną $V$ w punkcie ciągłości $\alpha$. Analogiczny argument dla podziału $[a,x]$ daje ciągłość lewostronną. Na końcach stosujemy tylko odpowiedni kierunek.

Niech $r>0$. Zbiór punktów $D_r$, w których lokalna oscylacja $f$ jest co najmniej $r$, jest skończony. Dla dowolnej skończonej rodziny różnych punktów $D_r$ wybieramy rozłączne względne otoczenia przedziałowe, w każdym z nich dwa punkty o różnicy wartości $f$ większej od $r/2$. Suma tych różnic nie przekracza $\Var_{[a,b]}f$, więc liczba wybranych punktów nie przekracza $2\Var_{[a,b]}f/r$. To daje skończoność $D_r$.

Ponieważ $r>0$, punkty $D_r$ są punktami nieciągłości $f$, a więc punktami ciągłości $\alpha$ i $V$. Wybieramy dla każdego z nich względne otwarte otoczenie przedziałowe $O_x$ o domknięciu tak małym, że suma wariacji $\alpha$ na domknięciach wszystkich $O_x$ jest mniejsza od $\eta$. Dla każdego punktu spoza $D_r$ wybieramy względne otwarte otoczenie, na którym oscylacja $f$ jest mniejsza od $r$. Otrzymujemy otwarte pokrycie zwartego przedziału. Z lematu o liczbie Lebesgue'a istnieje $h>0$ takie, że każda komórka o długości mniejszej od $h$ leży w jednym z tych otoczeń. Komórki przypisane do $O_x$ mają łączną wariację $\alpha$ nie większą niż suma wariacji na domknięciach $O_x$, ponieważ mają rozłączne wnętrza i wariacja jest addytywna. Pozostałe komórki mają oscylację $f$ mniejszą od $r$. Zatem dla każdego takiego podziału $P$,
\[
 \sum_{I\in P}\osc_I f\,\Var_I\alpha
 \le 2\|f\|_\infty\eta+rT.
\]
Wybór najpierw małego $r$, a następnie małego $\eta$, pozwala uczynić prawą stronę dowolnie małą (także gdy $T=0$ lub $f=0$).

Dla dowolnego rozdrobnienia $P'$ podziału $P$ różnica sum Stieltjesa, niezależnie od znaczników, ma moduł nie większy niż lewa strona ostatniej nierówności: rozpisujemy przyrost $\alpha$ na każdej starej komórce jako sumę przyrostów na nowych komórkach i stosujemy nierówność trójkąta. Porównanie dwóch dostatecznie drobnych podziałów ze wspólnym rozdrobnieniem daje kryterium Cauchy'ego dla sum. Istnienie wspólnej granicy wynika z zupełności $\R$: wybieramy ciąg podziałów o średnicach dążących do zera, ich sumy tworzą ciąg Cauchy'ego, a każdą dostatecznie drobną sumę porównujemy z jego granicą. Tak otrzymujemy $\int f\,\mathrm d\alpha$. Zamiana ról $f,\alpha$ daje istnienie $\int\alpha\,\mathrm df$.

Dla dowolnego podziału zachodzi dokładna tożsamość
\[
 \sum_i f(x_{i-1})(\alpha(x_i)-\alpha(x_{i-1}))
 +\sum_i\alpha(x_i)(f(x_i)-f(x_{i-1}))
 =f(b)\alpha(b)-f(a)\alpha(a).
\]
Wynika ona z rozpisania każdego składnika jako różnicy kolejnych iloczynów i teleskopowania. Przechodząc do granicy przy średnicy dążącej do zera, otrzymujemy wzór całkowania przez części. W (c) twierdzenie Lagrange'a daje $\alpha(x_i)-\alpha(x_{i-1})=\alpha'(\xi_i)(x_i-x_{i-1})$; z jednostajnej ciągłości $\alpha'$ wynika, że możemy zastąpić $\xi_i$ przez $t_i$. Ostatnia teza wynika bezpośrednio z definicji.
\end{proof}
\begin{twierdzenie}[Kryterium oscylacyjne]\label{it:thm:rs-oscillation}
Niech $\alpha:[a,b]\to\R$ będzie niemalejąca, prawostronnie ciągła, a $f:[a,b]\to\R$ ograniczona. Niech $\mu_\alpha$ będzie uzupełnioną miarą Lebesgue'a--Stieltjesa określoną przez $\mu_\alpha((s,t])=\alpha(t)-\alpha(s)$ dla $a\le s<t\le b$, przy $\mu_\alpha(\{a\})=0$. Następujące warunki są równoważne:
\begin{enumerate}[label=(\roman*)]
\item $f\in\mathcal R(\alpha)$;
\item $\mu_\alpha(D(f))=0$;
\item dla każdego $\varepsilon>0$ istnieje $\delta>0$ takie, że dla każdego podziału $P$ o $\mesh P<\delta$ zachodzi $\sum_i\osc_{[x_{i-1},x_i]}f\,(\alpha(x_i)-\alpha(x_{i-1}))<\varepsilon$.
\end{enumerate}
Wówczas $\int_a^b f\,\mathrm d\alpha=\int_{(a,b]}f\,\mathrm d\mu_\alpha$.
\end{twierdzenie}
\begin{proof}
Dla niemalejącej $\alpha$ suma górna minus dolna jest dokładnie sumą z (iii). Porównywanie sum dla dowolnych znaczników dowodzi (i)$\Leftrightarrow$(iii), tak samo jak w kryterium Darboux.

Niech $M=\|f\|_\infty$. Jeżeli $\mu_\alpha(D(f))=0$, to dla $r>0$ zwarty zbiór $D_r=\{x:\osc_x f\ge r\}$ można pokryć zbiorem otwartym $G$ o dowolnie małej mierze $\mu_\alpha$. Poza $G$ oscylacje $f$ są lokalnie mniejsze od $r$. Pokrycie przez $G$ i względne otwarte otoczenia punktów spoza $D_r$, na których oscylacja jest mniejsza od $r$, ma dodatnią liczbę Lebesgue'a. Przy średnicy mniejszej od tej liczby każda komórka mieści się w jednym z tych zbiorów. Przedziały przypisane do $G$ wnoszą do sumy (iii) nie więcej niż $2M\mu_\alpha(G)$, a pozostałe nie więcej niż $r(\alpha(b)-\alpha(a))$. To dowodzi (iii).

Jeśli zachodzi (iii), wybieramy podziały z coraz mniejszą średnicą i punktami podziału poza atomami $\mu_\alpha$ (można to zrobić, ponieważ atomów jest co najwyżej przeliczalnie wiele). Komórki o wnętrzach przecinających $D_r$ mają oscylację co najmniej $r$, więc ich suma miar jest arbitralnie mała. Punkty wewnętrzne leżące na granicach komórek nie są atomami, zaś obecność punktu $b$ w $D_r$ powoduje oscylację w ostatniej komórce. Przy konwencji $\mu_\alpha(\{a\})=0$ otrzymujemy $\mu_\alpha(D_r)=0$. Sumowanie po $r=1/n$ daje (ii).

Dla dowolnego ciągu podziałów o średnicach dążących do zera definiujemy funkcje schodkowe dolne i górne na komórkach półotwartych $(x_{i-1},x_i]$, używając infimum i supremum po domkniętych komórkach. W każdym punkcie ciągłości $f$ obie funkcje zbiegają do $f$. Ponieważ $D(f)$ jest zbiorem $\mu_\alpha$-zerowym, uzupełnienie miary zapewnia mierzalność $f$ (poza tym zbiorem jest granicą funkcji schodkowych). Funkcje schodkowe zbiegają zatem prawie wszędzie do $f$ i mają całki Lebesgue'a--Stieltjesa równe sumom dolnym i górnym. Zbieżność zmajoryzowana daje ostatnią równość.
\end{proof}
\begin{uwaga}[Atomy i końce przedziału]
Jeżeli $\mu_\alpha(\{c\})=\alpha(c)-\alpha(c-)>0$ dla $c\in(a,b]$, warunek z poprzedniego twierdzenia wymusza ciągłość $f$ w $c$ (przy $c=b$ w sensie lewostronnym). Konwencja $\mu_\alpha(\{a\})=0$ wynika z użycia przyrostów od $\alpha(a)$; przy innej konwencji konieczna jest osobna poprawka za atom w $a$.
\end{uwaga}
\begin{przyklad}[Część ciągła i atomy]
Na $[0,1]$ weźmy $\alpha(x)=x+\tfrac12\ind{[1/3,1]}(x)+\ind{[2/3,1]}(x)$. Dla $f(x)=x^2$ istnieje całka Stieltjesa i
\[
\int_0^1x^2\,\mathrm d\alpha=\int_0^1x^2\,\mathrm dx+\frac12\left(\frac13\right)^2+\left(\frac23\right)^2=\frac13+\frac1{18}+\frac49=\frac56.
\]
Dla $\alpha(x)=\lfloor x\rfloor$ na $[0,n]$, $n\in\N$, oraz ciągłej $f$ otrzymujemy analogicznie $\int_0^n f\,\mathrm d\alpha=\sum_{k=1}^n f(k)$.
\end{przyklad}
\begin{przyklad}[Wspólny skok]
Niech $c\in(a,b)$ i $f=\alpha=\ind{[c,b]}$. Zwykła całka Riemanna--Stieltjesa nie istnieje: na przedziale zawierającym $c$ suma zależy od tego, czy znacznik wybrano przed czy po $c$. Natomiast pojęcie całki Kurzweila--Stieltjesa pozwoli otrzymać wartość $1$.
\end{przyklad}
\begin{przyklad}[Dlaczego nie wystarczy wybrać jednego podziału]
Niech $a<c<b$, $\alpha=\ind{[c,b]}$ i $f=\ind{(c,b]}$, przy czym $f(c)=0$. Dla podziału zawierającego $c$ każda suma dolna i górna Stieltjesa może być równa zeru: dodatni przyrost funkcji całkującej kończy się w $c$, a na $[x_{i-1},c]$ funkcja $f$ jest równa zeru. Podziały przecinające $c$ mogą jednak dawać sumy $0$ albo $1$, zależnie od znacznika. Warunek „dla wszystkich dostatecznie drobnych podziałów” w definicji Riemanna--Stieltjesa jest więc istotny; zagęszczanie jednego wybranego podziału nie wystarcza.
\end{przyklad}
\begin{przyklad}[Wpływ funkcji całkującej na całkowalność]
Na $[0,1]$ niech $f=\ind{\Q\cap(0,1/4]}$ i $\alpha=\ind{[1/2,1]}$. Funkcja $f$ nie jest całkowalna Riemanna względem długości, gdyż jest nieciągła na całym $(0,1/4)$. Jest jednak ciągła i równa zeru w otoczeniu jedynego atomu $\mu_\alpha$, zatem $\int_0^1 f\,\mathrm d\alpha=0$ w sensie Riemanna--Stieltjesa. Z kolei funkcja z przykładu 3.7 jest Riemanna całkowalna względem długości, ale nie względem wskazanego $\alpha$. Dla różnych funkcji całkujących klasy całkowalności nie muszą być porównywalne przez inkluzję.
\end{przyklad}

% ============================================================
% UZUPELNIENIE MATEMATYCZNE: ROZDZIAL 3
% ============================================================
\subsection*{Rachunek praktyczny: skoki i kontrola błędu}
\begin{twierdzenie}[Jeden skok w zwykłej całce Stieltjesa]\label{tw:rs-jeden-skok}
Niech $a<c<b$, $A\ne0$, $\alpha(x)=A\ind{[c,b]}(x)$ dla $x\in[a,b]$, a $f:[a,b]\to\R$ będzie ograniczona. Całka Riemanna--Stieltjesa $\int_a^b f\,\mathrm d\alpha$ istnieje wtedy i tylko wtedy, gdy $f$ jest ciągła w $c$. Wówczas $\int_a^b f\,\mathrm d\alpha=Af(c)$. Jeśli skok jest umieszczony w $b$, warunkiem jest ciągłość lewostronna $f$ w $b$, a dla skoku umieszczonego w $a$ funkcją $\ind{[a,b]}$ integrator jest stały i całka wynosi zero.
\end{twierdzenie}
\begin{proof}
W każdej sumie Stieltjesa wszystkie przyrosty $\alpha$ są zerowe poza jedną komórką $[u,v]$, dla której $u<c\le v$. Suma wynosi więc $Af(t)$, gdzie $t\in[u,v]$. Jeśli $f$ jest ciągła w $c$, dostatecznie mała średnica podziału wymusza $|f(t)-f(c)|<\varepsilon/|A|$.

Odwrotnie, załóżmy istnienie całki i rozważmy dowolny ciąg $t_n\to c$. Dobieramy komórki $[u_n,v_n]$ o długościach dążących do zera, zawierające zarówno $c$, jak i $t_n$, przy czym $u_n<c\le v_n$. Uzupełniamy je do podziałów o średnicach dążących do zera. Sumy ze znacznikiem $t_n$ oraz ze znacznikiem $c$ muszą mieć tę samą granicę. Zatem $A(f(t_n)-f(c))\to0$. Ponieważ $(t_n)$ był dowolny, $f$ jest ciągła w $c$. Wariant dla $c=b$ ma identyczny dowód z granicą jednostronną. W punkcie $a$ integrator $A\ind{[a,b]}$ ma wszystkie przyrosty równe zeru.
\end{proof}
\begin{wniosek}[Skończona liczba atomów]\label{wn:rs-suma-skokow}
Niech $\alpha(x)=h(x)+\sum_{k=1}^m A_k\ind{[c_k,b]}(x)$ na $[a,b]$, gdzie $h\in\BV[a,b]$ jest ciągła, $a<c_1<\cdots<c_m\le b$, a $A_k\in\R$. Jeśli $f\in C[a,b]$, to
\[
\int_a^b f\,\mathrm d\alpha=\int_a^b f\,\mathrm dh+\sum_{k=1}^m A_kf(c_k).
\]
W przypadku samej części skokowej wystarczy ciągłość $f$ we wszystkich $c_k$ (jednostronna w $b$).
\end{wniosek}
\begin{proof}
Stosujemy liniowość całki i twierdzenie \ref{tw:rs-jeden-skok}. Dla ciągłego $f$ całka względem $h\in\BV$ istnieje na mocy twierdzenia z początku rozdziału.
\end{proof}
\begin{twierdzenie}[Oszacowanie błędu sumy Stieltjesa]\label{tw:rs-blad-kwadratury}
Niech $f\in C[a,b]$, $\alpha\in\BV[a,b]$, a $\mathcal P$ będzie dowolnym podziałem znaczonym o średnicy $h$. Wówczas
\[
\left|S(f,\alpha;\mathcal P)-\int_a^b f\,\mathrm d\alpha\right|
\le \omega_f(h)\Var_a^b(\alpha),
\qquad
\omega_f(h):=\sup_{\substack{x,y\in[a,b]\\|x-y|\le h}}|f(x)-f(y)|.
\]
Jeżeli $f$ jest lipschitzowska ze stałą $L$, błąd nie przekracza $Lh\Var_a^b(\alpha)$.
\end{twierdzenie}
\begin{proof}
Ustalmy podział $\mathcal P$ i zagęszczajmy każdą jego komórkę niezależnie, wybierając dowolne nowe znaczniki. Wewnątrz komórki $[x_{i-1},x_i]$ oba rodzaje znaczników są oddalone najwyżej o $h$; zatem różnica starego wkładu i sumy po jej zagęszczeniu ma moduł nie większy niż $\omega_f(h)\Var_{x_{i-1}}^{x_i}(\alpha)$. Po zsumowaniu otrzymujemy podane oszacowanie dla różnicy sum. Przy średnicy zagęszczenia dążącej do zera nowa suma dąży do całki Riemanna--Stieltjesa.
\end{proof}
\begin{przyklad}[Skończona kwadratura ze znanym błędem]
Niech $f(x)=e^x$ na $[0,1]$ i $\alpha(x)=x+2\ind{[1/2,1]}(x)$. Wówczas
$\int_0^1e^x\,\mathrm d\alpha=e-1+2\sqrt e$.
Dla dowolnego podziału znaczonego o średnicy $h$ twierdzenie \ref{tw:rs-blad-kwadratury} daje
$|S-(e-1+2\sqrt e)|\le 3eh$, gdyż $\Var_0^1(\alpha)=3$ i $\omega_f(h)\le eh$.
\end{przyklad}
\begin{uwaga}[Zgodność trzech konwencji]
W zwykłej całce Riemanna--Stieltjesa przy skoku w $c$ dowolne znaczniki wymuszają ciągłość $f$ w $c$. W całce Kurzweila--Stieltjesa cechowanie może wymusić znacznik $c$ i wówczas poprawny wkład skoku wynosi $f(c)\Delta^-\alpha(c)$, nawet jeśli $f$ nie jest ciągła. Całka Lebesgue'a--Stieltjesa wymaga mierzalności i całkowalności względem miary, lecz nie ciągłości w atomie: wkład wynosi również $f(c)\mu_\alpha(\{c\})$. Zbieżność wartości całek nie oznacza więc równości ich dziedzin.
\end{uwaga}


\section{Całka Kurzweila--Henstocka}
\providecommand{\atlasref}[2]{#2}
\label{it:sec:kurzweil-henstock}\label{ma1:ch:kurzweil-henstock}
% Import z AM1-36, zachowane etykiety źródłowe i historyczne IT.
Symbol $\lambda_1$ oznacza miarę Lebesgue'a na prostej. Dla ograniczonej funkcji $F$ i niepustego podprzedziału $I$ piszemy $\osc_I F=\sup_I F-\inf_I F$. Gdy $F$ jest ciągła, a $I$ zwarty, supremum i infimum są osiągane, więc $\osc_I F=\max_I F-\min_I F$. Symbole $\AC$ i $\BV$ oznaczają klasy funkcji bezwzględnie ciągłych i o ograniczonej wariacji; $\Var_I F$ jest supremum sum modułów przyrostów po podziałach $I$. Funkcja $F$ jest bezwzględnie ciągła na $[a,b]$, jeżeli dla każdego $\varepsilon>0$ istnieje $\eta>0$ takie, że dla każdej skończonej rodziny przedziałów $[u_j,v_j]$ o rozłącznych wnętrzach z $\sum_j(v_j-u_j)<\eta$ zachodzi $\sum_j|F(v_j)-F(u_j)|<\varepsilon$. Przez $\mathcal R[a,b]$ oznaczamy funkcje całkowalne w sensie Riemanna, a przez $\mathcal L^1$ konkretne funkcje całkowalne w sensie Lebesgue'a; $L^1$ oznacza ich klasy równości prawie wszędzie.
\subsection{Cechowania i definicja całki}
\begin{definicja}
\label{ma1:kh:def:kh-gauge}\label{it:def:kh-gauge}
Dla $a<b$ podział znaczony to skończona rodzina komórek $[x_{i-1},x_i]$, gdzie $a=x_0<x_1<\cdots<x_m=b$, ze znacznikami $t_i\in[x_{i-1},x_i]$. Jego suma ma postać
\[
 S(f,\mathcal P)=\sum_{i=1}^m f(t_i)(x_i-x_{i-1}).
\]
Podział częściowy składa się z przedziałów o dodatnich długościach, parami rozłącznych wnętrzach i znacznikach należących do odpowiednich przedziałów. Na przedziale jednopunktowym przyjmujemy całkę zero.


\emph{Cechowaniem} na $[a,b]$ nazywamy dowolną funkcję $\delta:[a,b]\to(0,\infty)$. Podział znaczony $\mathcal P=\{(t_i,[x_{i-1},x_i])\}$ jest $\delta$-drobny, jeśli $t_i\in[x_{i-1},x_i]\subset(t_i-\delta(t_i),t_i+\delta(t_i))$ dla każdego $i$.
\end{definicja}
\begin{lemat}[Cousina]\label{ma1:kh:lem:kh-cousin}\label{it:lem:kh-cousin}
Niech $a<b$. Dla każdego cechowania $\delta$ na $[a,b]$ istnieje $\delta$-drobny podział znaczony tego przedziału.
\end{lemat}
\begin{proof}
Przypuśćmy, że przedział nie ma takiego podziału. Dzielimy go na połowy; co najmniej jedna również nie ma drobnego podziału, bo dwa drobne podziały połówek można połączyć. Powtarzając, dostajemy zstępujący ciąg domkniętych przedziałów $J_n$ długości dążących do zera, z których żaden nie ma podziału drobnego. Ich przecięcie jest punktem $t$. Dla dostatecznie dużego $n$ mamy $J_n\subset(t-\delta(t),t+\delta(t))$, więc pojedynczy znaczony przedział $(t,J_n)$ jest drobny. Sprzeczność.
\end{proof}
\begin{definicja}
\label{ma1:kh:def:kh-integral}\label{it:def:kh-integral}
Funkcja $f:[a,b]\to\R$ jest całkowalna w sensie Kurzweila--Henstocka z całką $I$, jeżeli dla każdego $\varepsilon>0$ istnieje cechowanie $\delta$ takie, że dla każdego $\delta$-drobnego podziału znaczonego $\mathcal P$ zachodzi $|S(f,\mathcal P)-I|<\varepsilon$.
\end{definicja}
\begin{stwierdzenie}[Jednoznaczność]
\label{ma1:kh:prop:kh-unique}\label{it:prop:kh-unique}
Wartość całki HK jest wyznaczona jednoznacznie.
\end{stwierdzenie}
\begin{proof}
Jeśli $I,J$ spełniają definicję, dla $\varepsilon>0$ bierzemy minimum cechowań odpowiadających błędom $\varepsilon/2$ dla obu wartości. Lemat Cousina zapewnia istnienie podziału drobnego dla tego minimum. Dla jego sumy $S$ mamy $|I-J|\le |I-S|+|S-J|<\varepsilon$. Dowolność $\varepsilon$ daje $I=J$.
\end{proof}
\begin{twierdzenie}[Kryterium Cauchy'ego]\label{ma1:kh:thm:kh-cauchy}\label{it:thm:kh-cauchy}
Funkcja $f$ jest całkowalna w sensie HK wtedy i tylko wtedy, gdy dla każdego $\varepsilon>0$ istnieje cechowanie $\delta$ takie, że $|S(f,\mathcal P)-S(f,\mathcal Q)|<\varepsilon$ dla wszystkich $\delta$-drobnych podziałów $\mathcal P,\mathcal Q$.
\end{twierdzenie}
\begin{proof}
Konieczność wynika z nierówności trójkąta. Dla dostateczności wybieramy cechowania $\delta_n$ odpowiadające błędom $2^{-n}$, zastępujemy je kolejno minimami z poprzednimi i wybieramy $\delta_n$-drobne podziały $\mathcal P_n$ z lematu Cousina. Dla $m\ge n$ mamy $\delta_m\le\delta_n$, więc oba podziały $\mathcal P_n,\mathcal P_m$ są $\delta_n$-drobne i $|S(f,\mathcal P_n)-S(f,\mathcal P_m)|<2^{-n}$. Sumy tworzą zatem ciąg Cauchy'ego. Niech $I$ będzie jego granicą. Dla danego $n$ i dowolnego $\delta_n$-drobnego podziału $\mathcal P$ porównujemy jego sumę z sumami $S(f,\mathcal P_m)$, $m\ge n$, i przechodzimy z $m\to\infty$. Otrzymujemy $|S(f,\mathcal P)-I|\le2^{-n}$. Dla dowolnego $\varepsilon>0$ wybieramy $n$ z $2^{-n}<\varepsilon$; cechowanie $\delta_n$ spełnia definicję całki.
\end{proof}
\begin{twierdzenie}[Liniowość, porządek i lokalność]\label{ma1:kh:thm:kh-locality}\label{it:thm:kh-locality}
Klasa funkcji HK-całkowalnych jest przestrzenią wektorową, całka jest liniowa i monotoniczna. Jeśli $a<c<b$, to $f$ jest całkowalna w sensie HK na $[a,b]$ wtedy i tylko wtedy, gdy jest całkowalna na obu $[a,c]$ i $[c,b]$; wówczas $\int_a^b f=\int_a^c f+\int_c^b f$.
\end{twierdzenie}
\begin{proof}
Dla funkcji całkowalnych $f,g$ i liczb $\alpha,\beta\in\R$ wybieramy minimum cechowań dla $f,g$ z błędem $\varepsilon/(1+|\alpha|+|\beta|)$. Liniowość sum daje
\[
 |S(\alpha f+\beta g,\mathcal P)-(\alpha I_f+\beta I_g)|
 <\varepsilon.
\]
Jeżeli $f\le g$, to dla każdego podziału $S(f,\mathcal P)\le S(g,\mathcal P)$. Wybierając wspólne cechowanie dla błędu $\varepsilon/2$ i podział z lematu Cousina, uzyskujemy $I_f\le I_g+\varepsilon$. Dowolność $\varepsilon$ daje monotoniczność.

Jeżeli $f$ jest całkowalna na $[a,b]$, ustalamy cechowanie $\delta$ z kryterium Cauchy'ego dla błędu $\varepsilon$. Na $[a,c]$ wybieramy jeden $\delta$-drobny podział $\mathcal R$ z lematu Cousina. Dla dwóch dowolnych drobnych podziałów $\mathcal P,\mathcal Q$ przedziału $[c,b]$ ich połączenia z $\mathcal R$ są drobnymi podziałami $[a,b]$. Po odjęciu wspólnego składnika $S(f,\mathcal R)$ otrzymujemy $|S(f,\mathcal P)-S(f,\mathcal Q)|<\varepsilon$. Kryterium Cauchy'ego daje całkowalność na $[c,b]$; zamiana ról obu przedziałów daje całkowalność na $[a,c]$. Stosując ten argument dwukrotnie, otrzymujemy całkowalność na każdym podprzedziale. W drugą stronę wybieramy cechowania na obu podprzedziałach dla błędu $\varepsilon/2$, dla $t\ne c$ dodatkowo mniejsze od $|t-c|$, a w punkcie $c$ bierzemy minimum cechowań. Każdy przedział drobnego podziału przecinający $c$ ma znacznik $c$; rozcinamy go w $c$. Powstałe części dają drobne podziały z tą samą sumą, ponieważ przy rozcięciu zachowujemy znacznik $c$, a długości dodają się. Każda z dwóch sum różni się od odpowiedniej całki lokalnej o mniej niż $\varepsilon/2$. Nierówność trójkąta daje błąd mniejszy od $\varepsilon$ względem sumy całek lokalnych, co dowodzi całkowalności i addytywności. Jednoznaczność zapewnia tę samą równość także wtedy, gdy zaczynaliśmy od całkowalności na całym przedziale.
\end{proof}
\begin{stwierdzenie}
$\mathcal R[a,b]\subset\HK[a,b]$ i obie całki są równe.
\end{stwierdzenie}
\begin{proof}
Z twierdzenia \ref{ma1:thm:riemann-darboux-equivalence} dla każdego $\varepsilon>0$ istnieje $\delta_0>0$ takie, że wszystkie sumy Riemanna dla podziałów o średnicy mniejszej od $\delta_0$ różnią się od całki Riemanna o mniej niż $\varepsilon$. Wybieramy cechowanie stałe $\delta(t)=\delta_0/2$. Każda komórka drobnego podziału ma długość mniejszą od $2\delta(t)=\delta_0$, więc jego suma spełnia żądane oszacowanie.
\end{proof}
\begin{przyklad}[Funkcja Dirichleta jest HK-całkowalna]\label{ma1:kh:ex:dirichlet}
Niech $\Q\cap[0,1]=\{q_1,q_2,\dots\}$ i $D=\chi_{\Q\cap[0,1]}$. Dla $\varepsilon>0$ kładziemy $\delta(q_k)=\varepsilon 2^{-k-3}$, a w punktach niewymiernych np. $\delta(t)=1$. Grupujemy wszystkie komórki o tym samym znaczniku $q_k$. Ich wnętrza są rozłączne, a wszystkie leżą w $(q_k-\delta(q_k),q_k+\delta(q_k))$, więc ich \emph{łączna} długość nie przekracza $2\delta(q_k)$, również gdy znacznik jest wspólnym końcem dwóch komórek. Stąd dla każdego drobnego podziału $0\le S(D,\mathcal P)\le\sum_k 2\delta(q_k)<\varepsilon$. Całka HK jest równa $0$.
\end{przyklad}
\begin{lemat}[Zmiany na zbiorze miary zero]\label{ma1:kh:lem:kh-null-modification}\label{it:lem:kh-null-modification}
Jeśli $h:[a,b]\to\R$ ma skończone wartości i jest równa zeru poza zbiorem $N$ miary Lebesgue'a zero, to jej całka HK istnieje i wynosi $0$. To samo dotyczy całki McShane'a. Zmiana funkcji całkowalnej na zbiorze miary zero nie zmienia wartości żadnej z tych całek.
\end{lemat}
\begin{proof}
Rozbijamy $N$ na zbiory $N_k=\{t\in N:k-1<|h(t)|\le k\}$, $k\ge1$. Dla ustalonego $\varepsilon>0$ wybieramy otwarte pokrycia $U_k\supset N_k$ z $\lambda_1(U_k)<\varepsilon2^{-k-1}/k$. Jeśli $t\in N_k$, dobieramy $\delta(t)>0$ tak, aby $(t-\delta(t),t+\delta(t))\subset U_k$; w pozostałych punktach, również w punktach $N$ o $h(t)=0$, cechowanie jest dowolne, gdyż ich wkład do sumy wynosi zero. W podziale drobnym, również wolnym w sensie McShane'a, komórki z oznaczonymi punktami należącymi do $N_k$ mają rozłączne wnętrza i leżą w $U_k$. Ich wkład do modułu sumy jest co najwyżej $k\lambda_1(U_k)$. Sumując po $k$, dostajemy błąd mniejszy od $\varepsilon$. Jeśli dwie funkcje o skończonych wartościach różnią się tylko na zbiorze zerowym, ich różnica spełnia powyższy warunek. Liniowość całki HK (twierdzenie \ref{it:thm:kh-locality}) i całki McShane'a (twierdzenie \ref{it:thm:mcshane-basic}) daje ostatnią tezę, w obu kierunkach.
\end{proof}

\begin{center}
\begin{tikzpicture}[x=1.25cm,y=1cm,>=Stealth]
 \draw[->] (-.25,0)--(7.4,0);
 \draw[very thick,blue!65!black] (0,0)--(2,0);
 \draw[very thick,teal!70!black] (2,0)--(4.8,0);
 \draw[very thick,blue!65!black] (4.8,0)--(7,0);
 \foreach \x/\name in {0/{a=x_0},2/{x_1},4.8/{x_2},7/{x_3=b}}
   {\draw (\x,.10)--(\x,-.1) node[below] {$\name$};}
 \foreach \x/\name in {1/t_1,3.2/t_2,6/t_3}
   {\fill (\x,0) circle (2pt) node[above] {$\name$};}
 \draw[blue!65!black,<->] (-.2,.65)--(2.2,.65);
 
 \draw[teal!70!black,<->] (1.5,1.2)--(4.9,1.2);
 \draw[blue!65!black,<->] (4.6,.65)--(7.4,.65);
 \node[above] at (3.2,1.25) {$t_2\pm\delta(t_2)$};
\end{tikzpicture}

\small Rysunek: komórki podziału i otoczenia cechowania; ich rozmiary zależą od znacznika. Każda komórka musi zawierać się w swoim otwartym otoczeniu.
\end{center}
\subsection{Pochodne i lemat Saksa--Henstocka}
\begin{przyklad}[Wszystkie pochodne są całkowalne HK]
Jeśli $F:[a,b]\to\R$ jest ciągła i różniczkowalna na $(a,b)$, to funkcja $f$ równa $F'$ we wnętrzu i mająca dowolne skończone wartości w $a,b$ jest HK-całkowalna oraz $\int_a^b f(x)\,\mathrm dx=F(b)-F(a)$. Dla $\varepsilon>0$ przy każdym $t\in(a,b)$ wybieramy $\delta(t)$ tak, aby $|F(y)-F(t)-f(t)(y-t)|\le\varepsilon|y-t|/(4(b-a))$ dla bliskich $y$. Na komórce $[u,v]$ ze znacznikiem $t$ odejmujemy przyrosty $F(v)-F(t)$ oraz $F(u)-F(t)$; nierówność trójkąta daje
\[
 |F(v)-F(u)-f(t)(v-u)|
 \le \frac{\varepsilon}{4(b-a)}(|v-t|+|u-t|)
 =\frac{\varepsilon}{4(b-a)}(v-u).
\]
Dla $t=a$ i $t=b$ korzystamy tylko z ciągłości $F$: dobieramy $\delta(t)>0$ tak, aby $|F(y)-F(t)|+|f(t)||y-t|<\varepsilon/4$ przy $|y-t|<\delta(t)$, $y\in[a,b]$. Komórka ze znacznikiem $a$ lub $b$ musi mieć ten punkt jako koniec; każdego z tych znaczników można użyć co najwyżej raz. Łączny błąd komórek ze znacznikami wewnętrznymi jest nie większy niż $\varepsilon/4$, a komórek ze znacznikami brzegowymi mniejszy niż $\varepsilon/2$. Suma przyrostów $F$ teleskopuje się do $F(b)-F(a)$, więc całkowity błąd jest mniejszy od $\varepsilon$. Nie potrzeba istnienia pochodnych jednostronnych na końcach.
\end{przyklad}
\begin{przyklad}[Pochodna nieabsolutnie całkowalna]\label{ma1:kh:ex:conditional-derivative}
Niech $F(x)=x^2\cos(\pi/x^2)$ dla $x\in(0,1]$ i $F(0)=0$. Wtedy $F'(0)=0$ oraz $F'(x)=2x\cos(\pi/x^2)+(2\pi/x)\sin(\pi/x^2)$. Pochodna jest całkowalna HK i $\int_0^1F'(x)\,\mathrm dx=F(1)-F(0)=-1$. Natomiast rozszerzona całka Lebesgue'a $\int_0^1|F'(x)|\,\mathrm d\lambda_1(x)=\infty$:  składnik $2x\cos(\pi/x^2)$ ma całkowalny moduł, ograniczony przez $2x$. Dla drugiego składnika podstawienie $u=\pi/x^2$ daje
\[
 \int_c^1\frac{2\pi}{x}|\sin(\pi/x^2)|\,\mathrm dx
 =\pi\int_\pi^{\pi/c^2}\frac{|\sin u|}{u}\,\mathrm du.
\]
Na każdym przedziale $[k\pi,(k+1)\pi]$, $k\ge1$, całka po prawej bez czynnika $\pi$ jest nie mniejsza niż $2/((k+1)\pi)$. Szereg harmoniczny jest rozbieżny. Z nierówności $|A+B|\ge|B|-|A|$ wynika rozbieżność całki modułu całej pochodnej.
\end{przyklad}
\begin{lemat}[Saksa--Henstocka: uzupełnianie podziału częściowego]\label{ma1:kh:lem:kh-saks-henstock}\label{it:lem:kh-saks-henstock}\label{ma1:kh:lem:hk-uzupelnianie}\label{lem:hk-uzupelnianie}
Niech $f\in\HK[a,b]$, $F(x)=\int_a^x f$, a $\delta$ będzie cechowaniem, dla którego wszystkie $\delta$-drobne sumy na $[a,b]$ różnią się od $F(b)-F(a)$ o mniej niż $\varepsilon$. Jeśli $\mathcal Q=\{(t_j,[u_j,v_j])\}_{j=1}^m$ jest dowolnym $\delta$-drobnym podziałem częściowym z rozłącznymi wnętrzami, to
\[
\left|\sum_{j=1}^m\bigl(f(t_j)(v_j-u_j)-F(v_j)+F(u_j)\bigr)\right|\le\varepsilon.
\]
Ponadto suma modułów wyrazów w nawiasach jest nie większa niż $2\varepsilon$.
\end{lemat}
\begin{proof}
Uporządkujmy przedziały częściowego podziału. Domknięcia luk pomiędzy tymi przedziałami oraz skrajnych luk przy $a,b$ tworzą skończoną rodzinę co najwyżej $m+1$ przedziałów $J_k$ o dodatniej długości i rozłącznych wnętrzach. Wraz z komórkami częściowego podziału pokrywają $[a,b]$; luki jednopunktowe pomijamy. Dla dowolnego $\eta>0$ lokalność całki HK i lemat Cousina pozwalają wybrać na każdym $J_k$ podział znaczony, drobny względem $\delta$ oraz względem cechowania dającego błąd $\eta/(m+1)$ dla lokalnej całki po $J_k$. Połączenie tych podziałów z $\mathcal Q$ daje pełny podział $\delta$-drobny. Przyrosty $F$ na wszystkich komórkach sumują się do $F(b)-F(a)$. Błąd na $J_k$ jest mniejszy od $\eta/(m+1)$, zatem suma błędów na $\mathcal Q$ ma moduł nie większy niż $\varepsilon+\eta$. Przechodzimy z $\eta\downarrow0$.

Ten sam argument działa dla każdego podzbioru komórek $\mathcal Q$. Stosując go osobno do komórek o błędzie dodatnim i ujemnym, widzimy, że suma błędów dodatnich oraz modułów błędów ujemnych nie przekracza osobno $\varepsilon$. Ich suma jest więc nie większa niż $2\varepsilon$.
\end{proof}
\begin{lemat}[Pięciokrotne powiększenie przedziałów]\label{ma1:kh:lem:fivefold-cover}
Niech $\mathcal I$ będzie rodziną domkniętych przedziałów o dodatnich długościach ograniczonych z góry przez $L<\infty$. Istnieje co najwyżej przeliczalna podrodzina $\mathcal J$ przedziałów o parami rozłącznych wnętrzach taka, że
\[
 \bigcup_{I\in\mathcal I} I\subseteq\bigcup_{J\in\mathcal J}5J,
\]
gdzie $5J$ jest domkniętym przedziałem o tym samym środku co $J$ i pięciokrotnie większej długości.
\end{lemat}
\begin{proof}
Dzielimy rodzinę na klasy $2^{-k-1}L<|I|\le 2^{-k}L$, $k=0,1,\ldots$. W każdej kolejnej klasie wybieramy maksymalną względem inkluzji podrodzinę o rozłącznych wnętrzach, rozłącznych także z wnętrzami wcześniej wybranych przedziałów. Taka rodzina istnieje na mocy lematu Zorna, ponieważ suma łańcucha rodzin dopuszczalnych jest dopuszczalna. Suma wszystkich wybranych rodzin jest przeliczalna lub skończona: każdemu wnętrzu przypisujemy pierwszą należącą do niego liczbę wymierną w ustalonej enumeracji $\Q$, a rozłączność daje różnowartościowość tego przyporządkowania.
Jeśli $I$ z klasy $k$ nie został wybrany, jego wnętrze przecina wnętrze pewnego wybranego $J$ z klasy o indeksie co najwyżej $k$ (w przeciwnym razie można byłoby go dołączyć na etapie $k$). Wtedy $|J|>|I|/2$. Dla środków $c_I,c_J$ mamy $|c_I-c_J|<(|I|+|J|)/2$. Dla $x\in I$ stąd
\[
 |x-c_J|\le |x-c_I|+|c_I-c_J|<|I|+|J|/2<5|J|/2.
\]
Zatem $I\subset 5J$. Przedział wybrany zawiera się oczywiście we własnym pięciokrotnym powiększeniu. To dowodzi tezy także dla rodziny pustej.
\end{proof}

\begin{twierdzenie}[Prawie wszędzie pochodna pierwotnej]\label{ma1:kh:thm:kh-primitive-derivative}\label{it:thm:kh-primitive-derivative}
Jeśli $f\in\HK[a,b]$, a $F(x)=\int_a^x f$, to $F$ jest ciągła, $F'(x)=f(x)$ prawie wszędzie oraz $f$ jest mierzalna w sensie Lebesgue'a.
\end{twierdzenie}
\begin{proof}
Ustalmy $t\in[a,b]$ i $\eta>0$. Wybieramy cechowanie z lematu Saksa--Henstocka dla błędu $\eta/2$ i bierzemy $x$ z $|x-t|<\delta(t)$. Pojedyncza komórka o końcach $x,t$ ze znacznikiem $t$ jest drobnym podziałem częściowym. Oszacowanie modułu sumy błędów przez $\eta/2$ z lematu \ref{it:lem:kh-saks-henstock} daje
$|F(x)-F(t)-f(t)(x-t)|\le\eta/2$, niezależnie od kolejności $x,t$. Jeśli ponadto $|x-t|<\eta/(2(1+|f(t)|))$, to $|F(x)-F(t)|<\eta$. Dowodzi to ciągłości, także jednostronnej na końcach. Z lematu Saksa--Henstocka otrzymujemy następujące oszacowanie: dla ustalonego $r>0$ zbiór punktów, w których istnieją dowolnie małe przedziały $[u,v]$ zawierające $x$ i spełniające $|F(v)-F(u)-f(x)(v-u)|>r(v-u)$, ma miarę zewnętrzną zero. Istotnie, ustalmy dowolne $\varepsilon>0$ i wybierzmy cechowanie $\delta$ zapewniające błąd mniejszy od $\varepsilon$ w definicji całki na $[a,b]$. Lemat \ref{it:lem:kh-saks-henstock} ogranicza wówczas sumę modułów błędów każdego drobnego podziału częściowego przez $2\varepsilon$. Wybieramy tylko takie przedziały ze znacznikiem $x$, które leżą w $(x-\delta(x),x+\delta(x))$. Nadal istnieją dowolnie krótkie przedziały tego rodzaju przy każdym punkcie rozpatrywanego zbioru, więc tworzą drobne pokrycie Vitaliego. Stosujemy lemat o pięciokrotnym powiększeniu (twierdzenie \atlasref{ME-24::thm:5r-24}{ME-24: lemat pokryciowy 5r}) do wnętrz tych przedziałów. Otrzymujemy co najwyżej przeliczalną rodzinę przedziałów o rozłącznych wnętrzach. Ich pięciokrotne powiększenia pokrywają także końce wszystkich przedziałów z pierwotnej rodziny: w dowodzie lematu, jeżeli przedział o środku $z$ i promieniu $r_0$ przecina wybrany przedział o środku $w$ i promieniu $s_0>r_0/2$, to dla każdego $y$ z domknięcia pierwszego przedziału mamy $|y-w|\le r_0+|z-w|<2r_0+s_0<5s_0$. Dla przedziału wybranego jego domknięcie również mieści się w pięciokrotnym powiększeniu. W szczególności pokryty jest cały rozpatrywany zbiór punktów, nawet gdy znacznik leży na końcu komórki. Każdemu wybranemu przedziałowi przypisujemy jeden znacznik, dla którego zachodzi nierówność definiująca zbiór wyjątkowy. Dla każdej skończonej podrodziny lemat Saksa--Henstocka daje $r\sum_j(v_j-u_j)<2\varepsilon$. Biorąc supremum po skończonych podrodzinach, otrzymujemy $\sum_j(v_j-u_j)\le2\varepsilon/r$. Z przeliczalnej podaddytywności miary zewnętrznej wynika, że miara zewnętrzna zbioru jest nie większa niż $10\varepsilon/r$, a $\varepsilon$ jest dowolne. Stosujemy ten argument do $r=1/n$, $n\ge1$, i wykluczamy przeliczalną sumę otrzymanych zbiorów zerowych oraz końce $a,b$. Dla każdego pozostałego $x$ i każdego $n$ istnieje promień, poniżej którego żaden przedział zawierający $x$ nie narusza oszacowania z $r=1/n$. W szczególności przedział o końcach $x,x+h$ daje
\[
 \left|\frac{F(x+h)-F(x)}h-f(x)\right|\le1/n
\]
dla dostatecznie małych niezerowych $h$ obu znaków. To dowodzi $F'=f$ prawie wszędzie. Przedłużamy $F$ przez stałe poza $[a,b]$. Ciągłe funkcje $n(F(x+1/n)-F(x))$ zbiegają do $f(x)$ prawie wszędzie; ich granica górna jest mierzalna, a po zastąpieniu nieskończonych wartości zerem otrzymujemy skończoną funkcję mierzalną równą $f$ prawie wszędzie. Zupełność miary Lebesgue'a daje mierzalność $f$.
\end{proof}
\begin{uwaga}[Dlaczego ograniczamy każdą sumę częściową]
Nie wystarcza samo oszacowanie modułu sumy wszystkich błędów: dodatnie i ujemne błędy mogą się znosić. Dopiero możliwość wybrania dowolnego podzbioru komórek w lemacie \ref{ma1:kh:lem:hk-uzupelnianie} pozwala otrzymać sumę modułów. Jest ona potrzebna w dowodzie całkowania przez części i przy własnościach funkcji pierwotnych.
\end{uwaga}
\subsection{Całki niewłaściwe i osobliwości}
\noindent\emph{Porównanie w Bibliotece matematycznej:} \atlasref{AM1-36::ma1:kh:thm:kh-hake}{AM1-36: twierdzenie Hake'a}.

\begin{twierdzenie}[Hake'a: związek z całkami niewłaściwymi]\label{ma1:kh:thm:kh-hake}\label{it:thm:kh-hake}
Niech $f:[a,b]\to\R$ ma skończone wartości. Wówczas $f\in\HK[a,b]$ wtedy i tylko wtedy, gdy $f\in\HK[c,b]$ dla każdego $c\in(a,b)$ oraz istnieje skończona granica $L=\lim_{c\downarrow a}\int_c^b f$. W tym przypadku $\int_a^b f=L$, niezależnie od wartości $f(a)$.
\end{twierdzenie}
\begin{proof}
Jeśli $f\in\HK[a,b]$, lokalność daje całkowalność na $[c,b]$ i $\int_c^b f=F(b)-F(c)$, gdzie $F(x)=\int_a^x f$. Ciągłość $F$ w $a$ i $F(a)=0$ dają granicę $F(b)$. Dla dostateczności definiujemy $F(a)=0$ i $F(x)=L-\int_x^b f$ dla $x>a$. Funkcja jest ciągła, a jej przyrosty na przedziałach oddalonych od $a$ są lokalnymi całkami. Wybieramy ciąg $b=c_0>c_1>\dots\downarrow a$. Na $J_n=[c_n,c_{n-1}]$ wybieramy cechowanie $\delta_n$ zapewniające błąd definicyjny mniejszy od $\varepsilon2^{-n-3}$. Lemat Saksa--Henstocka daje wówczas sumę modułów błędów na podziale częściowym nie większą niż $\varepsilon2^{-n-2}$. Dla $t\in(c_n,c_{n-1})$ bierzemy minimum $\delta_n(t)$ oraz połowy odległości od obu końców. W punkcie $c_n$, $n\ge1$, bierzemy minimum $\delta_n(c_n)$, $\delta_{n+1}(c_n)$ i połowy odległości od $c_{n-1},c_{n+1}$. W punkcie $b$ wystarczy minimum $\delta_1(b)$ i $(b-c_1)/2$. Komórka przecinająca $c_n$ musi mieć znacznik $c_n$ i nie może sięgać kolejnego punktu podziału; po rozcięciu obie części są drobne względem lokalnych cechowań. W $a$ wybieramy $\delta(a)>0$ z $|F(y)|+|f(a)|(y-a)<\varepsilon/2$ dla $a<y<a+\delta(a)$, a poza $a$ zmniejszamy cechowania tak, aby przedziały z innymi znacznikami nie zawierały $a$. Przedział początkowy dowolnego drobnego podziału ma znacznik $a$, pozostałe można pogrupować według $J_n$ (i ewentualnie rozciąć w $c_n$). Rozcięcie nie zmienia sumy Riemanna ani sumy przyrostów $F$. Na każdym $J_n$ otrzymujemy podział częściowy, więc suma modułów błędów nie przekracza $\varepsilon2^{-n-2}$. Komórka początkowa wnosi błąd mniejszy od $\varepsilon/2$. Ponieważ $\sum_{n\ge1}2^{-n-2}=1/4$, całkowity błąd jest mniejszy od $\varepsilon$. Suma przyrostów wynosi $F(b)-F(a)=L$.
\end{proof}
\begin{przyklad}[Szeregi i całkowalność]
Dla $I_n=(2^{-n},2^{-n+1}]$ i ciągu $(a_n)_{n\ge1}$ kładziemy $f(x)=2^n a_n$ na $I_n$, a $f(0)=0$. Wówczas
\[
f\in\HK[0,1]\iff\sum_{n\ge1}a_n\ \text{jest zbieżny},\qquad
f\in\mathcal L^1[0,1]\iff\sum_{n\ge1}|a_n|<\infty.
\]
Na każdym $[c,1]$, $c>0$, funkcja jest schodkowa z co najwyżej skończoną liczbą skoków. Dla $c\in I_n$ całka po $[c,1]$ jest równa $\sum_{k=1}^{n-1}a_k+2^na_n(2^{-n+1}-c)$, przy czym ostatni składnik ma moduł co najwyżej $|a_n|$. Gdy szereg jest zbieżny, $a_n\to0$, więc ta całka ma granicę równą sumie szeregu. Odwrotnie, biorąc $c=2^{-n}$ otrzymujemy sumy częściowe szeregu. Twierdzenie Hake'a daje pierwszą równoważność; druga z $\int_{I_n}|f|\,\mathrm d\lambda_1=|a_n|$. Ponadto $f$ jest całkowalna Riemanna dokładnie wtedy, gdy $\sup_n2^n|a_n|<\infty$: konieczność wynika z ograniczoności funkcji całkowalnej Riemanna, a przy skończonym supremum funkcja jest ograniczona i jej punkty nieciągłości zawierają się w przeliczalnym zbiorze $\{0\}\cup\{2^{-n}:n\ge0\}$. Kryterium Lebesgue'a całkowalności Riemanna daje dostateczność. Dla $a_n=(-1)^{n+1}/n$ otrzymujemy $\int_0^1f=\log2$ w sensie HK, lecz $\int_0^1|f|\,\mathrm d\lambda_1=\infty$.
\end{przyklad}
\begin{przyklad}[Nie każda osobliwość jest całkowalna]
Funkcja $f(x)=1/x$ na $(0,1]$, niezależnie od skończonej wartości w $0$, nie jest całkowalna HK na $[0,1]$, bo $\int_c^1\mathrm dx/x=-\log c\to\infty$. Na $[-1,1]$ funkcja $1/x$ z dowolną skończoną wartością w $0$ ma wartość główną Cauchy'ego równą $0$, ale nie ma całki HK, gdyż jej ograniczenie do $[0,1]$ nie jest HK-całkowalne.
\end{przyklad}
\begin{twierdzenie}[Rodzina całek oscylacyjnych]\label{ma1:kh:tw:hk-oscylacyjne}\label{tw:hk-oscylacyjne}
Niech $p\in\R$, $q>0$ i niech $f_{p,q}:[0,1]\to\R$ będzie określona wzorem $f_{p,q}(0)=0$ oraz $f_{p,q}(x)=x^{-p}\sin(x^{-q})$ dla $x\in(0,1]$. Wówczas
\[
 f_{p,q}\in\HK[0,1]\quad\Longleftrightarrow\quad p<q+1,
 \qquad f_{p,q}\in\mathcal L^1[0,1]\quad\Longleftrightarrow\quad p<1.
\]
W szczególności dla $1\le p<q+1$ całka HK istnieje, lecz nie jest bezwzględnie zbieżna.
\end{twierdzenie}
\begin{proof}
Na $[\varepsilon,1]$ funkcja jest ciągła. Podstawiając $t=x^{-q}$, uzyskujemy
\[
\int_\varepsilon^1x^{-p}\sin(x^{-q})\,\mathrm dx
=\frac1q\int_1^{\varepsilon^{-q}}t^{(p-1)/q-1}\sin t\,\mathrm dt.
\]
Oznaczmy wykładnik przez $\beta=(p-1)/q-1$. Gdy $\beta<0$, funkcja $t^\beta$ maleje do zera. Dla $1\le A<B$ całkowanie przez części daje
\[
\int_A^B t^\beta\sin t\,\mathrm dt
=[-t^\beta\cos t]_A^B+\beta\int_A^B t^{\beta-1}\cos t\,\mathrm dt.
\]
Moduł prawej strony jest nie większy niż $A^\beta+B^\beta+|\beta|\int_A^B t^{\beta-1}\,\mathrm dt=2A^\beta$. Zatem kryterium Cauchy'ego daje zbieżność całki niewłaściwej. Gdy $\beta=0$, pierwotna $-\cos t$ nie ma granicy. Jeśli $\beta>0$, na półokresie $[k\pi,(k+1)\pi]$, $k\ge1$, sinus ma stały znak. Stąd
\[
 \left|\int_{k\pi}^{(k+1)\pi}t^\beta\sin t\,\mathrm dt\right|
 \ge 2(k\pi)^\beta,
\]
gdyż całka z $|\sin t|$ po tym półokresie wynosi $2$. Narusza to kryterium Cauchy'ego całki niewłaściwej. Twierdzenie Hake'a daje pierwszą równoważność, ponieważ $\beta<0$ dokładnie wtedy, gdy $p<q+1$.

Dla $\beta<-1$ całka modułu jest skończona przez oszacowanie $t^\beta|\sin t|\le t^\beta$ dla $t\ge1$. Dla $-1\le\beta<0$ mamy na każdym półokresie
\[
 \int_{k\pi}^{(k+1)\pi}t^\beta|\sin t|\,\mathrm dt
 \ge 2((k+1)\pi)^\beta.
\]
Szereg tych dolnych ograniczeń jest rozbieżny, ponieważ $(k+1)^\beta\ge 1/(k+1)$. Dla $\beta\ge0$ dolne ograniczenie wynosi $2(k\pi)^\beta\ge2$ i również daje rozbieżność. Zatem bezwzględna całkowalność zachodzi dokładnie dla $\beta<-1$, czyli $p<1$.
\end{proof}
\begin{przyklad}[Dokładny próg przy całce $x^{-2}\sin(x^{-2})$]
Dla $p=q=2$ mamy $2<3$, więc całka $\int_0^1x^{-2}\sin(x^{-2})\,\mathrm dx$ istnieje w sensie HK, ale nie w sensie Lebesgue'a. Podstawienie z dowodu daje użyteczną reprezentację obliczeniową
$\int_0^1x^{-2}\sin(x^{-2})\,\mathrm dx=\tfrac12\int_1^\infty t^{-1/2}\sin t\,\mathrm dt$.
\end{przyklad}
\subsection{Podstawienie oraz całkowanie przez części}
\begin{twierdzenie}[Mnożniki o ograniczonej wariacji]\label{ma1:kh:thm:bv-multiplier}
Niech $f\in\HK[a,b]$, $F(x)=\int_a^x f$, a $g:[a,b]\to\R$ ma ograniczoną wariację. Wtedy $fg\in\HK[a,b]$ oraz
\[
\int_a^b f(x)g(x)\,\mathrm dx=F(b)g(b)-F(a)g(a)-\int_a^b F\,\mathrm dg,
\]
gdzie ostatnia całka jest zwykłą całką Riemanna--Stieltjesa.
\end{twierdzenie}
\begin{proof}
Najpierw uzasadnimy istnienie całki Riemanna--Stieltjesa. Dla podziału $P$ ze znacznikami definiujemy sumę $\sum_i F(t_i)(g(x_i)-g(x_{i-1}))$. Niech $V=\Var_{[a,b]}g$ oraz $\omega_F(s)=\sup_{|x-y|\le s}|F(x)-F(y)|$ dla $x,y\in[a,b]$. Jeśli $P'$ jest zagęszczeniem $P$, dowolna suma na $P'$ różni się od danej sumy na $P$ o co najwyżej $\omega_F(\mesh P)V$: rozpisujemy przyrosty $g$ jako sumy przyrostów na mniejszych komórkach i stosujemy nierówność trójkąta. Dla dwóch podziałów o średnicach mniejszych od $s$ porównanie ze wspólnym zagęszczeniem daje różnicę nie większą niż $2\omega_F(s)V$. Jednostajna ciągłość $F$ daje $\omega_F(s)\to0$. Wybierając dowolny ciąg podziałów o średnicach dążących do zera, otrzymujemy ciąg Cauchy'ego sum, a porównanie z jego granicą dowodzi istnienia wspólnej granicy wszystkich sum. Nazywamy ją całką Riemanna--Stieltjesa $\int_a^b F\,\mathrm dg$. Dla $\delta$ z lematu Saksa--Henstocka i $\Delta F_i=F(x_i)-F(x_{i-1})$ mamy
$|S(fg,\mathcal P)-\sum_i g(t_i)\Delta F_i|\le2\varepsilon\|g\|_\infty$.
Ponieważ $\sum_i|g(t_i)-g(x_{i-1})|\le \Var_{[a,b]}g$ i $|\Delta F_i|\le\omega_F(\mesh P)$, różnica między $\sum_i g(t_i)\Delta F_i$ i $\sum_i g(x_{i-1})\Delta F_i$ nie przekracza $\omega_F(\mesh P)\Var_{[a,b]}g$. Dokładne sumowanie przez części daje
$\sum_i g(x_{i-1})\Delta F_i=F(b)g(b)-F(a)g(a)-\sum_iF(x_i)(g(x_i)-g(x_{i-1}))$.
Ostatnia suma dąży do $\int F\,\mathrm dg$. Bierzemy $\varepsilon$ małe i dodatkowo ograniczamy cechowanie stałą zapewniającą dostatecznie małą średnicę podziałów.
\end{proof}
\begin{uwaga}[Ograniczenia całki HK]
Klasa $\HK[a,b]$ nie jest kratą: z całkowalności $f$ nie wynika całkowalność $|f|$. Gdyby moduł pochodnej z przykładu o funkcji $x^2\cos(\pi/x^2)$ był HK-całkowalny, to na mocy twierdzenia \ref{it:thm:hk-absolute-integrability} byłby całkowalny w sensie Lebesgue'a, wbrew wykazanej rozbieżności jego całki. Nie jest też algebrą, ponieważ $x^{-1/2}$ na $(0,1]$ po odpowiednim określeniu w zerze ma całkę HK, a jej kwadrat $1/x$ nie. Seminorma Alexiewicza $\|f\|_A=\sup_{x\in[a,b]}\left|\int_a^x f\right|$ jest normą po utożsamieniu funkcji równych prawie wszędzie, lecz otrzymana przestrzeń nie jest zupełna (dowód i opis uzupełnienia podano w stwierdzeniu o normie Alexiewicza w niniejszym rozdziale).
\end{uwaga}
\begin{twierdzenie}[Podstawienie przez dyfeomorfizm klasy $C^1$]\label{ma1:kh:tw:hk-podstawienie}\label{tw:hk-podstawienie}
Niech $a<b$, $u<v$, $f\in\HK[a,b]$ i niech $\phi:[u,v]\to[a,b]$ będzie rosnącym dyfeomorfizmem klasy $C^1$, to znaczy $\phi\in C^1[u,v]$, $\phi(u)=a$, $\phi(v)=b$ i $\phi'(t)>0$ dla $t\in[u,v]$. Wówczas $t\mapsto f(\phi(t))\phi'(t)$ jest HK-całkowalna na $[u,v]$ oraz
\[
\int_u^v f(\phi(t))\phi'(t)\,\mathrm dt=\int_a^b f(x)\,\mathrm dx.
\]
Dla malejącego dyfeomorfizmu z $\phi(u)=b$, $\phi(v)=a$ i $\phi'<0$ całka po prawej ma odwrócone granice.
\end{twierdzenie}
\begin{proof}
Niech $F(x)=\int_a^x f$. Funkcja $F$ jest ciągła i $F'=f$ prawie wszędzie. Najpierw uzasadnimy, że $F\circ\phi$ należy do $\AC G_*[u,v]$. Z charakteryzacji pierwotnych HK (twierdzenie \ref{it:thm:kh-acgstar-characterization} w rozdziale o wąskiej całce Denjoya) mamy $[a,b]=\bigcup_n E_n$, gdzie zbiory $E_n$ są domknięte i $F\in\AC_*(E_n)$. Zbiory $\phi^{-1}(E_n)$ są domknięte i pokrywają $[u,v]$. Ponieważ $\phi$ jest lipschitzowska, dla rozłącznych przedziałów o końcach w $\phi^{-1}(E_n)$ i o małej łącznej długości obrazy mają także małą sumę długości i rozłączne wnętrza. Równość
$\osc_J(F\circ\phi)=\osc_{\phi(J)}F$ dowodzi $F\circ\phi\in\AC_*(\phi^{-1}(E_n))$.

Dyfeomorfizm i jego odwrotność są lipschitzowskie na zwartych przedziałach, więc przeciwobraz zbioru miary zero ma miarę zero. W punktach, w których $F'(\phi(t))=f(\phi(t))$, zwykła reguła łańcuchowa daje $(F\circ\phi)'(t)=f(\phi(t))\phi'(t)$. Twierdzenie \ref{it:thm:acgstar-derivative-integrable} daje całkowalność HK tej pochodnej i wzór Newtona--Leibniza; ostatni przyrost wynosi $F(b)-F(a)$. Dla malejącego dyfeomorfizmu argument jest taki sam: obrazy komórek zapisujemy z końcami w kolejności rosnącej, ich długości kontroluje $\max|\phi'|$, a odwrotność jest lipschitzowska z oszacowania $\min|\phi'|>0$. Przyrost funkcji złożonej wynosi tym razem $F(a)-F(b)$.
\end{proof}
\begin{uwaga}[Zakres twierdzenia]
Warunek $\phi'(t)>0$ zabezpiecza przed trudnościami z przenoszeniem zbiorów wyjątkowych. Istnieją mocniejsze twierdzenia o podstawieniu dla mniej regularnych funkcji $\phi$, ale bez dodatkowych warunków na składanie pierwotnych i na zbiory wyjątkowe nie wolno automatycznie zastępować dyfeomorfizmu dowolną funkcją bezwzględnie ciągłą.
\end{uwaga}
\begin{wniosek}[Całkowanie przez części z mnożnikiem bezwzględnie ciągłym]\label{ma1:kh:wn:hk-ibp-ac}\label{wn:hk-ibp-ac}
Niech $f\in\HK[a,b]$, $F(x)=\int_a^x f$, a $g\in\AC[a,b]$. Wtedy $fg\in\HK[a,b]$ i
\[
\int_a^b f(x)g(x)\,\mathrm dx=F(b)g(b)-F(a)g(a)-\int_a^b F(x)g'(x)\,\mathrm d\lambda_1(x).
\]
Ostatnia całka jest całką Lebesgue'a, ponieważ $F$ jest ograniczona i $g'\in L^1[a,b]$.
\end{wniosek}
\begin{proof}
Funkcja $g$ ma wahanie ograniczone, gdyż $\Var(g)\le\int_a^b|g'|\,\mathrm d\lambda_1$. Stosujemy twierdzenie o mnożnikach z tej sekcji. Z twierdzenia \ref{ma1:kh:thm:ac-lebesgue-primitive} mamy $g(v)-g(u)=\int_u^v g'\,\mathrm d\lambda_1$. Różnica sumy Stieltjesa i całki Lebesgue'a jest zatem ograniczona przez
\[
 \sum_i\int_{x_{i-1}}^{x_i}|F(t_i)-F(x)|\,|g'(x)|\,\mathrm d\lambda_1(x)
 \le \omega_F(\mesh P)\int_a^b|g'|\,\mathrm d\lambda_1.
\]
Prawa strona dąży do zera, co daje żądaną identyfikację.
\end{proof}
\begin{przyklad}[Warunkowo zbieżna całka niewłaściwa]
Niech $f(x)=\sin(1/x)/x$ dla $x\in(0,1]$ i $f(0)=0$. Podstawienie $t=1/x$ na $[\varepsilon,1]$ daje
$\int_\varepsilon^1f(x)\,\mathrm dx=\int_1^{1/\varepsilon}(\sin t)/t\,\mathrm dt$.
Całka po prawej jest zbieżna na mocy kryterium Dirichleta, zatem twierdzenie Hake'a daje
$\int_0^1f(x)\,\mathrm dx=\int_1^\infty (\sin t)/t\,\mathrm dt$ w sensie HK. Nie zachodzi całkowalność Lebesgue'a wartości bezwzględnej, gdyż $\int_1^\infty|\sin t|/t\,\mathrm dt=\infty$.
\end{przyklad}
\subsection{Twierdzenia o zbieżności}
\begin{twierdzenie}[Zbieżność z majorantą Lebesgue'a]\label{ma1:kh:tw:hk-zbieznosc-dct}\label{tw:hk-zbieznosc-dct}
Niech $f_n:[a,b]\to\R$ będą mierzalne i HK-całkowalne, $f_n\to f$ prawie wszędzie dla $f:[a,b]\to\R$, a $g\in L^1[a,b]$ niech spełnia $|f_n|\le g$ prawie wszędzie dla każdego $n$. Wówczas $f,f_n\in\mathcal L^1[a,b]$ i
$\lim_{n\to\infty}\int_a^bf_n\,\mathrm dx=\int_a^bf\,\mathrm dx$, przy czym po obu stronach całkę można rozumieć zarówno w sensie Lebesgue'a, jak i HK.
\end{twierdzenie}
\begin{proof}
Wybieramy skończonego mierzalnego reprezentanta $g$ klasy $L^1$. Poza jednym zbiorem miary zero (sumą zbiorów wyjątkowych dla wszystkich nierówności i zbieżności) zachodzą jednocześnie $|f_n|\le g$ i $f_n\to f$, stąd $|f|\le g$. Funkcja $f$ jest mierzalna, ponieważ poza tym zbiorem jest granicą funkcji mierzalnych, a miara Lebesgue'a jest zupełna. Oszacowania przez $g$ dają bezwzględną całkowalność Lebesgue'a każdej $f_n$ i $f$. Stosujemy twierdzenie \ref{thm:dct} oraz zgodność całek Lebesgue'a i HK (twierdzenie \ref{it:thm:mcshane-lebesgue} i inkluzja całki McShane'a w HK). Nieujemna majoranta HK-całkowalna jest automatycznie całkowalna w sensie Lebesgue'a na mocy kryterium bezwzględnej całkowalności. Dwustronne ograniczenie przez funkcje HK-całkowalne, które same nie muszą należeć do $L^1$, oraz zbieżność monotoniczną omawiamy w niniejszym rozdziale (twierdzenie \ref{it:thm:hk-limit-theorems}).
\end{proof}
\begin{przyklad}[Dlaczego sama zbieżność punktowa nie wystarcza]
Niech $f_n(0)=0$ i $f_n(x)=n\chi_{(0,1/n]}(x)$ na $[0,1]$. Każda $f_n$ jest całkowalna w sensie Riemanna, a więc HK-całkowalna, oraz $f_n(x)\to0$ w każdym punkcie $[0,1]$. Mimo to $\int_0^1f_n\,\mathrm dx=1$ dla każdego $n$. Nie istnieje jedna całkowalna w sensie Lebesgue'a funkcja $g$ z $|f_n|\le g$ prawie wszędzie dla wszystkich $n$: w przeciwnym razie twierdzenie o zbieżności zmajoryzowanej wymuszałoby granicę całek równą zeru.
\end{przyklad}
\begin{stwierdzenie}[Jednostajna zbieżność]\label{ma1:kh:st:hk-zbieznosc-jednostajna}\label{st:hk-zbieznosc-jednostajna}
Jeśli $f_n\in\HK[a,b]$ i $f_n\to f$ jednostajnie, to $f\in\HK[a,b]$ oraz $\int_a^bf_n\to\int_a^bf$.
\end{stwierdzenie}
\begin{proof}
Dla dwóch funkcji HK-całkowalnych $h,k$ o ograniczonej różnicy zachodzi $|\int_a^b(h-k)|\le(b-a)\|h-k\|_\infty$, ponieważ nierówność dla sum przechodzi na granicę. Dla dostatecznie dużych $n,m$ różnica $f_n-f_m$ jest ograniczona i jej norma supremum dąży do zera. Stąd całki $I_n=\int_a^bf_n$ tworzą ciąg Cauchy'ego. Dla ustalonego $\varepsilon>0$ wybieramy $n$ z $(b-a)\|f-f_n\|_\infty<\varepsilon/3$ oraz $|I-I_n|<\varepsilon/3$, gdzie $I=\lim I_n$. Następnie bierzemy cechowanie dla $f_n$ zapewniające błąd sumy mniejszy niż $\varepsilon/3$. Nierówność trójkąta kończy dowód bez wcześniejszego zakładania całkowalności $f$.
\end{proof}


\begin{twierdzenie}[Jednostajna całkowalność względem cechowania]
\label{ma1:kh:thm:kh-equiintegrability}\label{it:thm:kh-equiintegrability}
Niech $f_n\in\HK[a,b]$ i $f_n(t)\to f(t)\in\R$ dla każdego $t\in[a,b]$. Załóżmy, że dla każdego $\varepsilon>0$ istnieje jedno cechowanie $\delta$ takie, że
\[
 |S(f_n,\mathcal P)-I_n|<\varepsilon
 \quad\text{dla wszystkich }n\ge1
 \text{ i wszystkich }\delta\text{-drobnych }\mathcal P,
 \qquad I_n=\int_a^b f_n.
\]
Wtedy $f\in\HK[a,b]$ oraz $I_n\to\int_a^b f$.
\end{twierdzenie}
\begin{proof}
Dla $\varepsilon>0$ wybieramy wspólne cechowanie dla błędu $\varepsilon/3$ i ustalamy jeden drobny podział $\mathcal P$ z lematu Cousina. Ma on skończenie wiele znaczników, więc $S(f_n,\mathcal P)\to S(f,\mathcal P)$ dzięki zbieżności punktowej. Dla dostatecznie dużych $n,m$ różnica tych dwóch sum ma moduł mniejszy od $\varepsilon/3$. Nierówność trójkąta daje $|I_n-I_m|<\varepsilon$. Zupełność $\R$ zapewnia granicę $I$. Teraz dla dowolnej dokładności $\eta>0$ bierzemy wspólne cechowanie dla błędu $\eta/2$. Dla każdego jego drobnego podziału $\mathcal Q$ przechodzimy z $n$ do nieskończoności w nierówności $|S(f_n,\mathcal Q)-I_n|<\eta/2$. Otrzymujemy $|S(f,\mathcal Q)-I|\le\eta/2<\eta$, zatem $f$ jest całkowalna z całką $I$. Nie zakładaliśmy z góry istnienia całki z granicy.
\end{proof}

\begin{uwaga}[Całka Henstocka--Kurzweila--Stieltjesa]
Zastępując w sumach drobnych $x_i-x_{i-1}$ przez $\alpha(x_i)-\alpha(x_{i-1})$, otrzymujemy całkę HK--Stieltjesa. W szczególności dla $f=\alpha=\ind{[c,b]}$, $a<c<b$, całka istnieje i wynosi $1$. Cechowanie można dobrać tak, aby każdy przedział o niezerowym przyroście $\alpha$ miał znacznik $c$. Reguła całkowania przez części wymaga tu poprawek na wspólne skoki.
\end{uwaga}

\subsection{Norma Alexiewicza i twierdzenia graniczne dla HK}
\begin{stwierdzenie}[Uzupełnienie w normie Alexiewicza]\label{it:thm:alexiewicz-completion}
Niech $a<b$. Przestrzeń funkcji HK-całkowalnych modulo równość prawie wszędzie, z normą $\|[f]\|_A=\sup_{x\in[a,b]}|(\HK)\int_a^x f|$, nie jest zupełna. Jej uzupełnienie izometryczne można utożsamić z $C_0[a,b]=\{F\in C[a,b]:F(a)=0\}$ z normą supremum (symbol $C_0[a,b]$ w tym stwierdzeniu oznacza funkcje zerujące się w $a$, nie przestrzeń funkcji znikających w nieskończoności).
\end{stwierdzenie}
\begin{proof}
Przyporządkowanie $f\mapsto F(x)=(\HK)\int_a^x f$ przyjmuje wartości w $C_0[a,b]$: pierwotna jest ciągła, a jej wartość w $a$ wynosi zero. Wynika to z twierdzenia \ref{it:thm:kh-primitive-derivative}. Jest liniowe dzięki liniowości i lokalności HK, a jego jądro stanowią dokładnie funkcje zerowe prawie wszędzie: jeden kierunek wynika z lematu \ref{ma1:kh:lem:kh-null-modification}, drugi z twierdzenia \ref{ma1:kh:thm:kh-primitive-derivative}. Na przestrzeni ilorazowej jest zatem różnowartościową izometrią, ponieważ $\|F\|_\infty=\|f\|_A$. W szczególności $\|\cdot\|_A$ jest na niej normą. Jego obraz zawiera wszystkie funkcje przedziałami liniowe zerujące się w $a$: dla skończonego podziału wybieramy funkcję podcałkową równą nachyleniu na wnętrzu każdej komórki, a na końcach przyjmujemy dowolne skończone wartości. Jest to ograniczona mierzalna funkcja całkowalna w sensie Lebesgue'a, więc z twierdzenia \ref{it:thm:mcshane-lebesgue} i inkluzji McShane'a w HK ma tę samą całkę w sensie HK. Dokładniej, dla węzłów $a=x_0<\cdots<x_m=b$ oraz zadanej funkcji przedziałami liniowej $P$ o $P(a)=0$ wybieramy $p_j=(P(x_j)-P(x_{j-1}))/(x_j-x_{j-1})$. Dla $x\in[x_{k-1},x_k]$ zgodność całek na podprzedziałach i zerowa miara końców dają
\[
 (\HK)\int_a^x f
 =\sum_{j<k}p_j(x_j-x_{j-1})+p_k(x-x_{k-1})
 =P(x_{k-1})-P(a)+P(x)-P(x_{k-1})=P(x).
\]
 Dla $F\in C_0[a,b]$ i $\varepsilon>0$ z jednostajnej ciągłości $F$ na zwartym przedziale wybieramy podział o tak małej średnicy, że oscylacja $F$ na każdej komórce jest mniejsza od $\varepsilon$. Interpolant liniowy $P$ wartości $F$ na końcach komórek zeruje się w $a$. Dla $x$ w komórce $[u,v]$ mamy $P(x)=(1-\theta)F(u)+\theta F(v)$, gdzie $\theta=(x-u)/(v-u)$, zatem $|P(x)-F(x)|\le (1-\theta)|F(u)-F(x)|+\theta|F(v)-F(x)|<\varepsilon$. Domknięciem obrazu jest więc całe $C_0[a,b]$. Nie każda funkcja z tej przestrzeni jest pierwotną HK: funkcja $x\mapsto C((x-a)/(b-a))$ dla funkcji Cantora $C$ z przykładu \ref{ex:cantor-measure-17} jest ciągła, zeruje się w $a$ i przyjmuje wartość $1$ w $b$. Poza afinicznym przeciwobrazem zbioru Cantora jest lokalnie stała, a więc ma pochodną zero. Przeciwobraz tego zbioru ma miarę zero na mocy \atlasref{ME-13::eq:linear-volume-13}{ME-13: wzór na zmianę objętości}. Gdyby była pierwotną HK funkcji $f$, twierdzenie \ref{it:thm:kh-primitive-derivative} dawałoby $f=0$ prawie wszędzie, więc lemat \ref{it:lem:kh-null-modification} zastosowany do ograniczenia $f$ na każdym $[a,x]$, $x>a$, wymuszałby zerową pierwotną (w $a$ również równą zero). Sprzeczność. Z twierdzenia o zupełności $C(K)$ (stwierdzenie \ref{prop:ck-banach12}) przestrzeń $C[a,b]$ jest przestrzenią Banacha z normą supremum. Ewaluacja $F\mapsto F(a)$ jest liniowa i ciągła, ponieważ $|F(a)|\le\|F\|_\infty$. Jej jądro $C_0[a,b]$ jest zatem domknięte i jako domknięta podprzestrzeń przestrzeni Banacha jest zupełne. Obraz jest gęsty, lecz nie jest całym $C_0[a,b]$, a więc nie jest domknięty i nie jest zupełny.
\end{proof}
\begin{twierdzenie}[Twierdzenia graniczne dla HK]\label{it:thm:hk-limit-theorems}\label{ma1:kh:thm:kh-order-limits}
Niech $(f_n)$ będzie ciągiem rzeczywistych funkcji
HK-całkowalnych na $[a,b]$, a $f:[a,b]\to\R$.
W przypadku granicy określonej tylko prawie wszędzie
wybieramy dowolne skończone wartości na zbiorze
wyjątkowym. Wszystkie całki bez oznaczenia miary
w poniższych trzech warunkach rozumiemy w sensie HK.
\begin{enumerate}[label=(\alph*)]
\item Jeśli $f_n\uparrow f$ prawie wszędzie i $\sup_n(\HK)\int_a^b f_n<\infty$, to $f$ jest HK-całkowalna oraz $\int_a^b f_n\to\int_a^b f$.
\item Jeśli $f_n\to f$ prawie wszędzie i istnieją $g,h\in\HK[a,b]$ takie, że $g\le f_n\le h$ prawie wszędzie dla wszystkich $n$, to $f\in\HK[a,b]$ i $\int_a^b f_n\to\int_a^b f$.
\item Jeśli $f_n\downarrow f$ prawie wszędzie i
$\inf_n(\HK)\int_a^b f_n>-\infty$, to $f\in\HK[a,b]$
i całki zbiegają do całki $f$.
\end{enumerate}
We wszystkich trzech przypadkach $f_n-f\in\mathcal L^1$
oraz dla pierwotnych $F_n(x)=(\HK)\int_a^x f_n$,
$F(x)=(\HK)\int_a^x f$ zachodzi
\[
 \|[f_n-f]\|_A
 =\|F_n-F\|_\infty
 \le V_a^b(F_n-F)
 =\int_{[a,b]}|f_n-f|\,\mathrm d\lambda_1\longrightarrow0.
\]
Nie zakładamy, że same $f_n$ lub $f$ są całkowalne
w sensie Lebesgue'a.
\end{twierdzenie}
\begin{proof}
W każdej z trzech części zbieramy wszystkie zbiory
wyjątkowe (także dla poszczególnych $n$) w jeden
mierzalny zbiór miary zero. Na tym zbiorze zastępujemy
wszystkie rozpatrywane funkcje zerem. Z lematu
\ref{ma1:kh:lem:kh-null-modification} wynika, że taka
zmiana zachowuje całkowalność HK i wartość całki
w obie strony: jeżeli dopiero zmodyfikowana funkcja
graniczna okaże się całkowalna, ten sam lemat
pozwoli przywrócić jej pierwotne wartości.
Możemy zatem zakładać punktowe nierówności
i punktową zbieżność na całym przedziale.

(a) Funkcje $u_n=f_n-f_1$ są nieujemne i HK-całkowalne, zatem należą do $\mathcal L^1$ na mocy twierdzenia \ref{it:thm:hk-absolute-integrability}. Rosną punktowo do $u=f-f_1$, więc $u$ jest mierzalna. Z liniowości HK i zgodności z całką Lebesgue'a wynika
\[
 \sup_n\int_{[a,b]}u_n\,\mathrm d\lambda_1
 =\sup_n(\HK)\int_a^b f_n-(\HK)\int_a^b f_1<\infty.
\]
Twierdzenie \ref{thm:mct} daje $u\in\mathcal L^1$ i zbieżność całek. Ponieważ $f=f_1+u$, liniowość HK daje całkowalność $f$ oraz żądaną zbieżność, po dodaniu $(\HK)\int_a^b f_1$.

(b) Funkcja $h-g$ jest nieujemna i HK-całkowalna, więc należy do $\mathcal L^1$ z twierdzenia \ref{it:thm:hk-absolute-integrability}. Także $u_n=f_n-g$ jest nieujemna, HK-całkowalna i ograniczona przez $h-g$. Jej granica punktowa $u=f-g$ jest mierzalna. Twierdzenie \ref{thm:dct} i wniosek \ref{cor:dct-l1} dają $u\in\mathcal L^1$ i $\int_{[a,b]}|u_n-u|\,\mathrm d\lambda_1\to0$. Z $f=g+u$ otrzymujemy całkowalność HK i zbieżność całek. W obu częściach przywrócenie początkowych wartości na zbiorze zerowym jest dozwolone przez lemat \ref{it:lem:kh-null-modification}. Ponadto w (a) $\int|u-u_n|\,\mathrm d\lambda_1=\int u\,\mathrm d\lambda_1-\int u_n\,\mathrm d\lambda_1\to0$, a w (b) zbieżność w $L^1$ już wykazano. W obu przypadkach $f_n-f=u_n-u$ należy do $\mathcal L^1$,
także po przywróceniu wartości na zbiorze zerowym,
gdyż miara Lebesgue'a jest zupełna.
W (c) stosujemy (a) do $-f_n\uparrow-f$; warunek
ograniczenia całek od dołu dla $f_n$ jest dokładnie
warunkiem ograniczenia ich przeciwieństw od góry.
Otrzymujemy także $\int|f_n-f|\,\mathrm d\lambda_1\to0$.

W każdym przypadku zgodność całek Lebesgue'a i HK
(twierdzenie \ref{it:thm:mcshane-lebesgue} i inkluzja
McShane'a w HK) stosujemy do różnicy na każdym
podprzedziale $[a,x]$. Zatem
$F_n(x)-F(x)=\int_{[a,x]}(f_n-f)\,\mathrm d\lambda_1$.
Jest to funkcja bezwzględnie ciągła na mocy
twierdzenia \atlasref{ME-25::thm:FTC-lebesgue-25}{ME-25: podstawowe twierdzenie rachunku całkowego}, zerująca się w $a$,
a w szczególności ma ograniczone wahanie, jak wykazano
w dowodzie twierdzenia \ref{it:thm:ac-lebesgue-primitive}.
Jej miarą Stieltjesa ze znakiem jest
$(f_n-f)\lambda_1|_{(a,b]}$: obie miary mają
te same przyrosty na przedziałach $(s,t]$.
Z jednoznaczności i równości wahania w twierdzeniu
\ref{it:thm:stieltjes-signed-bv} otrzymujemy
$V_a^b(F_n-F)=\int_{[a,b]}|f_n-f|\,\mathrm d\lambda_1$.
Istotnie wahanie miary o rzeczywistej gęstości $q$
jest miarą $|q|\lambda_1$: górne oszacowanie wynika
z nierówności trójkąta na każdym mierzalnym podziale,
a dolne z podziału na $\{q\ge0\}$ i $\{q<0\}$.
Ponieważ $F_n(a)-F(a)=0$, norma supremum różnicy
nie przekracza jej wahania. To daje wszystkie
deklarowane równości i oszacowania.
\end{proof}
\begin{przyklad}[Konieczność założeń]
Dla $f_n=n\chi_{(0,1/n)}$ mamy $f_n\to0$ punktowo na $[0,1]$, lecz $\int_0^1f_n=1$ dla wszystkich $n$; brak wspólnej całkowalnej majoranty. Z kolei $g_n(x)=\min\{n,1/x\}$ dla $x>0$, $g_n(0)=0$, rośnie punktowo do funkcji równej $1/x$ na $(0,1]$, a $0$ w zerze, i $\int_0^1g_n=1+\log n\to\infty$. Sama monotoniczność bez skończonej granicy całek nie wystarcza.
\end{przyklad}


\section{Całka Perrona}
Całka Perrona zastępuje sumy całkowe porównywaniem funkcji, których pochodne Diniego leżą po obu stronach funkcji całkowanej. Daje inną charakterystykę tej samej skalarnej klasy co całka HK.
\begin{definicja}[Pochodne Diniego i funkcje porównawcze]\label{it:def:perron}
Zakładamy $a<b$; na przedziale jednopunktowym przyjmujemy całkę równą zeru.
Dla ciągłej $H:[a,b]\to\R$ definiujemy
\[
\Dn H(x)=\liminf_{\substack{y\to x\\y\ne x}}\frac{H(y)-H(x)}{y-x},\qquad
\Du H(x)=\limsup_{\substack{y\to x\\y\ne x}}\frac{H(y)-H(x)}{y-x}.
\]
Na końcach przedziału granice są jednostronne; dopuszczamy wartości nieskończone. Dla funkcji $f:[a,b]\to\R$ o skończonych wartościach \emph{funkcją górną} nazywamy ciągłą funkcję $M$ z $M(a)=0$ i $\Dn M(x)\ge f(x)$ dla wszystkich $x$; \emph{funkcją dolną} nazywamy ciągłą $m$ z $m(a)=0$ i $\Du m(x)\le f(x)$ dla wszystkich $x$. Całki górną i dolną Perrona definiujemy przez
\[
\overline{(\mathrm{Pe})\!\int_a^b}f=\inf_M M(b),\qquad
\underline{(\mathrm{Pe})\!\int_a^b}f=\sup_m m(b).
\]
Jeśli obie rodziny funkcji porównawczych są niepuste, a te kresy są równe i skończone, nazywamy $f$ całkowalną w sensie Perrona i oznaczamy tę wspólną wartość przez $(\mathrm{Pe})\int_a^b f$.
\end{definicja}
\begin{lemat}[Porządek funkcji porównawczych]\label{it:lem:perron-order}
Jeżeli $H:[a,b]\to\R$ jest ciągła i $\Dn H(x)\ge0$ dla każdego $x$, to $H$ jest niemalejąca. Dla każdej funkcji górnej $M$ i dolnej $m$ funkcja $M-m$ jest niemalejąca, a $m(b)\le M(b)$.
\end{lemat}
\begin{proof}
Dla $\eta>0$ funkcja $G(x)=H(x)+\eta x$ ma $\Dn G\ge\eta$. Gdyby $G(v)<G(u)$ przy $u<v$, jej maksimum na $[u,v]$ byłoby osiągnięte w pewnym $c<v$. W dostatecznie bliskich punktach po prawej od $c$ ilorazy różnicowe są dodatnie, co przeczy maksymalności. Tak więc $G$ jest niemalejąca; przechodząc z $\eta\downarrow0$, otrzymujemy tezę dla $H$. Z definicji ograniczeń Diniego wynika $\Dn(M-m)\ge0$, gdyż dla dostatecznie bliskich $y$ ilorazy różnicowe $M$ i $m$ są odpowiednio $\ge f(x)-\eta$ i $\le f(x)+\eta$. Zastosowanie pierwszej części i $M(a)=m(a)=0$ kończy dowód.
\end{proof}
\begin{lemat}[Mała ciągła korekta na zbiorze wyjątkowym]\label{it:lem:perron-null-correction}
Niech $F\in\AC G_*[a,b]$, $N\subset[a,b]$ będzie zbiorem
miary Lebesgue'a zero, a $f:[a,b]\to\R$ ma skończone wartości.
Dla każdego $\varepsilon>0$ istnieje ciągła niemalejąca funkcja
$H$ z $H(a)=0$, $H(b)<\varepsilon$ taka, że dla każdego $t\in N$
i wszystkich dostatecznie bliskich $y\in[a,b]$ zachodzi
\[
 |H(y)-H(t)|\ge |F(y)-F(t)|+|f(t)|\,|y-t|.
\]
\end{lemat}
\begin{proof}
Niech $[a,b]=\bigcup_{n\ge1}E_n$ będzie domkniętym pokryciem
z definicji $\AC G_*$; puste zbiory pomijamy. Najpierw dla
ustalonego $E=E_n$ konstruujemy skończoną dodatnią miarę
borelowską $\nu$ bez atomów, taką że $\nu(N\cap E)=0$ oraz
\[
 |F(y)-F(t)|\le\nu([\min(t,y),\max(t,y)])
 \qquad(t\in E,\ y\in[a,b]).
\]
Z lematu \ref{it:lem:acstar-affine-extension} wybieramy
bezwzględnie ciągłe przedłużenie $G$ funkcji $F|_E$, afiniczne
w lukach wewnętrznych i stałe w lukach brzegowych.
Miara $|G'|\lambda_1$ kontroluje moduł każdego przyrostu $G$,
na mocy twierdzenia \ref{it:thm:ac-lebesgue-primitive}.
Połóżmy $R=|F-G|$. Funkcja $R$ jest ciągła i znika na $E$.

Dla domknięcia $[l,r]$ luki wewnętrznej definiujemy na nim
\[
 L(x)=\max_{l\le z\le x}R(z),\qquad
 B(x)=\max_{x\le z\le r}R(z).
\]
Funkcje $L$ i $-B$ są niemalejące i ciągłe. Ciągłość wynika
z jednostajnej ciągłości $R$: przy przesunięciu końca
przedziału o $h$ zmiana maksimum nie przekracza $\omega_R(|h|)$.
Ponieważ $R(l)=R(r)=0$, funkcja
$K(x)=L(x)+B(l)-B(x)$ jest ciągła, niemalejąca, $K(l)=0$,
$K(r)=2\max_{[l,r]}R$ oraz
\[
 K(y)-K(l)\ge R(y),\qquad K(r)-K(y)\ge R(y).
\]
Przedłużamy $K$ przez stałe poza $[l,r]$ i oznaczamy przez
$\sigma_{[l,r]}$ jej miarę Lebesgue'a--Stieltjesa (bez atomu
w $a$). Jest bezatomowa i skupiona w otwartej luce.
W luce brzegowej z prawym końcem $r\in E$ stosujemy tylko
$K(x)=B(l)-B(x)$, gdzie $l=a$; w luce z lewym końcem $l\in E$
stosujemy tylko $K=L$, gdzie $r=b$. Zachowujemy odpowiednie
oszacowanie przyrostu od końca należącego do $E$.
Przedłużenia również są ciągłe i ich miary nie obciążają $E$.

Na każdej luce $\max R\le\osc_{[l,r]}F$: wartości $G$ leżą
między wartościami $F$ na jej końcach (w luce brzegowej
$G$ jest stała i równa wartości $F$ na końcu w $E$).
Z dowodu lematu \ref{it:lem:acstar-null-variation} suma
oscylacji na domknięciach wszystkich luk jest skończona.
Zatem miara
\[
 \nu=|G'|\lambda_1+\sum_J\sigma_J
\]
jest skończona, bezatomowa i $\nu(N\cap E)=0$; składnik
bezwzględnie ciągły nie obciąża zbiorów zerowych, a pozostałe
składniki nie obciążają $E$. Dla $t\in E$ mamy
$|F(y)-F(t)|\le |G(y)-G(t)|+R(y)$. Gdy $y\in E$,
drugi składnik znika. Gdy $y$ leży w luce, jej koniec
należący do $E$ leży pomiędzy $t$ i $y$; odpowiednia
nierówność dla $K$ kontroluje $R(y)$ przez miarę
$\sigma_J$ od tego końca do $y$. Otrzymujemy żądane
oszacowanie przez $\nu$.

Wracamy do pokrycia $E_n$ i odpowiadających mu miar $\nu_n$.
Z regularności skończonej miary borelowskiej na zwartym
przedziale, po jej uzupełnieniu, wybieramy względnie
otwarte $O_n\supset N\cap E_n$ takie, że
$\nu_n(O_n)<\varepsilon2^{-n-2}$. Nie wymaga to
borelowskości $N$: część $N\cap E_n$ leży w borelowskim
zbiorze zerowym dla $\nu_n$, otrzymanym przez przecięcie
borelowskiego zbioru $\lambda_1$-zerowego zawierającego $N$
ze zbiorem $E_n$.
Definiujemy $H_n(x)=\nu_n(O_n\cap[a,x])$.
Jest to funkcja ciągła i niemalejąca, gdyż miara jest
bezatomowa, a $H_n(a)=0$. Dla $t\in N\cap E_n$ i $y$
dostatecznie bliskiego $t$ cały odcinek między nimi leży
w $O_n$. Wtedy
$|H_n(y)-H_n(t)|\ge |F(y)-F(t)|$.

Pozostaje składnik zależny od $f$. Dla $k\ge1$ niech
$N_k=\{t\in N:k-1\le |f(t)|<k\}$. Wybieramy względnie
otwarte $U_k\supset N_k$ z
$k\lambda_1(U_k)<\varepsilon2^{-k-2}$ i kładziemy
\[
 H_0(x)=\sum_{k\ge1} k\lambda_1(U_k\cap[a,x]).
\]
Szereg jest jednostajnie zbieżny przez oszacowanie
każdego składnika; jego suma jest ciągła, niemalejąca
i zeruje się w $a$. W punkcie $t\in N_k$ odpowiedni
składnik ma lokalnie nachylenie $k$, więc
$|H_0(y)-H_0(t)|\ge |f(t)|\,|y-t|$ dla bliskich $y$.

Wreszcie $H=H_0+\sum_{n\ge1}H_n$. Także ten szereg jest
jednostajnie zbieżny, a $H(b)\le\varepsilon/2<\varepsilon$.
Wszystkie składniki są niemalejące, więc dla dowolnych
$t,y$ ich przyrosty mają ten sam znak. Wybierając $n$
z $t\in E_n$ i $k$ z $t\in N_k$, dodajemy oba lokalne
oszacowania, co dowodzi tezy.
\end{proof}

\begin{twierdzenie}[Perron i HK]\label{it:thm:perron-hk}
Funkcje całkowalne w sensie Perrona i HK tworzą tę samą klasę, a wartości całek są równe.
\end{twierdzenie}
\begin{proof}
Dowodzimy implikacji Perron $\Rightarrow$ HK. Niech $I$ będzie całką Perrona. Dla $\varepsilon>0$ wybieramy funkcje porównawcze $M,m$ z $M(b)-m(b)<\varepsilon/2$ i kładziemy $\eta=\varepsilon/(2(b-a))$. Z definicji granic Diniego dla każdego $t$ istnieje $\delta(t)>0$ taka, że dla $0<|y-t|<\delta(t)$ zachodzą nierówności
\[
\frac{M(y)-M(t)}{y-t}\ge f(t)-\eta,\qquad
\frac{m(y)-m(t)}{y-t}\le f(t)+\eta.
\]
Dla przedziału $[u,v]$ z $u\le t\le v$ zawartego w otoczeniu cechowania mnożymy pierwszą nierówność przez $v-t>0$ dla $y=v>t$, a dla $y=u<t$ przez $u-t<0$, odwracając jej znak. Dostajemy $M(v)-M(t)\ge(f(t)-\eta)(v-t)$ i $M(t)-M(u)\ge(f(t)-\eta)(t-u)$. Gdy $u=t$ lub $v=t$, odpowiednia nierówność jest równością $0=0$ i nie dzielimy przez zero. Dodanie daje $M(v)-M(u)\ge(f(t)-\eta)(v-u)$. Z drugiej nierówności analogicznie wynika $m(v)-m(u)\le(f(t)+\eta)(v-u)$. Zatem
\[
m(v)-m(u)-\eta(v-u)\le f(t)(v-u)\le M(v)-M(u)+\eta(v-u).
\]
Sumujemy po komórkach dowolnego $\delta$-drobnego podziału, uzyskując $m(b)-\eta(b-a)\le S(f,\mathcal P)\le M(b)+\eta(b-a)$. Ponieważ $m(b)\le I\le M(b)$ i $M(b)-m(b)<\varepsilon/2$, każda suma różni się od $I$ o mniej niż $\varepsilon$. Jest to całkowalność HK.

Dla implikacji HK $\Rightarrow$ Perron niech $f\in\HK[a,b]$
i $F(x)=\int_a^x f$. Z twierdzeń
\ref{it:thm:kh-primitive-derivative} i
\ref{it:thm:kh-acgstar-characterization} wiemy, że
$F\in\AC G_*[a,b]$ i $F'=f$ poza zbiorem zerowym $N$;
dołączamy do $N$ końce $a,b$. Dla $\varepsilon>0$
wybieramy korektę $H$ z lematu \ref{it:lem:perron-null-correction}
i definiujemy $M=F+H$, $m=F-H$. Obie funkcje są ciągłe
i zerują się w $a$. Poza $N$ ilorazy różnicowe $H$ są
nieujemne, dlatego $\Dn M\ge F'=f$ i $\Du m\le F'=f$.
Dla $t\in N$ oraz dostatecznie bliskiego $y\ne t$,
monotoniczność $H$ daje
\[
 \frac{H(y)-H(t)}{y-t}
 =\frac{|H(y)-H(t)|}{|y-t|}
 \ge \left|\frac{F(y)-F(t)}{y-t}\right|+|f(t)|.
\]
Zatem ilorazy różnicowe $M$ są nie mniejsze od $f(t)$,
a ilorazy $m$ nie większe od $f(t)$, także przy
zbliżaniu jednostronnym na końcach. Są to funkcje
porównawcze w sensie definicji. Ich wartości końcowe
wynoszą $F(b)\pm H(b)$. Lemat o porządku oraz
$H(b)<\varepsilon$ ściskają kresy Perrona między
$F(b)-\varepsilon$ i $F(b)+\varepsilon$. Dowolność
$\varepsilon$ daje całkowalność Perrona i wartość $F(b)$.
Dla $a=b$ obie całki przyjmujemy równe zeru.
\end{proof}
\begin{stwierdzenie}[Odtwarzanie pochodnej]
Jeżeli $F:[a,b]\to\R$ jest ciągła i różniczkowalna na $(a,b)$, a $f:[a,b]\to\R$ jest równa $F'$ we wnętrzu i ma dowolne skończone wartości na końcach, to $f$ jest całkowalna w sensie Perrona i $(\mathrm{Pe})\int_a^b f=F(b)-F(a)$.
\end{stwierdzenie}
\begin{proof}
Z twierdzenia o całkowalności pochodnych w rozdziale o całce HK wynika $f\in\HK[a,b]$ oraz $(\HK)\int_a^b f=F(b)-F(a)$, również bez istnienia skończonych pochodnych jednostronnych. Twierdzenie \ref{it:thm:perron-hk} daje tezę. Gdy pochodne jednostronne istnieją i są skończone oraz $f$ przyjmuje ich wartości na końcach, można także bezpośrednio zastosować funkcję $F-F(a)$ jako funkcję górną i dolną; lemat \ref{it:lem:perron-order} ściska obie całki do jej przyrostu.
\end{proof}
\begin{przyklad}[Funkcja charakterystyczna liczb wymiernych]
Niech $f=\ind{\Q\cap[0,1]}$. Dla $\varepsilon>0$ wybieramy otwarty zbiór $O\supset\Q\cap[0,1]$ o mierze Lebesgue'a mniejszej od $\varepsilon$ i definiujemy $M(x)=\lambda(O\cap[0,x])$, $m(x)=0$. Funkcja $M$ jest ciągła, niemalejąca i ma pochodną równą $1$ w każdym punkcie wymiernym, ponieważ punkt ten należy do wnętrza $O$. W pozostałych punktach $\Dn M\ge0=f$, a $\Du m=0\le f$ wszędzie. Rodziny funkcji porównawczych są niepuste. Z lematu \ref{it:lem:perron-order} wynika, że kres dolny wartości funkcji górnych nie jest mniejszy od kresu górnego wartości funkcji dolnych. Ponieważ $m(1)=0$ i $M(1)<\varepsilon$, oba kresy należą do $[0,\varepsilon)$. Dowolność $\varepsilon>0$ daje ich równość i skończoność, zatem całka Perrona istnieje i wynosi zero.
\end{przyklad}
\begin{uwaga}[Dlaczego warunek punktowy jest ważny]
Nie wolno bezkrytycznie zastępować warunków na pochodne funkcji porównawczych nierównościami prawie wszędzie. Funkcja Cantora i funkcja zerowa mają pochodną równą zero prawie wszędzie, ale różne przyrosty na $[0,1]$. Potrzebna jest kontrola zbiorów wyjątkowych.
\end{uwaga}

% ============================================================
% UZUPELNIENIE MATEMATYCZNE: ROZDZIAL 5
% ============================================================
\subsection*{Konstrukcja funkcji porównawczych}
\begin{lemat}[Ścisłe funkcje górne i dolne]\label{lem:perron-scisle}
Niech $f:[a,b]\to\R$ będzie funkcją Perrona-całkowalną, a $I=(\mathrm{Pe})\int_a^b f$. Dla każdego $\varepsilon>0$ istnieją funkcja górna $M$ i funkcja dolna $m$ takie, że
$0\le M(b)-I<\varepsilon$, $0\le I-m(b)<\varepsilon$ i $0\le M(x)-m(x)$ dla wszystkich $x\in[a,b]$.
\end{lemat}
\begin{proof}
Kresy definiujące całki górną i dolną Perrona są równe $I$, zatem z ich definicji można wybrać $M,m$ o wymaganych wartościach końcowych. Funkcja $M-m$ jest niemalejąca na mocy lematu o porządku z początku rozdziału i znika w $a$, więc pozostaje nieujemna na całym przedziale.
\end{proof}
\begin{twierdzenie}[Bezpośrednia konstrukcja dla funkcji z kontrolowanym zbiorem wyjątkowym]\label{tw:perron-wyjatki}
Niech $E\subset(a,b)$ będzie zwartym zbiorem miary Lebesgue'a zero, a $F\in C[a,b]$ będzie różniczkowalna w każdym punkcie $[a,b]\setminus E$ (jednostronnie na końcach). Załóżmy, że $F$ jest lipschitzowska ze stałą $L$ na pewnym otwartym otoczeniu $U$ zbioru $E$, a $f:[a,b]\to\R$ jest funkcją o skończonych wartościach, równą $F'$ poza $E$ oraz ograniczoną przez $K$ na $E$. Wtedy $f$ jest Perrona-całkowalna i
$(\mathrm{Pe})\int_a^b f=F(b)-F(a)$.
\end{twierdzenie}
\begin{proof}
Ustalmy $\varepsilon>0$ i $B>L+K+1$. Wybieramy zbiór otwarty $O$ taki, że $E\subset O$, $\overline O\subset U$ oraz $B\lambda(O)<\varepsilon$; jest to możliwe dzięki zwartości $E$ i $\lambda(E)=0$. Definiujemy $H(x)=B\lambda(O\cap[a,x])$. Funkcja $H$ jest ciągła, niemalejąca, $H(a)=0$ i $H(b)<\varepsilon$. W każdym punkcie $x\in E$ funkcja $H$ jest lokalnie afiniczna o nachyleniu $B$, a nachylenia siecznych $F$ mają moduły nie większe niż $L$. Stąd $\Dn(F+H)(x)\ge B-L>f(x)$ oraz $\Du(F-H)(x)\le L-B<f(x)$.

W punktach $x\notin E$ funkcja $F$ jest różniczkowalna; monotoniczność $H$ daje $\Dn(F+H)(x)\ge F'(x)=f(x)$ i $\Du(F-H)(x)\le F'(x)=f(x)$. Funkcje $M=F-F(a)+H$ i $m=F-F(a)-H$ są więc odpowiednio górna i dolna. Ich wartości końcowe różnią się o $2H(b)<2\varepsilon$. Z dowolności $\varepsilon$ i z definicji całki Perrona wynika teza.
\end{proof}
\begin{uwaga}[Granica argumentu]
W poprzednim dowodzie kontrola lipschitzowska na otoczeniu $E$ była potrzebna, aby wybrać \emph{jedno} skończone nachylenie $B$ dominujące wszystkie sieczne w punktach wyjątkowych. Dla dowolnej pierwotnej HK taka kontrola nie musi istnieć. Ogólna konstrukcja z lematu \ref{it:lem:perron-null-correction} zastępuje jednolite ograniczenie lipschitzowskie miarami kontrolującymi przyrosty na zbiorach pokrycia $\AC G_*$; sama tożsamość $F'=f$ prawie wszędzie nie wystarcza.
\end{uwaga}
\begin{wniosek}[Funkcje kawałkami gładkie]\label{wn:perron-kawalkami}
Jeśli $F\in C[a,b]$ i każde ograniczenie $F|_{[x_{j-1},x_j]}$ jest klasy $C^1$ względem skończonego podziału $a=x_0<\dots<x_m=b$, to dowolna funkcja $f$ równa $F'$ poza punktami podziału i mająca w tych punktach skończone wartości jest Perrona-całkowalna, a jej całka wynosi $F(b)-F(a)$.
\end{wniosek}
\begin{proof}
Niech $L=\max_{1\le j\le m}\sup_{[x_{j-1},x_j]}|F'_j|$, gdzie $F'_j$ jest pochodną odpowiedniego ograniczenia (na końcach jednostronną). Dla $u<v$ wstawiamy między nie wszystkie punkty podziału. Twierdzenie Lagrange'a na każdej części oraz nierówność trójkąta dają $|F(v)-F(u)|\le L(v-u)$, więc $F$ jest globalnie lipschitzowska. Na skończonym zbiorze $E=\{x_1,\ldots,x_{m-1}\}$ funkcja $f$ jest ograniczona. Najpierw wybieramy wartości końcowe $f(a),f(b)$ równe pochodnym jednostronnym; twierdzenie \ref{tw:perron-wyjatki} daje tezę (także przy pustym $E$). Dowolna zmiana skończonych wartości $f$ w punktach końcowych nie zmienia całki, na mocy twierdzenia \ref{it:thm:perron-hk} i lematu \ref{it:lem:kh-null-modification}.
\end{proof}
\begin{przyklad}[Wartość bezwzględna]
Dla $F(x)=|x|$ na $[-1,2]$ niech $f(x)=-1$ dla $x<0$, $f(x)=1$ dla $x>0$, a $f(0)=17$. Wówczas $f$ jest Perrona-całkowalna oraz $(\mathrm{Pe})\int_{-1}^{2}f=F(2)-F(-1)=1$. Dowód nie wymaga nadawania sztucznej pochodnej $F'(0)$.
\end{przyklad}
\subsection*{Porównanie konstrukcji Perrona i HK}
\begin{stwierdzenie}[Kontrola sum z funkcji porównawczych]\label{st:perron-sumy}
Niech $f:[a,b]\to\R$ ma skończone wartości, $M,m$ będą ciągłymi funkcjami porównawczymi z definicji Perrona, a $\eta>0$. Istnieje cechowanie $\delta$ takie, że każdy $\delta$-drobny podział znaczony $\mathcal P$ spełnia
\[
 m(b)-\eta(b-a)\le S(f,\mathcal P)\le M(b)+\eta(b-a).
\]
\end{stwierdzenie}
\begin{proof}
W punkcie $t$ wybieramy promień cechowania tak, aby ilorazy różnicowe funkcji $M$ były większe od $f(t)-\eta$, a ilorazy różnicowe funkcji $m$ mniejsze od $f(t)+\eta$, niezależnie od kierunku zbliżania się do $t$. Dla każdego $u\le t\le v$ z tej nierówności, po uwzględnieniu znaków mianowników dla $u<t$ i $v>t$, otrzymujemy
$m(v)-m(u)-\eta(v-u)\le f(t)(v-u)\le M(v)-M(u)+\eta(v-u)$. Sumowanie po komórkach kończy dowód.
\end{proof}
W ten sposób implikacja Perron $\Rightarrow$ HK jest konstruktywna: para funkcji $M,m$ bliska w punkcie $b$ daje cechowanie kontrolujące \emph{każdą} sumę całkową. Kierunek odwrotny udowodniono w twierdzeniu \ref{it:thm:perron-hk}, wykorzystując charakteryzację pierwotnych HK i małą ciągłą korektę na zbiorze wyjątkowym.


\section{Całka Kurzweila--Stieltjesa}
\label{sec:kurzweil-stieltjes}
Całka Kurzweila--Stieltjesa (KS) łączy przyrosty funkcji całkującej z cechowaniami znanymi z całki Henstocka--Kurzweila. Jej przewaga nad całką Riemanna--Stieltjesa ujawnia się przede wszystkim przy wspólnych nieciągłościach funkcji podcałkowej i całkującej. W tej sekcji, oprócz kryteriów istnienia, wyprowadzimy wzory obliczeniowe i oszacowania błędu przybliżeń.

\subsection*{Definicja, jednoznaczność i podstawowe oszacowania}
\begin{definicja}[Całka Kurzweila--Stieltjesa]
Niech $a<b$ i niech $f,\alpha:[a,b]\to\R$ będą funkcjami o skończonych wartościach. \emph{Cechowaniem} nazywamy funkcję $\delta:[a,b]\to(0,\infty)$. Podział znaczony $\mathcal P=\{(\xi_i,[x_{i-1},x_i])\}_{i=1}^n$, gdzie $a=x_0<\cdots<x_n=b$ i $\xi_i\in[x_{i-1},x_i]$, jest \emph{$\delta$-drobny}, jeżeli $[x_{i-1},x_i]\subset(\xi_i-\delta(\xi_i),\xi_i+\delta(\xi_i))$ dla każdego $i$. Definiujemy
\[
 S_\alpha(f,\mathcal P):=\sum_{i=1}^n f(\xi_i)\bigl(\alpha(x_i)-\alpha(x_{i-1})\bigr).
\]
Liczba $I\in\R$ jest całką KS z $f$ względem $\alpha$, jeżeli dla każdego $\varepsilon>0$ istnieje cechowanie $\delta$ takie, że $|S_\alpha(f,\mathcal P)-I|<\varepsilon$ dla każdego $\delta$-drobnego podziału znaczonego $\mathcal P$. Piszemy wówczas $I=(\KS)\int_a^b f\,\mathrm d\alpha$.
\end{definicja}
Istnienie podziału $\delta$-drobnego dla każdego cechowania wynika z lematu Cousina (lemat 4.2). Gdy $\alpha(x)=x$, definicja jest dokładnie definicją całki Henstocka--Kurzweila. W szczególności nie wolno utożsamiać dowolnej całki KS z całką względem miary Stieltjesa: potrzebne są dodatkowe założenia.

\begin{twierdzenie}[Kryterium Cauchy'ego, liniowość i zgodność]
\label{tw:ks-cauchy}
Dla funkcji $f,\alpha:[a,b]\to\R$ o skończonych wartościach całka $(\KS)\int_a^b f\,\mathrm d\alpha$ istnieje wtedy i tylko wtedy, gdy dla każdego $\varepsilon>0$ istnieje cechowanie $\delta$ takie, że
$|S_\alpha(f,\mathcal P)-S_\alpha(f,\mathcal Q)|<\varepsilon$ dla dowolnych dwóch $\delta$-drobnych podziałów znaczonych $\mathcal P,\mathcal Q$. Całka jest jednoznaczna i liniowa osobno w $f$ i $\alpha$, o ile rozważane całki istnieją. Jeżeli istnieje właściwa całka Riemanna--Stieltjesa $\int_a^b f\,\mathrm d\alpha$, to istnieje również całka KS i ma tę samą wartość.
\end{twierdzenie}
\begin{proof}
Jeżeli sumy są bliskie liczbie $I$, ich różnica jest mniejsza od $2\varepsilon$. Odwrotnie, dla każdego $m\ge1$ wybieramy cechowanie $\delta_m$ zapewniające różnice sum mniejsze niż $2^{-m}$ i zastępujemy je przez $\min(\delta_1,\ldots,\delta_m)$. Wybieramy odpowiadające im podziały $\mathcal P_m$. Gdy $k\ge m$, obie sumy $S_\alpha(f,\mathcal P_k)$ i $S_\alpha(f,\mathcal P_m)$ są porównywalne przez cechowanie $\delta_m$, toteż tworzą ciąg Cauchy'ego. Niech $I$ będzie jego granicą. Dla dowolnego $\delta_m$-drobnego podziału $\mathcal P$ porównujemy jego sumę z $\mathcal P_k$ i przechodzimy z $k$ do nieskończoności; otrzymujemy $|S_\alpha(f,\mathcal P)-I|\le2^{-m}$. Jednoznaczność wynika z istnienia podziałów drobnych dla minimum dwóch cechowań. Liniowość dostajemy z liniowości sum i ponownego użycia minimum cechowań. Dla całki Riemanna--Stieltjesa wystarcza cechowanie stałe mniejsze niż dopuszczalna średnica podziału.
\end{proof}

\begin{definicja}[Wahanie całkowite i funkcja regulowana]
Dla $\alpha:[a,b]\to\R$ definiujemy $V_a^b(\alpha):=\sup_P\sum_i|\alpha(x_i)-\alpha(x_{i-1})|$, gdzie supremum przebiega wszystkie podziały $P$ przedziału $[a,b]$. Piszemy $\alpha\in\BV[a,b]$, jeśli $V_a^b(\alpha)<\infty$. Funkcja $f:[a,b]\to\R$ jest \emph{regulowana}, gdy istnieją skończone granice $f(t-)$ dla $t\in(a,b]$ i $f(t+)$ dla $t\in[a,b)$. Równoważnie $f$ jest jednostajną granicą funkcji schodkowych, przy czym wartości w skończonej liczbie węzłów można określać niezależnie od wartości na otwartych przedziałach.
\end{definicja}
\begin{proof}[Uzasadnienie równoważności]
Każda funkcja schodkowa ma granice jednostronne, a jednostajna granica funkcji posiadających takie granice spełnia jednostronny warunek Cauchy'ego. W przeciwnym kierunku ustalmy $\varepsilon>0$. W każdym punkcie $t$ mamy jednostronne otoczenia, na których oscylacja $f$ jest mniejsza niż $\varepsilon$. Lemat Cousina dostarcza skończony podział, którego wnętrza komórek można rozciąć w znacznikach, tak aby na każdym z otrzymanych otwartych przedziałów oscylacja $f$ była mniejsza niż $\varepsilon$. Na tych przedziałach podstawiamy stałe, a w węzłach zachowujemy dokładne wartości $f$. Powstała funkcja schodkowa przybliża $f$ jednostajnie z błędem co najwyżej $\varepsilon$.
\end{proof}

\begin{lemat}[Oszacowanie przez wahanie]
\label{lem:ks-var}
Jeżeli $\alpha\in\BV[a,b]$ i $f:[a,b]\to\R$ jest ograniczona, to dla każdego podziału znaczonego $\mathcal P$ zachodzi $|S_\alpha(f,\mathcal P)|\le\|f\|_\infty V_a^b(\alpha)$. Jeżeli całka KS istnieje, to
\[
 \left|(\KS)\int_a^b f\,\mathrm d\alpha\right|\le\|f\|_\infty V_a^b(\alpha).
\]
\end{lemat}
\begin{proof}
Stosujemy nierówność trójkąta do sumy $S_\alpha(f,\mathcal P)$, a następnie przechodzimy do granicy po podziałach $\delta$-drobnych.
\end{proof}

\subsection*{Skoki, wartości w punktach i funkcje schodkowe}
Dla regulowanej $g:[a,b]\to\R$ stosujemy oznaczenia $\Delta^-g(t)=g(t)-g(t-)$ dla $t\in(a,b]$, $\Delta^+g(t)=g(t+)-g(t)$ dla $t\in[a,b)$ i $\Delta g(t)=g(t+)-g(t-)$ dla $t\in(a,b)$. Na końcach przyjmujemy $\Delta^-g(a)=\Delta^+g(b)=0$ oraz $\Delta g(a)=\Delta^+g(a)$, $\Delta g(b)=\Delta^-g(b)$. Rozróżnienie skoku lewostronnego i prawostronnego jest istotne: $g(t)$ nie musi równać się żadnej granicy jednostronnej.

\begin{lemat}[Punktowy wzór wykonawczy]
\label{lem:ks-point}
Niech $g:[a,b]\to\R$ będzie regulowana i $c\in[a,b]$. Dla dowolnej liczby $q\in\R$ zachodzi
\[
 (\KS)\int_a^b q\ind{\{c\}}(t)\,\mathrm dg(t)=q\Delta g(c).
\]
W szczególności wartość funkcji podcałkowej w punkcie skoku nie może być pomijana.
\end{lemat}
\begin{proof}
Rozpatrzmy $a<c<b$. Dla $t\ne c$ wybieramy $\delta(t)<|t-c|/2$; wartość $\delta(c)>0$ dobieramy tak małą, aby $|g(u)-g(c-)|<\varepsilon$ dla $c-\delta(c)<u<c$ i $|g(v)-g(c+)|<\varepsilon$ dla $c<v<c+\delta(c)$. W podziale drobnym każda komórka, której wnętrze zawiera $c$, musi mieć znacznik $c$. Jeżeli $c$ jest węzłem podziału, jedynie dwie przyległe do niego komórki mogą mieć znacznik $c$; obie rzeczywiście muszą go mieć przy dodatniej długości. Niezerowe składniki sumy teleskopują więc do $q(g(v)-g(u))$, gdzie $u\le c\le v$, a odległości odpowiednich końców od $c$ są mniejsze niż $\delta(c)$. Wartość sumy różni się od $q(g(c+)-g(c-))$ co najwyżej o $2|q|\varepsilon$. Gdy $c=a$ lub $c=b$, pozostaje tylko jedna z dwóch stron, co daje odpowiednio $q\Delta^+g(a)$ i $q\Delta^-g(b)$. W razie $q=0$ twierdzenie jest natychmiastowe.
\end{proof}

\begin{lemat}[Wskaźniki przedziałów]
\label{lem:ks-indicators}
Jeśli $g:[a,b]\to\R$ jest regulowana, to dla $a\le u<v\le b$ istnieją całki KS z funkcji charakterystycznych $[u,v]$, $(u,v)$, $[u,v)$ i $(u,v]$. Przykładowo dla $a<u<v<b$ otrzymujemy
\[
 \begin{aligned}
 (\KS)\int_a^b\ind{(u,v)}\,\mathrm dg&=g(v-)-g(u+),\\
 (\KS)\int_a^b\ind{[u,v]}\,\mathrm dg&=g(v+)-g(u-),\\
 (\KS)\int_a^b\ind{[u,v)}\,\mathrm dg&=g(v-)-g(u-),\\
 (\KS)\int_a^b\ind{(u,v]}\,\mathrm dg&=g(v+)-g(u+).
 \end{aligned}
\]
Gdy $u=a$ lub $v=b$, symbole $g(a-)$ i $g(b+)$ we wzorach zastępujemy odpowiednio $g(a)$ i $g(b)$.
\end{lemat}
\begin{proof}
Dla funkcji $\ind{(u,v)}$ wybieramy cechowanie tak, aby przedziały o znacznikach różnych od $u,v$ nie przecinały tych punktów. Komórki mające znaczniki $u$ lub $v$ wnoszą zero, a wszystkie komórki ze znacznikami wewnątrz $(u,v)$ tworzą jeden blok. Ich przyrosty teleskopują do $g(r)-g(s)$, gdzie $s\to u+$ i $r\to v-$, skąd pierwszy wzór. Następnie korzystamy z tożsamości funkcji charakterystycznych, np. $\ind{[u,v]}=\ind{(u,v)}+\ind{\{u\}}+\ind{\{v\}}$, oraz z lematu~\ref{lem:ks-point}. Pozostałe wzory uzyskujemy analogicznie.
\end{proof}

\begin{twierdzenie}[Całkowanie względem funkcji schodkowej]
\label{tw:ks-step}
Niech $g:[a,b]\to\R$ będzie funkcją schodkową, mającą skończony zbiór punktów nieciągłości $D$, i niech $f:[a,b]\to\R$ będzie dowolną funkcją o skończonych wartościach. Wtedy całka KS istnieje i
\[
 (\KS)\int_a^b f\,\mathrm dg
 =f(a)\Delta^+g(a)+\sum_{c\in D\cap(a,b)}f(c)\Delta g(c)+f(b)\Delta^-g(b).
\]
\end{twierdzenie}
\begin{proof}
Zapisujemy $g$ jako sumę stałej oraz skończonej liczby funkcji $\ind{[c,b]}$ i $\ind{(c,b]}$, uwzględniając oddzielnie skoki przy końcach. Z lematu~\ref{lem:ks-point} wynika dla każdej z tych funkcji całkujących tożsamość $\int f\,\mathrm d\ind{[c,b]}=f(c)$ i $\int f\,\mathrm d\ind{(c,b]}=f(c)$. Liniowość względem integratora daje wskazany wzór. Nie potrzebujemy ograniczoności ani regularności $f$, ponieważ występuje skończenie wiele wartości $f(c)$.
\end{proof}

\begin{przyklad}[Wspólny skok i nieistnienie zwykłej całki Stieltjesa]
Niech $a<c<b$ i $H=\ind{[c,b]}$. Wtedy dla każdej funkcji $f$ o skończonych wartościach zachodzi $(\KS)\int_a^b f\,\mathrm dH=f(c)$. Stąd $(\KS)\int_a^b H\,\mathrm dH=1$, a $(\KS)\int_a^b\ind{(c,b]}\,\mathrm dH=0$. Pierwsza z tych całek nie istnieje jako zwykła całka Riemanna--Stieltjesa, ponieważ swobodny wybór znacznika na komórce przecinającej $c$ daje różne sumy. Całka KS rozstrzyga tę niejednoznaczność za pomocą cechowania.
\end{przyklad}

\begin{przyklad}[Punktowy impuls integratora]
Niech $g=\ind{\{c\}}$ dla $a<c<b$. Wtedy $\Delta^-g(c)=1$ i $\Delta^+g(c)=-1$, zatem $\Delta g(c)=0$. Dla każdego regulowanego $f$ otrzymujemy $(\KS)\int_a^b f\,\mathrm dg=0$. Oznacza to, że dwie przeciwne zmiany funkcji całkującej, następujące po obu stronach tego samego punktu, znoszą się. Wartość $g(c)$ może więc nie wpływać na całkę, mimo że oddzielne skoki jednostronne są niezerowe.
\end{przyklad}

\subsection*{Dwa podstawowe twierdzenia o istnieniu}
\begin{twierdzenie}[Funkcja regulowana i całkująca o wahaniu ograniczonym]
\label{tw:ks-reg-bv}
Jeśli $f:[a,b]\to\R$ jest regulowana, a $g\in\BV[a,b]$, to całka $(\KS)\int_a^b f\,\mathrm dg$ istnieje. Ponadto
$|(\KS)\int_a^b f\,\mathrm dg|\le\|f\|_\infty V_a^b(g)$.
\end{twierdzenie}
\begin{proof}
Wybieramy funkcje schodkowe $f_n$ z $\|f_n-f\|_\infty\to0$. Każda całka $I_n=(\KS)\int_a^b f_n\,\mathrm dg$ istnieje na mocy lematu~\ref{lem:ks-indicators}, ponieważ $f_n$ jest skończoną kombinacją wskaźników przedziałów i punktów. Z lematu~\ref{lem:ks-var} i liniowości wynika $|I_n-I_m|\le\|f_n-f_m\|_\infty V_a^b(g)$. Niech $I=\lim_n I_n$. Dla dowolnego podziału znaczonego $\mathcal P$ mamy
\[
 |S_g(f,\mathcal P)-I|\le
 \|f-f_n\|_\infty V_a^b(g)+|S_g(f_n,\mathcal P)-I_n|+|I_n-I|.
\]
Dla ustalonego $\varepsilon$ najpierw wybieramy $n$ zapewniające małość pierwszego i trzeciego składnika, a następnie cechowanie dla $f_n$. Ostatnia nierówność wynika z lematu~\ref{lem:ks-var} po przejściu do granicy.
\end{proof}

\begin{lemat}[Nierówność sumowania przez części]
\label{lem:ks-summation}
Niech $f\in\BV[a,b]$ i niech $g:[a,b]\to\R$ będzie ograniczona. Dla każdego podziału znaczonego $\mathcal P$ zachodzi
\[
 |S_g(f,\mathcal P)|\le
 \bigl(2\|f\|_\infty+V_a^b(f)\bigr)\|g\|_\infty.
\]
\end{lemat}
\begin{proof}
Oznaczmy $f_i=f(\xi_i)$, $g_i=g(x_i)$. Przekształcenie skończonej sumy daje
$S_g(f,\mathcal P)=f_ng_n-f_1g_0+\sum_{i=1}^{n-1}(f_i-f_{i+1})g_i$. Ponieważ znaczniki są uporządkowane, $\xi_i\le\xi_{i+1}$, mamy $\sum_{i=1}^{n-1}|f_i-f_{i+1}|\le V_a^b(f)$. Oszacowanie wynika z nierówności trójkąta.
\end{proof}

\begin{twierdzenie}[Całkująca regulowana i funkcja podcałkowa o wahaniu ograniczonym]
\label{tw:ks-bv-reg}
Jeśli $f\in\BV[a,b]$ i $g:[a,b]\to\R$ jest regulowana, to całka $(\KS)\int_a^b f\,\mathrm dg$ istnieje. Jeżeli $g,h$ są regulowane, to
\[
 \left|(\KS)\int_a^b f\,\mathrm dg-(\KS)\int_a^b f\,\mathrm dh\right|
 \le\bigl(2\|f\|_\infty+V_a^b(f)\bigr)\|g-h\|_\infty.
\]
\end{twierdzenie}
\begin{proof}
Funkcję $g$ przybliżamy jednostajnie funkcjami schodkowymi $g_n$. Całki $I_n=(\KS)\int_a^b f\,\mathrm dg_n$ istnieją na mocy twierdzenia~\ref{tw:ks-step}. Z lematu~\ref{lem:ks-summation}, zastosowanego do $g_n-g_m$, otrzymujemy $|I_n-I_m|\le K_f\|g_n-g_m\|_\infty$, gdzie $K_f=2\|f\|_\infty+V_a^b(f)$. Istnieje zatem $I=\lim_nI_n$. Dla dowolnego podziału $\mathcal P$ mamy
$|S_g(f,\mathcal P)-I|\le K_f\|g-g_n\|_\infty+|S_{g_n}(f,\mathcal P)-I_n|+|I_n-I|$. Wybór najpierw $n$, a następnie cechowania dla $g_n$, dowodzi istnienia całki. To samo oszacowanie dla całki różnicy $g-h$ daje drugą tezę.
\end{proof}

\begin{uwaga}[Właściwa dziedzina obu twierdzeń]
Twierdzenia~\ref{tw:ks-reg-bv} i~\ref{tw:ks-bv-reg} są różne. Pierwsze pozwala na funkcję podcałkową regulowaną o nieskończonym wahaniu, o ile całkująca ma wahanie ograniczone; drugie pozwala na całkującą regulowaną o nieskończonym wahaniu, gdy funkcja podcałkowa ma wahanie ograniczone. Żadne z tych twierdzeń nie wymaga ciągłości obu funkcji ani braku wspólnych skoków.
\end{uwaga}

\begin{twierdzenie}[Konstruktywny rozkład funkcji o wahaniu ograniczonym]
\label{tw:ks-bv-split}
Dla $g\in\BV[a,b]$ zdefiniujmy
\[
 g_{\mathrm{sk}}(t)=
 \sum_{c\in(a,b]}\Delta^-g(c)\ind{[c,b]}(t)
 +\sum_{c\in[a,b)}\Delta^+g(c)\ind{(c,b]}(t).
\]
Oba szeregi są bezwzględnie i jednostajnie zbieżne, a funkcja $g_{\mathrm c}:=g-g(a)-g_{\mathrm{sk}}$ jest ciągła i ma wahanie ograniczone. Dla każdej regulowanej $f:[a,b]\to\R$ zachodzi
\[
 \begin{aligned}
 (\KS)\int_a^b f\,\mathrm dg
 ={}&(\KS)\int_a^b f\,\mathrm dg_{\mathrm c}
 +f(a)\Delta^+g(a)\\
 &+\sum_{c\in(a,b)}f(c)\bigl(\Delta^-g(c)+\Delta^+g(c)\bigr)
 +f(b)\Delta^-g(b).
 \end{aligned}
\]
W szczególności dla funkcji całkującej złożonej wyłącznie ze skoków całka jest szeregiem wartości $f$ w punktach skoków, niezależnie od ich liczby.
\end{twierdzenie}
\begin{proof}
Dla każdego skończonego zbioru punktów $E\subset[a,b]$ suma modułów wszystkich skoków jednostronnych $g$ w tych punktach nie przekracza $V_a^b(g)$. Aby to zobaczyć, wybieramy parami rozłączne, odpowiednio małe otoczenia punktów $E$, a następnie podziały uwzględniające punkty $c$ i po jednym punkcie z każdej dozwolonej strony $c$. W definicji wahania przechodzimy do granic jednostronnych. Biorąc supremum po skończonych $E$, otrzymujemy
$\sum_{c\in(a,b]}|\Delta^-g(c)|+\sum_{c\in[a,b)}|\Delta^+g(c)|\le V_a^b(g)$. Zatem szeregi definiujące $g_{\mathrm{sk}}$ są absolutnie i jednostajnie zbieżne.

Skok lewostronny funkcji $g_{\mathrm{sk}}$ w $c\in(a,b]$ jest równy $\Delta^-g(c)$, a skok prawostronny w $c\in[a,b)$ wynosi $\Delta^+g(c)$. Ponieważ $g$ jest regulowana, różnica $g_{\mathrm c}$ nie ma żadnych skoków jednostronnych i dlatego jest ciągła. Ma również wahanie ograniczone, jako różnica funkcji o wahaniu ograniczonym. Dla skończonych sum wzór na całkę wynika z twierdzenia~\ref{tw:ks-step}. Przy zwiększaniu liczby składników ogon ma wahanie ograniczone przez sumę modułów pominiętych skoków, a więc jego całka ma moduł nie większy niż $\|f\|_\infty$ razy ten ogon, na mocy lematu~\ref{lem:ks-var}. Przejście do granicy dowodzi wzoru.
\end{proof}

\subsection*{Wzory przez miarę Stieltjesa i rozkład Lebesgue'a}
\begin{twierdzenie}[Miara kanonicznie związana z funkcją $\BV$]
\label{tw:ks-measure}
Niech $g\in\BV[a,b]$. Definiujemy $G:[a,b]\to\R$ przez $G(t)=g(t+)-g(a)$ dla $a\le t<b$ oraz $G(b)=g(b)-g(a)$. Istnieje dokładnie jedna skończona \emph{podpisana miara borelowska} $\mu_g$ na $[a,b]$ spełniająca $\mu_g([a,t])=G(t)$ dla każdego $t\in[a,b]$. Jej atomy są dane wzorami
\[
 \mu_g(\{a\})=\Delta^+g(a),\qquad
 \mu_g(\{c\})=\Delta g(c)\ \ (a<c<b),\qquad
 \mu_g(\{b\})=\Delta^-g(b).
\]
Jeśli $f$ jest regulowana, to $f$ jest całkowalna względem $|\mu_g|$ i
\[
 (\KS)\int_a^b f\,\mathrm dg=\int_{[a,b]}f(t)\,\mathrm d\mu_g(t),
 \qquad |\mu_g|([a,b])\le V_a^b(g).
\]
\end{twierdzenie}
\begin{proof}
Funkcja $G$ jest prawostronnie ciągła na $[a,b)$ i ma wahanie ograniczone. Rozkłada się na różnicę dwóch funkcji niemalejących i prawostronnie ciągłych (stosujemy funkcję wahania całkowitego i rozkład Jordana). Każda z tych funkcji wyznacza skończoną miarę Stieltjesa; ich różnica daje $\mu_g$. Jednoznaczność wynika z jednoznaczności miar na półpierścieniu przedziałów $(s,t]$ wraz z jednoelementowym zbiorem $\{a\}$. Wzory na atomy otrzymujemy przez różnice granic jednostronnych $G$. Nierówność dla wahania miary wynika z oszacowania wszystkich sum bezwzględnych przyrostów przez $V_a^b(g)$.

Dla funkcji $\ind{\{c\}}$ równość całek jest lematem~\ref{lem:ks-point}; dla wskaźników czterech typów przedziałów otrzymujemy ją z lematu~\ref{lem:ks-indicators} oraz wzorów na masy przedziałów. Z liniowości wynika równość dla każdej funkcji schodkowej. Wybieramy funkcje schodkowe $f_n\to f$ jednostajnie. Z twierdzenia~\ref{tw:ks-reg-bv} i lematu~\ref{lem:ks-var} całki KS z $f_n$ zbiegają do całki KS z $f$, zaś całki względem $\mu_g$ zbiegają dzięki nierówności $|\int(f_n-f)\,\mathrm d\mu_g|\le\|f_n-f\|_\infty|\mu_g|([a,b])$. Stąd teza.
\end{proof}

\begin{uwaga}[Konwencja na końcach i równość prawie wszędzie]
Jeśli $g$ jest niemalejąca i prawostronnie ciągła na $[a,b]$, to $\mu_g(\{a\})=0$ i $\mu_g((s,t])=g(t)-g(s)$; jest to konwencja używana wcześniej przy całce Lebesgue'a--Stieltjesa. Dla funkcji, która nie jest prawostronnie ciągła, nie wolno stosować tego samego wzoru bez uwzględnienia przesunięcia do $g(t+)$. Równość $f=h$ prawie wszędzie względem miary Lebesgue'a nie wystarcza do równości całek KS względem $g$ z atomami: trzeba uwzględnić wartości w punktach, którym $\mu_g$ przypisuje masę.
\end{uwaga}

\begin{twierdzenie}[Wzór: część absolutnie ciągła, osobliwa i skoki]
\label{tw:ks-decomp}
Niech $g\in\BV[a,b]$ i $f:[a,b]\to\R$ będzie regulowana. Niech rozkład Lebesgue'a podpisanej miary $\mu_g$ względem $\lambda$ ma postać
$\mu_g=q\lambda+\mu_{\mathrm{sc}}+\mu_{\mathrm{at}}$, gdzie $q\in L^1([a,b],\lambda)$, $\mu_{\mathrm{sc}}$ jest bezatomową miarą podpisaną osobliwą względem $\lambda$, a $\mu_{\mathrm{at}}$ jest częścią czysto atomową. Wtedy
\[
 \begin{aligned}
 (\KS)\int_a^b f\,\mathrm dg
 ={}&\int_a^b f(t)q(t)\,\mathrm d\lambda(t)
 +\int_{[a,b]}f\,\mathrm d\mu_{\mathrm{sc}}\\
 &+f(a)\Delta^+g(a)+\sum_{c\in(a,b)}f(c)\Delta g(c)
 +f(b)\Delta^-g(b).
 \end{aligned}
\]
Szereg jest bezwzględnie zbieżny. Jeżeli $g$ nie ma osobliwej części ciągłej, drugi składnik znika; jeśli $g$ jest absolutnie ciągła, znikają również wszystkie składniki skokowe i $q=g'$ prawie wszędzie.
\end{twierdzenie}
\begin{proof}
Z twierdzenia~\ref{tw:ks-measure} całka KS jest całką względem skończonej miary podpisanej $\mu_g$. Stosujemy rozkład Lebesgue'a jej części bezatomowej względem $\lambda$ i liniowość całki. Każda skończona miara podpisana ma co najwyżej przeliczalnie wiele atomów niezerowej masy, a suma modułów ich mas nie przekracza $|\mu_g|([a,b])\le V_a^b(g)$. Stąd $\sum|f(c)\Delta g(c)|\le\|f\|_\infty V_a^b(g)$ i bezwzględna zbieżność. Ostatnia część wynika ze wzoru Newtona--Leibniza dla funkcji absolutnie ciągłej $g(t)=g(a)+\int_a^t g'(s)\,\mathrm d\lambda(s)$.
\end{proof}

\begin{wniosek}[Wzór dla funkcji absolutnie ciągłej z impulsami]
\label{wn:ks-ac-jumps}
Niech $h\in\AC[a,b]$, $a<c_1<\cdots<c_m<b$, $A_j,B_j\in\R$ oraz
$g(t)=h(t)+\sum_{j=1}^m\bigl(A_j\ind{[c_j,b]}(t)+B_j\ind{(c_j,b]}(t)\bigr)$. Dla każdej regulowanej $f$ zachodzi
\[
 (\KS)\int_a^b f\,\mathrm dg
 =\int_a^b f(t)h'(t)\,\mathrm d\lambda(t)
 +\sum_{j=1}^m(A_j+B_j)f(c_j).
\]
Gdy $h\in C^1[a,b]$, pierwsza całka jest zwykłą całką Riemanna z funkcji regulowanej $fh'$.
\end{wniosek}
\begin{proof}
Dla części $h$ stosujemy twierdzenie~\ref{tw:ks-decomp}; dla każdej funkcji skokowej stosujemy twierdzenie~\ref{tw:ks-step}. Liniowość względem funkcji całkującej kończy dowód.
\end{proof}

\begin{przyklad}[Skończona liczba impulsów: pełny rachunek]
Na $[0,1]$ niech $g(t)=t+3\ind{[1/3,1]}(t)-2\ind{(2/3,1]}(t)$ oraz $f(t)=t^2$. Ponieważ funkcja $t\mapsto t$ jest absolutnie ciągła, a skoki $g$ mają wielkości odpowiednio $3$ i $-2$, otrzymujemy
\[
 (\KS)\int_0^1 t^2\,\mathrm dg(t)
 =\int_0^1 t^2\,\mathrm d\lambda(t)
 +3\left(\frac13\right)^2-2\left(\frac23\right)^2
 =\frac13+\frac13-\frac89=-\frac29.
\]
Wartość $g(2/3)$ jest tutaj inna niż $g((2/3)+)$, ale wzór na całkę korzysta z całkowitego skoku $\Delta g(2/3)$.
\end{przyklad}

\begin{przyklad}[Dwa atomy na końcach przedziału]
Na $[0,1]$ niech $g(t)=\ind{(0,1]}(t)+2\ind{\{1\}}(t)$ i $f(t)=1+t$. Funkcja $g$ ma atom o masie $1$ w punkcie $0$ i atom o masie $2$ w punkcie $1$, mimo że wewnątrz przedziału jest stała. Z lematu~\ref{lem:ks-point} otrzymujemy
$(\KS)\int_0^1 f\,\mathrm dg=f(0)\cdot1+f(1)\cdot2=5$.
Ten rachunek pokazuje, dlaczego w twierdzeniach trzeba osobno określić sposób traktowania końców przedziału.
\end{przyklad}

\begin{przyklad}[Nieskończenie wiele impulsów i oszacowanie ogona]
Niech $c_n=1-2^{-n}$, $g(t)=\sum_{n\ge1}2^{-n}\ind{[c_n,1]}(t)$ na $[0,1]$ oraz $f(t)=t$. Wówczas $g$ jest niemalejąca, $V_0^1(g)=1$ i
\[
 (\KS)\int_0^1 t\,\mathrm dg(t)
 =\sum_{n=1}^\infty2^{-n}(1-2^{-n})
 =1-\frac13=\frac23.
\]
Jeśli liczymy tylko pierwszych $N$ impulsów, błąd nie przekracza $\|f\|_\infty\sum_{n>N}2^{-n}=2^{-N}$. Jest to konkretny, certyfikowany algorytm przybliżania całki.
\end{przyklad}

\begin{przyklad}[Miara Cantora: brak wzoru $g'\,\mathrm dx$]
Niech $C:[0,1]\to[0,1]$ będzie funkcją Cantora. Funkcja ta jest ciągła, niemalejąca, $C'=0$ prawie wszędzie i $C(1)-C(0)=1$, więc jej miara Stieltjesa jest osobliwa ciągła. Dla $f(t)=t$ dostajemy $(\KS)\int_0^1 t\,\mathrm dC(t)=1/2$ (z symetrii miary Cantora), a dla $f(t)=t^2$ wartość $3/8$ (obliczenie momentów w rozdziale 10). Zastąpienie $\mathrm dC$ przez $C'(t)\,\mathrm d\lambda(t)$ dałoby błędnie zero. Dlatego w twierdzeniu~\ref{tw:ks-decomp} osobliwa część ciągła jest niezbędna.
\end{przyklad}

\subsection*{Całkowanie przez części ze wszystkimi poprawkami skokowymi}
\begin{twierdzenie}[Wzór całkowania przez części]
\label{tw:ks-ibp}
Niech $f,g:[a,b]\to\R$ będą regulowane i niech przynajmniej jedna z funkcji $f,g$ ma wahanie ograniczone. Wtedy obie całki KS istnieją i
\[
 \begin{aligned}
 (\KS)\int_a^b f\,\mathrm dg+(\KS)\int_a^b g\,\mathrm df
 ={}&f(b)g(b)-f(a)g(a)\\
 &+\sum_{c\in[a,b]}\bigl(\Delta^-f(c)\Delta^-g(c)
 -\Delta^+f(c)\Delta^+g(c)\bigr).
 \end{aligned}
\]
Przyjmujemy $\Delta^-f(a)=\Delta^-g(a)=0$ i $\Delta^+f(b)=\Delta^+g(b)=0$. Szereg korekcyjny jest bezwzględnie zbieżny. W szczególności, jeśli jedna z funkcji jest ciągła, cały wyraz korekcyjny znika.
\end{twierdzenie}
\begin{proof}
\emph{Krok 1: obliczenie przy jednym wspólnym skoku.} Ustalmy $c\in(a,b)$ i zapiszmy $f_-=f(c-)$, $f_0=f(c)$, $f_+=f(c+)$ oraz analogicznie $g_-,g_0,g_+$. Punktowy wkład do sumy dwóch całek wynosi $f_0(g_+-g_-)+g_0(f_+-f_-)$. Tymczasem bezpośrednie wymnożenie daje
\[
 f_+g_+-f_-g_-=f_0(g_+-g_-)+g_0(f_+-f_-)
 +(f_+-f_0)(g_+-g_0)-(f_0-f_-)(g_0-g_-).
\]
Po przeniesieniu dwóch ostatnich składników na drugą stronę otrzymujemy wyraz $\Delta^-f(c)\Delta^-g(c)-\Delta^+f(c)\Delta^+g(c)$. Na końcach $a,b$ występuje tylko jedna z dwóch połówek skoku.

\emph{Krok 2: funkcje schodkowe.} Wybieramy wspólny skończony zbiór węzłów obu funkcji. Na otwartych przedziałach pomiędzy nimi obie funkcje są stałe. Całki liczymy z twierdzenia~\ref{tw:ks-step}, sumując wkłady w węzłach. Powyższa tożsamość algebraiczna po zsumowaniu teleskopuje się do różnicy $f(b)g(b)-f(a)g(a)$ i sumy poprawek.

\emph{Krok 3: przejście do funkcji regulowanych.} Możemy bez straty ogólności założyć $f\in\BV[a,b]$. Funkcję $g$ przybliżamy jednostajnie funkcjami schodkowymi $g_n$. Dla ustalonej funkcji schodkowej $g_n$ otrzymujemy tożsamość najpierw dla schodkowych przybliżeń $f_m$ funkcji $f$ i przechodzimy do granicy przy $m\to\infty$: całka $\int f_m\,\mathrm dg_n$ jest ciągła w normie jednostajnej z oszacowaniem przez $V(g_n)$, a całka $\int g_n\,\mathrm df_m$ jest ciągła względem normy jednostajnej integratora dzięki lematowi~\ref{lem:ks-summation}, zastosowanemu do $g_n$ jako funkcji podcałkowej o wahaniu ograniczonym. Sumy poprawek zawierają tylko skończenie wiele punktów skokowych $g_n$, więc również przechodzą do granicy. Istnienie potrzebnych schodkowych przybliżeń $f_m\to f$ wynika z regulowania $f$.

Przy $n\to\infty$ całki $\int f\,\mathrm dg_n$ zbiegają do $\int f\,\mathrm dg$ z twierdzenia~\ref{tw:ks-bv-reg}; całki $\int g_n\,\mathrm df$ zbiegają dzięki lematowi~\ref{lem:ks-var}. Skoki $g_n$ zbiegają jednostajnie do odpowiednich jednostronnych skoków $g$, a sumę ich iloczynów ze skokami $f$ kontrolujemy przez $V(f)\|g_n-g\|_\infty$: przykładowo $\sum_c|\Delta^-f(c)\Delta^-(g_n-g)(c)|\le2\|g_n-g\|_\infty V(f)$. Druga suma podlega temu samemu oszacowaniu. Wyrazy korekcyjne zbiegają więc bezwzględnie. Dowód zakończony.
\end{proof}

\begin{wniosek}[Klasyczna postać i wzór dla funkcji absolutnie ciągłej]
\label{wn:ks-ibp-ac}
Jeśli $f\in\AC[a,b]$, a $g:[a,b]\to\R$ jest regulowana, to
\[
 (\KS)\int_a^b f\,\mathrm dg
 =f(b)g(b)-f(a)g(a)-\int_a^b g(t)f'(t)\,\mathrm d\lambda(t).
\]
Jeśli $f\in C^1[a,b]$, ostatnią całkę można rozumieć jako zwykłą całkę Riemanna. Nie wymagamy w tym wniosku, aby $g$ miała wahanie ograniczone.
\end{wniosek}
\begin{proof}
Funkcja absolutnie ciągła ma wahanie ograniczone i jest ciągła, więc wzór z twierdzenia~\ref{tw:ks-ibp} nie ma poprawek skokowych. Ponadto $\mathrm df=f'(t)\,\mathrm d\lambda(t)$; z twierdzenia~\ref{tw:ks-decomp}, zastosowanego do całki $\int g\,\mathrm df$ (integrator $f\in\BV$, funkcja podcałkowa $g$ jest regulowana), otrzymujemy $\int g\,\mathrm df=\int gf'\,\mathrm d\lambda$. Iloczyn $gf'$ jest całkowalny w sensie Lebesgue'a, ponieważ $g$ jest ograniczona, a $f'\in L^1$.
\end{proof}

\begin{przyklad}[Rachunek bez znajomości pochodnej funkcji całkującej]
Niech $g:[0,1]\to\R$ będzie dowolną funkcją regulowaną (nie zakładamy $g\in\BV$). Dla $f(t)=t^2$ otrzymujemy
\[
 (\KS)\int_0^1 t^2\,\mathrm dg(t)=g(1)-2\int_0^1 t\,g(t)\,\mathrm d\lambda(t).
\]
Ponieważ funkcja regulowana jest ograniczona i Riemanna całkowalna, ostatnia całka jest zwykłą całką Riemanna. W szczególności wzór działa dla funkcji $g$ oscylującej z nieskończonym wahaniem, o ile ma granice jednostronne.
\end{przyklad}

\begin{przyklad}[Całkująca o nieskończonym wahaniu]
Dla $t\in(0,1]$ niech $g(t)=t\sin(1/t^2)$, a $g(0)=0$. Funkcja $g$ jest ciągła, ale nie ma wahania ograniczonego: na punktach $t_n=(\pi(n+1/2))^{-1/2}$ przyjmuje naprzemienne wartości $\pm t_n$, a suma sąsiednich oscylacji jest rozbieżna. Pomimo tego dla $f(t)=t$ całka KS istnieje na mocy twierdzenia~\ref{tw:ks-bv-reg}, a wzór całkowania przez części daje
\[
 (\KS)\int_0^1 t\,\mathrm dg(t)
 =\sin1-\int_0^1 t\sin(1/t^2)\,\mathrm d\lambda(t)
 =\sin1-\frac12\int_1^\infty\frac{\sin u}{u^2}\,\mathrm du.
\]
Ostatnia całka jest bezwzględnie zbieżna i może być obliczana metodami numerycznymi. Nie potrzebujemy ani wahania ograniczonego funkcji $g$, ani klasycznej całki Riemanna--Stieltjesa.
\end{przyklad}

\begin{przyklad}[Dlaczego potrzebujemy wyrazu korekcyjnego]
Dla $H=\ind{[c,b]}$ mamy $\Delta^-H(c)=1$ i $\Delta^+H(c)=0$. Twierdzenie~\ref{tw:ks-ibp} z $f=g=H$ daje
$2(\KS)\int_a^b H\,\mathrm dH=H(b)^2-H(a)^2+1=2$. Dla $K=\ind{(c,b]}$ mamy natomiast $\Delta^-K(c)=0$, $\Delta^+K(c)=1$, zatem $2(\KS)\int_a^b K\,\mathrm dK=1-1=0$. Oba wyniki są zgodne z bezpośrednim wzorem $\int f\,\mathrm dH=f(c)$ i $\int f\,\mathrm dK=f(c)$.
\end{przyklad}

\subsection*{Całka nieoznaczona, podstawienie i błędy numeryczne}
\begin{twierdzenie}[Skoki funkcji pierwotnej KS]
\label{tw:ks-primitive}
Niech $f$ będzie regulowana i $g\in\BV[a,b]$. Zdefiniujmy $F(x)=(\KS)\int_a^x f(t)\,\mathrm dg(t)$ dla $x\in[a,b]$. Wtedy $F\in\BV[a,b]$, $F(a)=0$, $V_a^b(F)\le\|f\|_\infty V_a^b(g)$, a ponadto
\[
 F(c+)-F(c)=f(c)\Delta^+g(c),\qquad
 F(c)-F(c-)=f(c)\Delta^-g(c)
\]
odpowiednio dla $c\in[a,b)$ i $c\in(a,b]$.
\end{twierdzenie}
\begin{proof}
Z addytywności całki po przyległych przedziałach i z lematu~\ref{lem:ks-var} mamy dla każdego podziału $a=x_0<\cdots<x_n=b$ nierówność $\sum_i|F(x_i)-F(x_{i-1})|\le\|f\|_\infty\sum_i V_{x_{i-1}}^{x_i}(g)=\|f\|_\infty V_a^b(g)$. Dla $x\downarrow c$ rozdzielamy całkę $\int_c^x f\,\mathrm dg$ na atom w punkcie $c$, mający wartość $f(c)(g(c+)-g(c))$, oraz całkę względem pozostałej miary na $(c,x]$. Jej moduł jest ograniczony przez $\|f\|_\infty|\mu_g|((c,x])\to0$. Otrzymujemy pierwszy wzór. Dowód wzoru lewostronnego jest analogiczny.
\end{proof}

\begin{twierdzenie}[Podstawienie w przedziale]
\label{tw:ks-substitution}
Niech $\varphi:[A,B]\to[a,b]$ będzie ściśle rosnącym homeomorfizmem, $\varphi(A)=a$, $\varphi(B)=b$. Niech $f,g:[a,b]\to\R$ będą funkcjami o skończonych wartościach. Wtedy całka KS po jednej stronie istnieje wtedy i tylko wtedy, gdy istnieje po drugiej, oraz
\[
 (\KS)\int_a^b f(t)\,\mathrm dg(t)
 =(\KS)\int_A^B f(\varphi(s))\,\mathrm d(g\circ\varphi)(s).
\]
\end{twierdzenie}
\begin{proof}
Każdy podział znaczony $\mathcal Q=\{(\eta_i,[s_{i-1},s_i])\}$ przedziału $[A,B]$ przenosi się na podział $\mathcal P=\{(\varphi(\eta_i),[\varphi(s_{i-1}),\varphi(s_i)])\}$ przedziału $[a,b]$. Obie sumy KS są równe składnik po składniku. Niech $\delta$ będzie cechowaniem na $[a,b]$. Dzięki ciągłości $\varphi$ w każdym punkcie $s$ można wybrać $\widetilde\delta(s)>0$ tak, aby $|u-s|<\widetilde\delta(s)$ implikowało $|\varphi(u)-\varphi(s)|<\delta(\varphi(s))$. Wtedy każdy podział $\widetilde\delta$-drobny przenosi się na podział $\delta$-drobny, zatem istnienie i wartość całki po lewej implikują istnienie i tę samą wartość po prawej. Dla implikacji odwrotnej powtarzamy argument z $\varphi^{-1}$.
\end{proof}

\begin{twierdzenie}[Podstawienie całki nieoznaczonej]
\label{tw:ks-chain}
Niech $f$ będzie regulowana, $g\in\BV[a,b]$, $F(x)=(\KS)\int_a^x f\,\mathrm dg$ i niech $h:[a,b]\to\R$ będzie regulowana. Wówczas
\[
 (\KS)\int_a^b h(t)\,\mathrm dF(t)
 =(\KS)\int_a^b h(t)f(t)\,\mathrm dg(t).
\]
\end{twierdzenie}
\begin{proof}
Na mocy twierdzenia~\ref{tw:ks-measure} funkcja $F$ ma wahanie ograniczone, a odpowiadająca jej miara Stieltjesa spełnia $\mu_F(E)=\int_E f\,\mathrm d\mu_g$ dla każdego zbioru borelowskiego $E\subset[a,b]$. Wystarczy zweryfikować tę tożsamość na przedziałach $[a,x]$, uwzględniając prawostronne wartości $F(x+)$ i $g(x+)$; wówczas wynika ona z definicji $F$ oraz wzorów na jej skoki z twierdzenia~\ref{tw:ks-primitive}. Jednoznaczność miar daje równość na wszystkich zbiorach borelowskich. W konsekwencji $\int h\,\mathrm d\mu_F=\int hf\,\mathrm d\mu_g$. Z twierdzenia~\ref{tw:ks-measure} obie strony są wskazanymi całkami KS.
\end{proof}

\begin{twierdzenie}[Oszacowanie błędu zwykłych sum]
\label{tw:ks-quadrature}
Niech $f\in C[a,b]$, $g\in\BV[a,b]$ oraz $P=\{a=x_0<\cdots<x_n=b\}$ będzie dowolnym podziałem o średnicy $h$. Dla dowolnych znaczników $\xi_i\in[x_{i-1},x_i]$ zachodzi
\[
 \left|(\KS)\int_a^b f\,\mathrm dg-
 \sum_{i=1}^n f(\xi_i)\bigl(g(x_i)-g(x_{i-1})\bigr)\right|
 \le\omega_f(h)\,V_a^b(g),
\]
gdzie $\omega_f(h)=\sup\{|f(x)-f(y)|:x,y\in[a,b],\ |x-y|\le h\}$ jest modułem ciągłości. Jeśli $f$ jest $L$-lipschitzowska, błąd jest nie większy niż $LhV_a^b(g)$; jeśli jest h\"olderowska ze stałą $K$ i wykładnikiem $\theta\in(0,1]$, błąd nie przekracza $Kh^\theta V_a^b(g)$.
\end{twierdzenie}
\begin{proof}
Z addytywności całki KS i liniowości mamy
\[
 (\KS)\int_a^b f\,\mathrm dg-S_g(f,\mathcal P)
 =\sum_{i=1}^n(\KS)\int_{x_{i-1}}^{x_i}\bigl(f(t)-f(\xi_i)\bigr)\,\mathrm dg(t).
\]
Funkcja $f-f(\xi_i)$ ma na $[x_{i-1},x_i]$ moduł nieprzekraczający $\omega_f(h)$. Lemat~\ref{lem:ks-var} na każdym przedziale daje oszacowanie przez $\omega_f(h)V_{x_{i-1}}^{x_i}(g)$. Sumując i korzystając z addytywności wahania całkowitego, otrzymujemy żądaną nierówność. Dwa ostatnie oszacowania wynikają z definicji modułu ciągłości.
\end{proof}

\begin{uwaga}[Czego nie daje sama średnica podziału]
Jeżeli $f$ ma skok w punkcie, w którym skacze $g$, mała średnica podziału nie gwarantuje przybliżenia całki przez dowolne sumy. Dla $f=g=\ind{[c,b]}$ sumy na bardzo drobnych podziałach mogą być równe $0$ albo $1$ w zależności od znaczników. Właśnie dlatego w definicji KS stosujemy cechowania zależne od punktu. Oszacowanie z twierdzenia~\ref{tw:ks-quadrature} wymaga ciągłości $f$.
\end{uwaga}

\begin{przyklad}[Przedziałowe oszacowanie dla kwadratury]
Na $[0,1]$ niech $f(t)=\sin t$ i $g$ będzie dowolną niemalejącą funkcją o $g(0)=0$, $g(1)=3$. Dla podziału o średnicy $h=10^{-3}$ dowolna suma Stieltjesa z wartościami $\sin\xi_i$ różni się od całki KS o co najwyżej $3\cdot10^{-3}$, gdyż $|f'|\le1$. Nie potrzebujemy pochodnej $g$ ani wiedzy o położeniu jej skoków.
\end{przyklad}

\subsection*{Granice wzorów i porównanie z całką HK}
\begin{stwierdzenie}[Nie każda całka KS jest całką Lebesgue'a--Stieltjesa]
Dla $g(t)=t$ na $[0,1]$ całka KS jest całką Henstocka--Kurzweila. Istnieją funkcje $f:[0,1]\to\R$, które są KS-całkowalne względem $g$, lecz nie należą do $L^1([0,1],\lambda)$. W szczególności równość z twierdzenia~\ref{tw:ks-measure} nie obowiązuje bez założeń o $f$.
\end{stwierdzenie}
\begin{proof}
Niech $F(0)=0$ i $F(t)=t^2\cos(\pi/t^2)$ dla $t\in(0,1]$ oraz $f(0)=0$, $f(t)=F'(t)=2t\cos(\pi/t^2)+(2\pi/t)\sin(\pi/t^2)$ dla $t>0$. Z twierdzenia Newtona--Leibniza dla całki HK otrzymujemy $(\KS)\int_0^1 f\,\mathrm dt=F(1)-F(0)=-1$. Drugi składnik pochodnej nie jest bezwzględnie całkowalny w otoczeniu zera, ponieważ po podstawieniu $u=\pi/t^2$ jego całka bezwzględna jest porównywalna z $\int^\infty|\sin u|/u\,\mathrm du=\infty$. Tymczasem pierwszy składnik należy do $L^1$. Zatem $\int_0^1|f|\,\mathrm d\lambda=\infty$, więc całka Lebesgue'a z $f$ jako funkcji rzeczywistej nie istnieje.
\end{proof}

\begin{uwaga}[Schemat wyboru wzoru do obliczenia]
Gdy $g$ jest schodkowa, stosujemy twierdzenie~\ref{tw:ks-step}. Gdy $g$ ma wahanie ograniczone, rozkładamy jej miarę według twierdzenia~\ref{tw:ks-decomp}. Gdy $f$ jest absolutnie ciągła, a $g$ jedynie regulowana, najlepiej zastosować wniosek~\ref{wn:ks-ibp-ac}. Przy ciągłej funkcji podcałkowej i całkującej o wahaniu ograniczonym można korzystać z oszacowania kwadraturowego~\ref{tw:ks-quadrature}. Dla dowolnego integratora i funkcji bez powyższej regularności nie istnieje jeden uniwersalny wzór redukujący całkę KS do zwykłej całki Lebesgue'a.
\end{uwaga}


\section{Całka McShane'a}
\providecommand{\atlasref}[2]{#2}
% Import z AM1-36, zachowane szczególne konstrukcje skryptu.
\begin{definicja}\label{it:def:mcshane}
Niech $a<b$ i niech $f:[a,b]\to\R$ ma skończone wartości. Podział znaczony $\mathcal P=\{(t_i,[x_{i-1},x_i])\}$ jest \emph{wolny $\delta$-drobny}, gdy $t_i\in[a,b]$ oraz $[x_{i-1},x_i]\subset(t_i-\delta(t_i),t_i+\delta(t_i))$, bez wymagania $t_i\in[x_{i-1},x_i]$. Funkcja $f$ jest całkowalna w sensie McShane'a, $f\in\mathcal M[a,b]$, jeśli istnieje $I\in\R$ takie, że dla każdego $\varepsilon>0$ istnieje cechowanie $\delta:[a,b]\to(0,\infty)$, dla którego każdy wolny $\delta$-drobny podział spełnia
\[
\left|S(f,\mathcal P)-I\right|<\varepsilon,\qquad
 S(f,\mathcal P)=\sum_i f(t_i)(x_i-x_{i-1}).
\]
Piszemy wówczas $I=(\mathcal M)\int_a^b f$. Na przedziale jednopunktowym przyjmujemy całkę zero.
\end{definicja}
\begin{twierdzenie}[Kryterium Cauchy'ego, liniowość i porządek]\label{it:thm:mcshane-basic}
Dla funkcji $f:[a,b]\to\R$ o skończonych wartościach
całka McShane'a istnieje wtedy i tylko wtedy, gdy dla każdego
$\varepsilon>0$ istnieje cechowanie $\delta$ takie, że
$|S(f,\mathcal P)-S(f,\mathcal Q)|<\varepsilon$ dla wszystkich
wolnych $\delta$-drobnych podziałów $\mathcal P,\mathcal Q$.
Całka jest jednoznaczna, liniowa i monotoniczna.
Na przedziale jednopunktowym przyjmujemy całkę zero.
\end{twierdzenie}
\begin{proof}
Każdy podział HK jest wolnym podziałem, zatem lemat Cousina
\ref{it:lem:kh-cousin} zapewnia istnienie wolnego podziału
dla dowolnego cechowania. Konieczność kryterium wynika
z nierówności trójkąta. Dla dostateczności wybieramy cechowania
$\delta_n$ dla błędów $2^{-n}$, zastępując je kolejno
minimami z poprzednimi. Wybieramy wolne $\delta_n$-drobne
podziały $\mathcal P_n$. Dla $m\ge n$ obie sumy są drobne
względem $\delta_n$, więc ich różnica ma moduł mniejszy
od $2^{-n}$. Sumy tworzą ciąg Cauchy'ego o granicy $I$.
Porównując dowolny wolny $\delta_n$-drobny podział
z $\mathcal P_m$ dla $m\ge n$ i przechodząc do granicy,
otrzymujemy błąd względem $I$ nie większy od $2^{-n}$. Dla dowolnej wymaganej dokładności $\varepsilon>0$ wybieramy $n$ takie, że $2^{-n}<\varepsilon$. Cechowanie $\delta_n$ spełnia wtedy definicję całki McShane'a z wartością $I$.

Jeśli dwie liczby spełniają definicję całki, wybieramy
wspólne cechowanie dla błędu $\varepsilon/2$ i jeden drobny
podział; różnica tych liczb ma moduł mniejszy od $\varepsilon$.
Dla dwóch funkcji całkowalnych $f,g$ i współczynników $\alpha,\beta\in\R$ wybieramy minimum ich cechowań dla błędu $\varepsilon/(1+|\alpha|+|\beta|)$. Liniowość skończonych sum daje
\[
 |S(\alpha f+\beta g,\mathcal P)-\alpha I_f-\beta I_g|
 \le |\alpha||S(f,\mathcal P)-I_f|
     +|\beta||S(g,\mathcal P)-I_g|<\varepsilon.
\]
To dowodzi istnienia całki kombinacji liniowej oraz jej wartości.
Dla $f\le g$ mamy $S(f,\mathcal P)\le S(g,\mathcal P)$
na każdym podziale, ponieważ długości komórek są nieujemne.
Dla $\varepsilon>0$ wybieramy wspólne cechowanie dla błędu $\varepsilon/2$ i wolny drobny podział (istniejący na mocy lematu Cousina). Wtedy $I_f\le S(f,\mathcal P)+\varepsilon/2\le S(g,\mathcal P)+\varepsilon/2\le I_g+\varepsilon$. Dowolność $\varepsilon$ daje $I_f\le I_g$.
\end{proof}

\begin{stwierdzenie}
$\mathcal R[a,b]\subset\mathcal M[a,b]\subset\HK[a,b]$, a wartości całek są zgodne.
\end{stwierdzenie}
\begin{proof}
Druga inkluzja wynika z warunków na znaczniki. Dla pierwszej ustalmy $M=\|f\|_\infty$ oraz podział $P_0$ z $U(f,P_0)-L(f,P_0)<\varepsilon/2$. Niech $B$ będzie skończonym zbiorem punktów jego podziału. Dla $t\notin B$ wybieramy $\delta(t)<\operatorname{dist}(t,B)$, zaś w punktach $t\in B$ dobieramy dodatnie cechowania tak małe, aby $4M\sum_{t\in B}\delta(t)<\varepsilon/2$. Każdy przedział wolnego podziału ze znacznikiem poza $B$ leży w tej samej komórce $P_0$ co jego znacznik. Oznaczmy komórki $P_0$ przez $J_k$, a oscylację $f$ na $J_k$ przez $\omega_k=\sup_{J_k}f-\inf_{J_k}f$. Z lokalności całki Riemanna (twierdzenie \ref{ma1:thm:riemann-additivity}) mamy dla każdej komórki $I_i$
\[
 S(f,\mathcal P)-\int_a^b f=\sum_i\left(f(t_i)|I_i|-\int_{I_i} f\right).
\]
Gdy $t_i\notin B$ oraz $t_i\in J_k$, $I_i\subset J_k$, oszacowanie całki przez kresy funkcji daje moduł składnika nie większy niż $\omega_k|I_i|$. Łączna długość tych komórek w $J_k$ jest nie większa niż $|J_k|$, więc suma błędów dla znaczników spoza $B$ jest nie większa niż $\sum_k\omega_k|J_k|=U(f,P_0)-L(f,P_0)$. Dla ustalonego $t\in B$ wszystkie komórki z tym znacznikiem mają rozłączne wnętrza i leżą w otoczeniu długości $2\delta(t)$, nawet jeśli znacznik do nich nie należy. Moduł każdego składnika jest nie większy niż $2M|I_i|$, więc ich łączny błąd nie przekracza $4M\delta(t)$. Dodając oszacowania, dostajemy błąd mniejszy niż $\varepsilon$.
\end{proof}
\noindent\emph{Porównanie w Bibliotece matematycznej:} \atlasref{AM1-36::it:thm:mcshane-absolute}{AM1-36: bezwzględna całkowalność w sensie McShane'a}.

\begin{twierdzenie}[Bezwzględna całkowalność]\label{it:thm:mcshane-absolute}
Jeżeli $f\in\mathcal M[a,b]$, to $|f|\in\mathcal M[a,b]$. Tak więc całka McShane'a jest bezwzględna, mimo że definiujemy ją przez sumy znaczonych przedziałów.
\end{twierdzenie}
\begin{proof}
Niech $I=\int_a^b f$ w sensie McShane'a. Wybieramy cechowanie $\delta$ zapewniające błąd mniejszy niż $\varepsilon/3$. Dla dwóch dowolnych wolnych $\delta$-drobnych podziałów $\mathcal P=\{(t_i,I_i)\}$ i $\mathcal Q=\{(s_j,J_j)\}$ rozważamy wspólne zagęszczenie $K_{ij}=I_i\cap J_j$, pomijając przecięcia zerowej długości. Każda komórka $K_{ij}$ leży w otoczeniach cechowania zarówno $t_i$, jak i $s_j$, więc obydwa znaczniki są dopuszczalne. Wybierając w jednej sumie znacznik z większą wartością $f$, a w drugiej z mniejszą, otrzymujemy dwie dopuszczalne sumy $S_+,S_-$, zatem
\[
\sum_{i,j}|f(t_i)-f(s_j)|\,|K_{ij}|=S_+-S_-<2\varepsilon/3.
\]
Ponieważ $|I_i|=\sum_j|K_{ij}|$, $|J_j|=\sum_i|K_{ij}|$ i $\bigl||u|-|v|\bigr|\le|u-v|$, mamy
$|S(|f|,\mathcal P)-S(|f|,\mathcal Q)|\le\sum_{i,j}|f(t_i)-f(s_j)|\,|K_{ij}|<\varepsilon$.
Kryterium Cauchy'ego z twierdzenia \ref{it:thm:mcshane-basic} dowodzi więc istnienia całki McShane'a z $|f|$. Dla HK podobny argument nie działa: znacznik pierwotnej komórki nie musi należeć do jej przecięcia z inną komórką, co narusza warunek HK.
\end{proof}

\begin{przyklad}[Inkluzje są właściwe]
Funkcja Dirichleta należy do $\mathcal M[0,1]$ (lemat \ref{ma1:kh:lem:kh-null-modification} działa także dla wolnych znaczników), ale nie do $\mathcal R[0,1]$, ponieważ jest nieciągła w każdym punkcie (twierdzenie \ref{thm:riemann-criterion}). Funkcja $x^{-1/2}$ po określeniu w $0$ również należy do $\mathcal M\setminus\mathcal R$. Pochodna niecałkowalna bezwzględnie z rozdziału o całce HK należy do $\HK\setminus\mathcal M$; wynika to z twierdzenia \ref{it:thm:mcshane-lebesgue}. To samo twierdzenie uzasadnia całkowalność McShane'a funkcji $x^{-1/2}$, której całka Lebesgue'a wynosi $2$: funkcje $x^{-1/2}\chi_{[1/n,1]}(x)$, przy wartości zero w $0$, rosną prawie wszędzie do tej funkcji i mają całki $2(1-1/\sqrt n)$, co przez twierdzenie \ref{thm:mct-ae} daje żądaną wartość. Brak całkowalności Riemanna wynika z nieograniczoności.
\end{przyklad}

\begin{uwaga}[Kolejność rozszerzeń]
Całka McShane'a rozszerza Riemanna i jest zawarta w HK. W teorii funkcji rzeczywistych opisuje dokładnie całkę Lebesgue'a; wykażemy to dalej. Wolność znaczników powiększa rodzinę sum, wobec których wymagamy jednakowych oszacowań, więc \emph{zmniejsza} klasę funkcji całkowalnych w porównaniu z HK.
\end{uwaga}

% ============================================================
% UZUPELNIENIE MATEMATYCZNE: ROZDZIAL 7
% ============================================================
\subsection*{Konstrukcja całki McShane'a dla funkcji charakterystycznych}
\begin{twierdzenie}[Funkcje charakterystyczne zbiorów mierzalnych]\label{tw:mcshane-indykator}
Niech $E\subset[a,b]$ będzie zbiorem mierzalnym w sensie Lebesgue'a. Wówczas $\ind E$ jest całkowalna w sensie McShane'a i
$(\mathcal M)\int_a^b\chi_E(t)\,\mathrm dt=\lambda_1(E)$.
\end{twierdzenie}
\begin{proof}
Ustalmy $\varepsilon>0$. Z regularności miary Lebesgue'a wybieramy zbiór zwarty $K\subset E$ i zbiór otwarty $O\supset E$ (w topologii względnej $[a,b]$), dla których $\lambda_1(O\setminus K)<\varepsilon$. Konstruujemy cechowanie: dla $t\in E$ wybieramy $\delta(t)>0$ tak, by $(t-\delta(t),t+\delta(t))\cap[a,b]\subset O$; dla $t\notin E$ dobieramy je tak, by wymienione otoczenie nie przecinało $K$ (jest to możliwe, ponieważ $K$ jest domknięty i $t\notin K$).

Niech $\mathcal P$ będzie dowolnym wolnym $\delta$-drobnym podziałem. Każda jego komórka z oznaczeniem $t\in E$ leży w $O$, a każda komórka ze znacznikiem $t\notin E$ nie przecina $K$. Ponieważ komórki pokrywają cały przedział, suma długości komórek ze znacznikami w $E$ mieści się między $\lambda_1(K)$ a $\lambda_1(O)$. Jest ona dokładnie sumą McShane'a $S(\ind E,\mathcal P)$. Wobec $\lambda_1(K)\le\lambda_1(E)\le\lambda_1(O)$ otrzymujemy $|S-\lambda_1(E)|<\varepsilon$.
\end{proof}
\begin{wniosek}[Funkcje proste]\label{wn:mcshane-proste}
Niech $s=\sum_{k=1}^m c_k\ind{E_k}$, gdzie $c_k\in\R$ i wszystkie $E_k\subset[a,b]$ są mierzalne. Wtedy $s\in\mathcal M[a,b]$ oraz $(\mathcal M)\int_a^b s(t)\,\mathrm dt=\sum_{k=1}^m c_k\lambda_1(E_k)$.
\end{wniosek}
\begin{proof}
Każdy składnik jest całkowalny na mocy twierdzenia \ref{tw:mcshane-indykator}; liniowość sum jest bezpośrednią konsekwencją definicji całki McShane'a.
\end{proof}
\begin{przyklad}[Zbiór Cantora]
Dla klasycznego zbioru Cantora $K\subset[0,1]$ mamy $\lambda_1(K)=0$, więc $(\mathcal M)\int_0^1\chi_K(t)\,\mathrm dt=0$. Ten sam wynik zachodzi dla funkcji Dirichleta z racji $\lambda_1(\Q\cap[0,1])=0$. Funkcja charakterystyczna klasycznego zbioru Cantora jest też całkowalna w sensie Riemanna, ponieważ zbiór jej nieciągłości jest zawarty w zbiorze Cantora i ma miarę zero. Funkcja Dirichleta nie jest całkowalna w sensie Riemanna, gdyż jest nieciągła w każdym punkcie.
\end{przyklad}
\begin{uwaga}[Zakres uogólnienia]
Twierdzenie o funkcjach charakterystycznych daje od razu zgodność całki McShane'a z całką Lebesgue'a na funkcjach prostych. Przejście do wszystkich funkcji z $L^1$ nie jest wyłącznie konsekwencją liniowości: wymaga dodatkowej kontroli sum dla przybliżeń i jest przeprowadzone w twierdzeniu \ref{it:thm:mcshane-lebesgue} w niniejszym rozdziale.
\end{uwaga}

\subsection{Własności i twierdzenia graniczne}
\begin{wniosek}[Działania kratowe]\label{it:cor:mcshane-lattice}
Dla $f,g\in\mathcal M[a,b]$ także $f^+$, $f^-$,
$\max(f,g)$ i $\min(f,g)$ należą do $\mathcal M[a,b]$.
\end{wniosek}
\begin{proof}
Stosujemy twierdzenia \ref{it:thm:mcshane-basic}
i \ref{it:thm:mcshane-absolute} do tożsamości
$f^+=(|f|+f)/2$ i $f^-=(|f|-f)/2$
oraz $\max(f,g)=(f+g+|f-g|)/2$ i
$\min(f,g)=(f+g-|f-g|)/2$.
\end{proof}

\begin{twierdzenie}[Wspólne cechowanie]\label{it:thm:mcshane-equiintegrability}
Niech $f_n\in\mathcal M[a,b]$ z całkami $I_n$ i niech
$f_n(t)\to f(t)\in\R$ dla każdego $t$.
Załóżmy, że dla każdego $\varepsilon>0$ istnieje jedno
cechowanie $\delta$ takie, że $|S(f_n,\mathcal P)-I_n|<\varepsilon$
dla wszystkich $n$ i wszystkich wolnych $\delta$-drobnych
podziałów $\mathcal P$. Wówczas $f\in\mathcal M[a,b]$
oraz $I_n\to(\mathcal M)\int_a^b f$.
\end{twierdzenie}
\begin{proof}
Dla ustalonego $\varepsilon>0$ wybieramy wspólne cechowanie
dla błędu $\varepsilon/3$ i jeden wolny drobny podział.
Jego skończona suma dla $f_n$ zbiega do sumy dla $f$,
ponieważ zbieżność zachodzi we wszystkich znacznikach.
Dla dostatecznie dużych $n,m$ moduł różnicy sum jest
mniejszy od $\varepsilon/3$, więc z nierówności trójkąta
$|I_n-I_m|<\varepsilon$. Niech $I=\lim_n I_n$.
Dla dowolnej dokładności $\eta>0$ wybieramy wspólne
cechowanie dla błędu $\eta/2$. Dla każdego wolnego
drobnego podziału $\mathcal Q$ przejście do granicy
w nierówności $|S(f_n,\mathcal Q)-I_n|<\eta/2$ daje
$|S(f,\mathcal Q)-I|\le\eta/2<\eta$.
\end{proof}

\begin{wniosek}[Zbieżność jednostajna]\label{it:cor:mcshane-uniform}
Jeżeli $f_n\in\mathcal M[a,b]$ i $f_n\to f$ jednostajnie,
to $f\in\mathcal M[a,b]$ i całki zbiegają.
\end{wniosek}
\begin{proof}
Dla dostatecznie dużych $n,m$ różnica $f_n-f_m$ jest
ograniczona. Z oszacowania sum oraz liniowości wynika
$|I_n-I_m|\le(b-a)\|f_n-f_m\|_\infty$.
Stąd $I=\lim I_n$ istnieje. Dla $\varepsilon>0$ wybieramy
$n$ tak, aby $(b-a)\|f-f_n\|_\infty<\varepsilon/3$
oraz $|I_n-I|<\varepsilon/3$. Na wolnych podziałach
drobnych względem cechowania dla $f_n$ z błędem
$\varepsilon/3$ mamy
\[
 |S(f,\mathcal P)-I|
 \le (b-a)\|f-f_n\|_\infty+|S(f_n,\mathcal P)-I_n|
     +|I_n-I|<\varepsilon.
\]
Nie zakładamy z góry całkowalności granicy ani ograniczoności
poszczególnych funkcji.
\end{proof}

\begin{wniosek}[Lokalność, zbieżność monotoniczna i zmajoryzowana]
\label{it:cor:mcshane-lebesgue-consequences}
Całkowalność McShane'a na $[a,b]$ pociąga całkowalność
na każdym podprzedziale; całka jest addytywna przy
sklejaniu dwóch sąsiednich przedziałów.
Jeśli $f_n\in\mathcal M[a,b]$ oraz $f_n\to f$ prawie
wszędzie, zbieżność całek i całkowalność $f$ zachodzą
w każdym z następujących przypadków:
\begin{enumerate}[label=(\alph*)]
\item $|f_n|\le g$ prawie wszędzie dla pewnej
$g\in\mathcal M[a,b]$;
\item $f_n\uparrow f$ prawie wszędzie i $\sup_n I_n<\infty$,
gdzie $I_n=(\mathcal M)\int_a^b f_n$.
\end{enumerate}
Zakładamy skończone wartości funkcji $f$; na zbiorze
wyjątkowym wolno nadać jej dowolne skończone wartości.
\end{wniosek}
\begin{proof}
Najpierw dowodzimy lokalności bez reprezentacji miarowej. Dla $a<c<b$ wybieramy cechowanie z kryterium Cauchy'ego na $[a,b]$. Dwa wolne drobne podziały $[a,c]$ uzupełniamy tym samym drobnym podziałem $[c,b]$, istniejącym dzięki lematowi Cousina. Po odjęciu wspólnej sumy kryterium Cauchy'ego daje całkowalność na $[a,c]$. Zamiana ról i powtórzenie argumentu daje całkowalność na każdym podprzedziale. Odwrotnie, jeżeli istnieją całki na obu częściach, wybieramy na nich cechowania dla błędu $\varepsilon/2$. Dla $t\ne c$ ograniczamy odpowiednie cechowanie przez $|t-c|/2$, a dla $t=c$ bierzemy minimum obu cechowań. Żadna komórka ze znacznikiem $t\ne c$ nie może leżeć po przeciwnej stronie $c$ niż jej znacznik, nawet gdy znacznik nie należy do komórki. Komórkę przecinającą $c$ rozcinamy w $c$, zachowując znacznik $c$. Po rozcięciu otrzymujemy wolne drobne podziały obu części o niezmienionej łącznej sumie. Nierówność trójkąta daje błąd mniejszy od $\varepsilon$ względem sumy całek lokalnych. To dowodzi również addytywności.

Dla twierdzeń granicznych stosujemy identyfikację \ref{it:thm:mcshane-lebesgue}: wszystkie funkcje McShane'a są całkowalne w sensie Lebesgue'a i odwrotnie. Wybieramy jeden mierzalny zbiór miary zero zawierający wszystkie zbiory wyjątkowe z założeń. Poza nim $f$ jest granicą funkcji mierzalnych; z zupełności miary Lebesgue'a wynika mierzalność $f$. W (a) stosujemy twierdzenie \ref{thm:dct} z majorantą $|g|$. W (b) funkcje $f_n-f_1$ są nieujemne prawie wszędzie i rosną do $f-f_1$ prawie wszędzie. Ich całki są ograniczone przez $\sup_n I_n-I_1$, więc twierdzenie \ref{thm:mct-ae} daje całkowalność $f-f_1$ i zbieżność jego całek. Dodając $f_1$, dostajemy całkowalność $f$ i zbieżność $I_n$.
Ponowna identyfikacja całek przez twierdzenie
\ref{it:thm:mcshane-lebesgue} kończy dowód.
\end{proof}

\subsection{Zgodność całek McShane'a, Lebesgue'a i HK}
\begin{lemat}[Bezwzględna ciągłość całki]\label{it:lem:integral-absolute-continuity}\label{ma1:kh:lem:integral-ac}
Jeśli $f\in\mathcal L^1(\mu)$, to dla każdego $\varepsilon>0$ istnieje $\eta>0$ takie, że $\mu(A)<\eta$ pociąga $\int_A|f|\,\mathrm d\mu<\varepsilon$ dla każdego $A\in\mathcal F$ w przestrzeni z miarą $(X,\mathcal F,\mu)$.
\end{lemat}
\begin{proof}
Wybieramy $M>0$ z $\int_{\{|f|>M\}}|f|\,\mathrm d\mu<\varepsilon/2$ oraz $\eta=\varepsilon/(2M)$. Wtedy $\int_A|f|\,\mathrm d\mu\le M\mu(A)+\int_{\{|f|>M\}}|f|\,\mathrm d\mu<\varepsilon$. Istotnie, funkcje $|f|\chi_{\{|f|\le n\}}$ rosną prawie wszędzie do $|f|$ (wartości nieskończone występują tylko na zbiorze zerowym). Z twierdzenia \ref{thm:mct-ae} ich całki dążą do skończonej całki $|f|$. Z addytywności całki po rozłącznych zbiorach otrzymujemy $\int_{\{|f|>n\}}|f|\,\mathrm d\mu\to0$, co uzasadnia wybór $M$.
\end{proof}

\begin{twierdzenie}[McShane = Lebesgue]\label{it:thm:mcshane-lebesgue}
Niech $a<b$. Dla funkcji rzeczywistych zachodzi
$\mathcal M[a,b]=\mathcal L^1([a,b],\lambda_1)$
jako równość zbiorów konkretnych funkcji o skończonych
wartościach, gdzie $\lambda_1$ jest miarą uzupełnioną.
Dowolne zmiany skończonych wartości na zbiorach
miary zero zachowują całkowalność i wartość całki.
Wartości obu całek są równe.
\end{twierdzenie}
\begin{proof}
\emph{Krok 1: Lebesgue $\Rightarrow$ McShane.} Niech $f\in\mathcal L^1$ i $\varepsilon>0$. Wybieramy $h>0$ z $h(b-a)<\varepsilon/3$ i rozbijamy dziedzinę na mierzalne zbiory $E_k=\{kh\le f<(k+1)h\}$, $k\in\mathbb Z$. Z bezwzględnej ciągłości całki wybieramy $\eta>0$ tak, że $\lambda_1(A)<\eta$ pociąga $\int_A|f|\,\mathrm d\lambda_1<\varepsilon/3$. Ustalamy dodatnie liczby $w_k=2^{-|k|}/3$ ($k\in\Z$), których suma wynosi $1$. Z regularności miary wybieramy względnie otwarte $G_k\supset E_k$ z
\[
 \lambda_1(G_k\setminus E_k)<
 w_k\min\left\{\eta/2,\frac{\varepsilon}{6(1+|kh|+h)}\right\}.
\]
Dla $U=\bigcup_k(G_k\setminus E_k)$ mamy wtedy $\lambda_1(U)\le\eta/2$ i $\int_U|f|\,\mathrm d\lambda_1<\varepsilon/3$, a suma $\sum_k(|kh|+h)\lambda_1(G_k\setminus E_k)$ nie przekracza $\varepsilon/6$. Dla $t\in E_k$ wybieramy $\delta(t)>0$ z $(t-\delta(t),t+\delta(t))\cap[a,b]\subset G_k$.

Niech $\mathcal P=\{(t_i,I_i)\}$ będzie dowolnym wolnym $\delta$-drobnym podziałem. Zatem $I_i\subset G_k$ dla $t_i\in E_k$, również gdy $t_i\notin I_i$. Ponieważ komórki mają rozłączne wnętrza, a ich wspólne końce są zbiorami zerowymi, z addytywności całki Lebesgue'a mamy
\[
 S(f,\mathcal P)-\int_a^b f(x)\,\mathrm d\lambda_1(x)
 =\sum_i\int_{I_i}\bigl(f(t_i)-f(x)\bigr)\,\mathrm d\lambda_1(x).
\]
Przy znaczniku $t_i\in E_k$ na części $I_i\cap E_k$ moduł różnicy wartości $f(t_i)$ i $f(x)$ jest mniejszy od $h$, ponieważ obie wartości należą do $[kh,(k+1)h)$. Pozostałe mierzalne części $I_i\setminus E_k$ leżą w $U$; ich przecięcia zawierają co najwyżej wspólne końce komórek, a więc mają miarę zero. Dla ustalonego $k$ suma ich miar po komórkach ze znacznikami w $E_k$ nie przekracza $\lambda_1(G_k\setminus E_k)$. Po rozdzieleniu sumy otrzymujemy
\[
\left|S(f,\mathcal P)-\int_a^b f\,\mathrm d\lambda_1\right|
\le h(b-a)+\sum_k(|kh|+h)\lambda_1(G_k\setminus E_k)+\int_U|f|\,\mathrm d\lambda_1<\varepsilon.
\]

\emph{Krok 2: McShane $\Rightarrow$ Lebesgue.} Ponieważ każdy podział ze znacznikami należącymi do komórek jest również podziałem z wolnymi znacznikami, całkowalność McShane'a pociąga całkowalność HK z tą samą wartością całki. Mierzalność $f$ wynika z twierdzenia \ref{it:thm:kh-primitive-derivative}. Z twierdzenia \ref{it:thm:mcshane-absolute} także $|f|\in\mathcal M$, z całką $J\ge0$. Obcięcia $h_n=\min(|f|,n)$ są ograniczone i mierzalne, zatem Lebesgue'a całkowalne, a na mocy kroku 1 także McShane'a całkowalne. Monotoniczność całki McShane'a (twierdzenie \ref{it:thm:mcshane-basic}) i zgodność z kroku 1 dają $\int h_n\,\mathrm d\lambda_1\le J$. Przez zbieżność monotoniczną $\int|f|\,\mathrm d\lambda_1\le J<\infty$. Krok 1 i jednoznaczność całki McShane'a
(twierdzenie \ref{it:thm:mcshane-basic})
dowodzą zgodności wartości. Jeżeli zmienimy
skończone wartości $f$ na zbiorze miary zero,
to zupełność miary zapewnia mierzalność zmodyfikowanej
funkcji, a twierdzenie \atlasref{ME-07::thm:ae-integral-invariance}{ME-07: niezmienniczość całki przy równości prawie wszędzie}
zapewnia równość całek Lebesgue'a i zachowanie
całkowalności. Z wykazanej równoważności wynika to
samo dla McShane'a. Ponieważ każdy podział HK jest
wolnym podziałem, w szczególności dla każdej
$f\in\mathcal L^1$ zachodzi $f\in\HK$ i
$(\HK)\int_a^b f=\int_{[a,b]}f\,\mathrm d\lambda_1$.
Nie korzystaliśmy tutaj z kryterium bezwzględnej
całkowalności HK, które udowodnimy poniżej.
Dla $a=b$ każda skończona funkcja na singletonie
jest mierzalna, całkowalna i obie całki wynoszą zero.
\end{proof}
\begin{twierdzenie}[Kryterium bezwzględnej całkowalności HK]\label{ma1:kh:thm:kh-absolute}\label{it:thm:hk-absolute-integrability}
Niech $f:[a,b]\to\R$ będzie funkcją o skończonych
wartościach na przedziale zwartym. Jeżeli $f\ge0$
jest HK-całkowalna, to $f\in\mathcal L^1$.
Dla dowolnej takiej funkcji mamy
\[
f\in\mathcal L^1\quad\Longleftrightarrow\quad f\in\HK\ \text{i}\ |f|\in\HK.
\]
Każda ograniczona funkcja na $[a,b]$ jest
HK-całkowalna wtedy i tylko wtedy, gdy jest
mierzalna w sensie Lebesgue'a. W przypadku Riemanna
konieczny i wystarczający jest warunek, aby zbiór
punktów nieciągłości miał miarę zero
(twierdzenie \ref{thm:riemann-criterion}).
\end{twierdzenie}
\begin{proof}
Z twierdzenia \ref{it:thm:kh-primitive-derivative} funkcja HK-całkowalna jest mierzalna. Dla nieujemnej $f$ obcięcia $h_n=\min(f,n)$ są mierzalne, ograniczone i całkowalne Lebesgue'a, więc należą także do HK. Z monotoniczności HK mamy $\int h_n\,\mathrm d\lambda_1\le(\HK)\int f<\infty$. Zbieżność monotoniczna daje $f\in\mathcal L^1$. Zastosowanie pierwszej części do $|f|$ oraz inkluzja $\mathcal L^1\subset\HK$ dają równoważność. Ostatnia teza wynika z mierzalności funkcji HK i skończoności miary przedziału.
\end{proof}

\begin{wniosek}[Części dodatnia i ujemna w kryterium HK]
\label{it:cor:hk-positive-negative-integrability}
Dla każdej funkcji rzeczywistej $f$ o skończonych
wartościach na $[a,b]$ następujące warunki są równoważne:
\begin{enumerate}[label=(\roman*)]
\item $f\in\mathcal L^1([a,b],\lambda_1)$;
\item $f^+,f^-\in\HK[a,b]$;
\item $f,|f|\in\HK[a,b]$;
\item $f\in\mathcal M[a,b]$.
\end{enumerate}
Jeżeli zachodzi którykolwiek z nich, to
\[
 (\HK)\int_a^b|f|=
 \int_{[a,b]}|f|\,\mathrm d\lambda_1
 =(\HK)\int_a^b f^++(\HK)\int_a^b f^-.
\]
Jeśli natomiast $f\in\HK[a,b]\setminus\mathcal L^1$,
to żadna z funkcji $f^+,f^-$ nie jest HK-całkowalna,
a obie ich całki Lebesgue'a wynoszą $+\infty$.
\end{wniosek}
\begin{proof}
Z (i) wynika (ii), ponieważ $0\le f^\pm\le|f|$,
więc obie części są Lebesgue'a całkowalne, a
$\mathcal L^1\subset\HK$ z twierdzenia
\ref{it:thm:mcshane-lebesgue} i jego dowodu.
Z (ii) wynika (iii) przez liniowość HK i tożsamości
$f=f^+-f^-$, $|f|=f^++f^-$. Implikacja
(iii)$\Rightarrow$(i) jest treścią twierdzenia
\ref{it:thm:hk-absolute-integrability}, a równoważność
(i) i (iv) wynika z twierdzenia
\ref{it:thm:mcshane-lebesgue}. Wzór na całkę modułu
otrzymujemy z zgodności wartości całek i addytywności
całki nieujemnej.

Jeżeli $f\in\HK$ i jedna z funkcji $f^+,f^-$
należy do HK, to druga również należy do HK,
ponieważ $f^-=f^+-f$, $f^+=f+f^-$.
Z (ii)$\Rightarrow$(i) wynika wówczas $f\in\mathcal L^1$.
Dla $f\in\HK\setminus\mathcal L^1$ obie części są
mierzalne (twierdzenie \ref{it:thm:kh-primitive-derivative}),
lecz żadna nie może mieć skończonej całki Lebesgue'a,
gdyż wówczas należałaby do HK. Stąd obie całki
Lebesgue'a są nieskończone. Całka Lebesgue'a ze
znakiem nie jest wtedy określona, mimo istnienia
skończonej całki HK.
\end{proof}


\section{Całka Denjoya w wersji wąskiej}
Absolutna ciągłość jest zbyt silna, aby obejmować funkcje pierwotne wszystkich pochodnych. Wprowadzamy warunek kontrolujący oscylacje na zbiorach, a następnie jego przeliczalną wersję.
\begin{definicja}
Niech $F:[a,b]\to\R$ będzie ciągła i $E\subset[a,b]$. Piszemy $F\in\AC_*(E)$, jeżeli dla każdego $\varepsilon>0$ istnieje $\eta>0$ takie, że dla każdej skończonej rodziny przedziałów $[u_j,v_j]$ o parami rozłącznych wnętrzach i końcach $u_j,v_j\in E$ spełniającej $\sum_j(v_j-u_j)<\eta$ mamy
\[
\sum_j\osc_{[u_j,v_j]}F<\varepsilon.
\]
Piszemy $F\in\AC G_*[a,b]$, jeżeli $[a,b]=\bigcup_{n\ge1}E_n$ dla domkniętych zbiorów $E_n$ takich, że $F\in\AC_*(E_n)$. Na całym przedziale $\AC_*([a,b])=\AC[a,b]$, ale na podzbiorze $E$ warunek oscylacyjny jest silniejszy od kontroli samych przyrostów między punktami $E$.
\end{definicja}
\begin{definicja}[Wąska całka Denjoya]
Funkcja $f:[a,b]\to\R$ o skończonych wartościach jest całkowalna w wąskim sensie Denjoya, jeśli istnieje ciągła $F\in\AC G_*[a,b]$ taka, że $F'=f$ prawie wszędzie. Wówczas definiujemy $(\mathrm D_*)\int_a^b f=F(b)-F(a)$. Jednoznaczność uzasadnimy poniżej.
\end{definicja}
\begin{twierdzenie}[Każda różniczkowalna funkcja należy do $\AC G_*$]
Jeśli $F:[a,b]\to\R$ jest ciągła i ma skończoną pochodną w każdym punkcie wnętrza oraz odpowiednie pochodne jednostronne na końcach, to $F\in\AC G_*[a,b]$.
\end{twierdzenie}
\begin{proof}
Dla $n\in\N$ niech $E_n$ będzie zbiorem tych $x$, dla których $|F(y)-F(x)|\le n|y-x|$ dla wszystkich $y\in[a,b]$ z $|y-x|\le1/n$. Różniczkowalność daje $\bigcup_nE_n=[a,b]$. Każdy $E_n$ jest domknięty: przy $x_k\in E_n$ i $x_k\to x$ przechodzimy do granicy w nierówności, wykorzystując ciągłość $F$. Jeśli $u,v\in E_n$ i $v-u<1/n$, to dla każdego $y\in[u,v]$ mamy $|F(y)-F(u)|\le n(v-u)$. Stąd $\osc_{[u,v]}F\le2n(v-u)$. Dla rodziny przedziałów z końcami w $E_n$ i $\sum_j(v_j-u_j)<\min\{1/n,\varepsilon/(2n)\}$ suma oscylacji jest mniejsza od $\varepsilon$.
\end{proof}
\subsection*{Kontrola oscylacji na zbiorach zerowych}
\begin{lemat}[Zaniedbywalne wahanie na zbiorze zerowym]
\label{it:lem:acstar-null-variation}
Niech $F$ będzie ciągła na $[a,b]$, $E\subset[a,b]$ domknięty i $F\in\AC_*(E)$. Jeżeli $N\subset E$ ma miarę Lebesgue'a zero, to dla każdego $\varepsilon>0$ istnieje cechowanie na $N$, dla którego dowolny skończony drobny podział częściowy $\{(t_j,[u_j,v_j])\}$ ze znacznikami w $N$ spełnia
\[
 \sum_j|F(v_j)-F(u_j)|<\varepsilon.
\]
\end{lemat}
\begin{proof}
Dla pustego $E$ teza jest bezprzedmiotowa. Składowe przedziałowe dopełnienia $E$ w $[a,b]$ są co najwyżej przeliczalne (każda niepusta składowa zawiera liczbę wymierną). Niech $J_k$ oznaczają ich domknięcia. Co najwyżej dwie składowe brzegowe nie mają obu końców w $E$. Suma długości wszystkich $J_k$ nie przekracza $b-a$, ponieważ ich wnętrza są rozłączne.

Twierdzimy, że $\sum_k\osc_{J_k}F<\infty$. Wybieramy $\eta_0>0$ z definicji $\AC_*(E)$ dla błędu $1$. Pomijając skończenie wiele składowych, w tym składowe brzegowe, możemy zapewnić, że suma długości pozostałych jest mniejsza niż $\eta_0$. Każda skończona podrodzina pozostałych ma końce w $E$ i spełnia warunek definicji, więc suma jej oscylacji jest mniejsza od $1$. Supremum po podrodzinach daje zbieżność sumy; skończenie wiele pominiętych oscylacji jest skończonych dzięki ciągłości $F$.

Wybieramy skończoną rodzinę $\mathcal J_0$ składowych (w tym brzegowe) tak, by $\sum_{J_k\notin\mathcal J_0}\osc_{J_k}F<\varepsilon/8$. Niech $m$ oznacza jej liczność. Z definicji $\AC_*(E)$ wybieramy $\eta>0$ dla błędu $\varepsilon/4$. Ponieważ $N$ ma miarę zero, istnieje otwarty zbiór $O\supset N$ o mierze mniejszej od $\eta$. Dla $t\in N$ dobieramy $\delta(t)>0$ tak, by jego otoczenie o promieniu $\delta(t)$ leżało w $O$. Dla każdej składowej $J\in\mathcal J_0$, która nie zawiera $t$, dodatkowo wymagamy $\delta(t)<\operatorname{dist}(t,J)$. Jeżeli $t\in J$, jest jego końcem, gdyż $t\in E$. Z ciągłości $F$ zmniejszamy $\delta(t)$ tak, by $|F(y)-F(t)|<\varepsilon/(8(m+1))$ dla wszystkich $y\in[a,b]$ z $|y-t|<\delta(t)$. Warunki są skończone i wszystkie dopuszczają dodatni promień.

Dla każdej komórki $I_j=[u_j,v_j]$ z rozpatrywanego podziału częściowego zbiór $E\cap I_j$ jest niepusty (zawiera $t_j$) i zwarty. Niech $p_j=\min(E\cap I_j)$, $q_j=\max(E\cap I_j)$. Rozkładamy przyrost:
\[
 |F(v_j)-F(u_j)|
 \le |F(q_j)-F(p_j)|+|F(p_j)-F(u_j)|+|F(v_j)-F(q_j)|.
\]
Przedziały $[p_j,q_j]$ mają rozłączne wnętrza, końce w $E$ i sumę długości nie większą od $\lambda(O)<\eta$. Suma pierwszych składników jest więc mniejsza od $\varepsilon/4$. Jeśli któryś z pozostałych odcinków ma dodatnią długość, leży w domknięciu jednej składowej dopełnienia $E$ i przylega do jej końca w $E$. Dla ustalonej składowej mogą wystąpić najwyżej dwa takie odcinki o dodatniej długości, po jednym przy każdym końcu: dwa przy tym samym końcu miałyby wspólne wnętrze. Składowe brzegowe mają co najwyżej jeden taki odcinek.

Dla składowych spoza $\mathcal J_0$ suma tych składników nie przekracza $2\sum_{J_k\notin\mathcal J_0}\osc_{J_k}F<\varepsilon/4$. Jeżeli komórka wchodzi w składową $J\in\mathcal J_0$, to jej znacznik musi należeć do $J$ (w przeciwnym przypadku warunek odległości wyklucza przecięcie). Znacznik jest więc końcem $J$ w $E$ i odcinek skrajny przy tym końcu ma przyrost mniejszy od $\varepsilon/(8(m+1))$. Łącznie składniki z $\mathcal J_0$ wnoszą mniej niż $\varepsilon/4$. Dodanie trzech oszacowań kończy dowód.
\end{proof}

\begin{wniosek}
\label{it:cor:acgstar-null-variation}
Jeśli $F\in\AC G_*[a,b]$ i $N\subset[a,b]$ ma miarę zero, to wniosek lematu \ref{it:lem:acstar-null-variation} zachodzi dla całego $N$.
\end{wniosek}
\begin{proof}
Niech $[a,b]=\bigcup_{n\ge1}E_n$ będzie domkniętym pokryciem z definicji. Dla $N\cap E_n$ stosujemy lemat z dokładnością $\varepsilon2^{-n-1}$, otrzymując cechowanie $\delta_n$. Każdemu $t\in N$ przypisujemy najmniejsze $n(t)$ z $t\in E_{n(t)}$ i przyjmujemy $\delta(t)=\delta_{n(t)}(t)$. W dowolnym skończonym podziale częściowym grupujemy komórki według $n(t_j)$. Każda grupa spełnia hipotezy odpowiedniego lematu, więc suma modułów przyrostów jest mniejsza od $\sum_{n\ge1}\varepsilon2^{-n-1}=\varepsilon/2$ (dla pustej rodziny suma jest zerowa). W szczególności jest mniejsza od $\varepsilon$.
\end{proof}

\begin{twierdzenie}[Całkowalność pochodnej prawie wszędzie funkcji $\AC G_*$]
\label{it:thm:acgstar-derivative-integrable}
Niech $F\in\AC G_*[a,b]$ oraz niech $f:[a,b]\to\R$ ma skończone wartości i spełnia $F'=f$ prawie wszędzie. Wtedy $f\in\HK[a,b]$ i dla każdego $x\in[a,b]$
\[
 F(x)-F(a)=(\HK)\int_a^x f(t)\,\mathrm dt.
\]
\end{twierdzenie}
\begin{proof}
Zakładamy $a<b$; przypadek przedziału jednopunktowego daje całkę zero. Niech $N$ będzie zbiorem miary zero, poza którym istnieje skończona pochodna równa $f$; dołączamy do $N$ oba końce przedziału. Dla $\varepsilon>0$ wybieramy na $N$ cechowanie z wniosku \ref{it:cor:acgstar-null-variation} dla błędu $\varepsilon/3$. Funkcja $h=f\chi_N$ jest równa zeru poza $N$ i ma skończone wartości. Lemat o zmianach na zbiorze miary zero z rozdziału o całce KH, zastosowany do $|h|$, daje cechowanie dla błędu $\varepsilon/3$. W punktach $N$ bierzemy minimum obu cechowań.

Dla $t\notin N$ z definicji pochodnej dobieramy $\delta(t)>0$, nie większe od cechowania dla $|h|$, aby
\[
 |F(y)-F(t)-f(t)(y-t)|\le
 \frac{\varepsilon}{3(b-a)}|y-t|
 \qquad (|y-t|<\delta(t),\ y\in[a,b]).
\]
Niech $\mathcal P$ będzie dowolnym pełnym podziałem drobnym. Dla komórek ze znacznikami poza $N$ odejmujemy oszacowania dla $y=u_j$ i $y=v_j$; nierówność trójkąta daje błąd nie większy od $\varepsilon(v_j-u_j)/(3(b-a))$. Suma tych błędów nie przekracza $\varepsilon/3$. Dla komórek ze znacznikami w $N$ suma modułów przyrostów $F$ jest mniejsza od $\varepsilon/3$, a suma $|f(t_j)|(v_j-u_j)$ jest mniejsza od $\varepsilon/3$: pełna suma dla $|h|$ jest właśnie tą drugą sumą i ma błąd mniejszy od $\varepsilon/3$ względem całki zero. W konsekwencji
\[
 |S(f,\mathcal P)-F(b)+F(a)|<\varepsilon,
\]
ponieważ suma wszystkich przyrostów $F$ jest teleskopowa. Definicja daje całkowalność i wartość na $[a,b]$. Ten sam argument stosujemy na każdym $[a,x]$: przecięcia zbiorów $E_n$ z podprzedziałem są domkniętym pokryciem, warunek $\AC_*$ pozostaje prawdziwy, a zbiór wyjątkowy pozostaje zerowy po dołączeniu końców. To daje wzór dla wszystkich $x$.
\end{proof}

\subsection*{Różniczkowalność prawie wszędzie}
\begin{lemat}[Przedłużenie afiniczne przez luki]
\label{it:lem:acstar-affine-extension}
Niech $E\subset[a,b]$ będzie niepustym zbiorem domkniętym, a $F:[a,b]\to\R$ funkcją ciągłą należącą do $\AC_*(E)$. Istnieje funkcja $G\in\AC[a,b]$, równa $F$ na $E$, afiniczna na każdej wewnętrznej składowej dopełnienia $E$ i stała na każdej składowej brzegowej dopełnienia.
\end{lemat}
\begin{proof}
Niech $(l_k,r_k)$ będą wewnętrznymi składowymi dopełnienia $E$; obydwa ich końce należą do $E$. Na $E$ przyjmujemy $G=F$, na $[l_k,r_k]$ kładziemy
\[
 G(x)=F(l_k)+\frac{x-l_k}{r_k-l_k}(F(r_k)-F(l_k)).
\]
Na $[a,\min E]$ i $[\max E,b]$ przyjmujemy wartości stałe, odpowiednio $F(\min E)$ i $F(\max E)$. Definicje są zgodne na końcach. Funkcja $G$ jest ciągła: we wnętrzu luk jest afiniczna lub stała, a dla $x\in E$ rozważamy $y\to x$. Jeżeli $y\in E$, korzystamy z ciągłości $F$. Jeżeli $y$ należy do wewnętrznej luki długości mniejszej od $\eta$, jej oba końce są oddalone od $x$ o mniej niż $|x-y|+\eta$, a $G(y)$ jest kombinacją wypukłą ich wartości. Jednostajna ciągłość $F$ pozwala uczynić różnicę $|G(y)-F(x)|$ dowolnie małą przez wybór $\eta$ i $|x-y|$. Luk długości co najmniej $\eta$ jest skończenie wiele; przy $y\to x$ możemy spośród nich rozpatrywać tylko luki, których domknięcia zawierają $x$, i na nich ciągłość w końcu wynika bezpośrednio ze wzoru afinicznego. To samo dotyczy co najwyżej dwóch luk brzegowych.

Z dowodu lematu \ref{it:lem:acstar-null-variation} wiemy, że $\sum_k\osc_{[l_k,r_k]}F<\infty$. W szczególności funkcja
\[
 w(t)=\sum_k\frac{|F(r_k)-F(l_k)|}{r_k-l_k}
             \chi_{(l_k,r_k)}(t)
\]
jest nieujemna, mierzalna i całkowalna w sensie Lebesgue'a: przez zbieżność monotoniczną jej całka jest sumą $|F(r_k)-F(l_k)|$, ograniczoną przez sumę oscylacji. Jeśli $[u,v]$ leży w domknięciu jednej luki, to $|G(v)-G(u)|=\int_u^v w\,\mathrm d\lambda$; w luce brzegowej obie strony są zerowe.

Dla dowolnego $[u,v]$ przecinającego $E$ oznaczamy $p=\min(E\cap[u,v])$, $q=\max(E\cap[u,v])$. Przedziały skrajne $[u,p]$ i $[q,v]$ leżą w lukach lub są jednopunktowe. Zatem
\[
 |G(v)-G(u)|\le |F(q)-F(p)|
       +\int_{[u,p]\cup[q,v]} w\,\mathrm d\lambda.
\]
Rozważmy teraz skończoną rodzinę przedziałów o rozłącznych wnętrzach. Dla tych przecinających $E$ środkowe przedziały $[p,q]$ także mają rozłączne wnętrza i końce w $E$. Przy dostatecznie małej łącznej długości całej rodziny suma $|F(q)-F(p)|$ jest mniejsza od $\varepsilon/2$ przez $\AC_*(E)$. Wszystkie odcinki skrajne i przedziały leżące w lukach mają rozłączne wnętrza, a ich suma długości nie przekracza sumy długości całej rodziny. Absolutna ciągłość całki Lebesgue'a funkcji $w\in L^1$ (rozdział o całce Lebesgue'a) daje dla dostatecznie małej tej sumy łączny wkład całek mniejszy od $\varepsilon/2$. Otrzymujemy warunek absolutnej ciągłości $G$.
\end{proof}

\begin{twierdzenie}[Różniczkowalność prawie wszędzie na zbiorach $\AC_*$]
\label{it:thm:acstar-ae-differentiable}
Jeśli $F$ jest ciągła na $[a,b]$ i $F\in\AC_*(E)$ dla domkniętego $E\subset[a,b]$, to $F$ ma skończoną pochodną prawie wszędzie na $E\cap(a,b)$. W szczególności każda funkcja $F\in\AC G_*[a,b]$ jest różniczkowalna prawie wszędzie.
\end{twierdzenie}
\begin{proof}
Dla pustego $E$ nie ma czego dowodzić. Niech $G$ będzie przedłużeniem z lematu \ref{it:lem:acstar-affine-extension} i $H=F-G$. Na $E$ mamy $H=0$. Ponieważ $G$ jest absolutnie ciągła, jest różniczkowalna prawie wszędzie z podstawowego twierdzenia rachunku całkowego dla całki Lebesgue'a (skrypt \emph{Teoria miary}, rozdział 25, twierdzenie o etykiecie \texttt{thm:FTC-lebesgue-25}). Wystarczy wykazać, że $H'=0$ prawie wszędzie na $E$.

Dla wewnętrznej luki $(l_k,r_k)$ połóżmy $d_k=\osc_{[l_k,r_k]}F$. Z konstrukcji $G$ obie wartości $F(y),G(y)$ leżą pomiędzy minimum i maksimum $F$ na $[l_k,r_k]$. Wobec tego $|H(y)|\le d_k$ dla $y\in[l_k,r_k]$. Ponadto $\sum_k d_k<\infty$ z dowodu lematu \ref{it:lem:acstar-null-variation}.

Dla dodatniej liczby całkowitej $m$ rozważmy zbiory
\[
 U_{k,m}=(l_k-md_k,l_k+md_k)\cup(r_k-md_k,r_k+md_k).
\]
Ich miary nie przekraczają $4md_k$. Jeżeli luk jest nieskończenie wiele, zbiór
\[
 B_m=\bigcap_{n\ge1}\bigcup_{k\ge n}U_{k,m}
\]
ma miarę zero, gdyż dla każdego $n$ jego miara nie przekracza $4m\sum_{k\ge n}d_k$, a prawa strona dąży do zera. Przy skończonej liczbie luk przyjmujemy $B_m=\varnothing$. Wykluczamy z $E$ sumę zbiorów $B_m$, wszystkie końce luk (także brzegowych), oraz $a,b$. Jest to zbiór miary zero.

Ustalmy pozostały punkt $x\in E$ i $m\ge1$. Dla wszystkich dostatecznie dużych $k$ punkt $x$ nie należy do $U_{k,m}$. Jeśli $y$ leży w takiej luce, to jeden z jej końców leży pomiędzy $x$ i $y$, ponieważ $x\in E$. Oznaczając go przez $z$, otrzymujemy $|x-y|\ge |x-z|\ge md_k$. Dla $d_k>0$ mamy więc $|H(y)|/|y-x|\le1/m$; dla $d_k=0$ iloraz jest zerowy. Skończenie wiele pozostałych luk, w tym brzegowe, ma domknięcia niezawierające $x$, zatem można wybrać otoczenie $x$, które nie przecina żadnej z nich. Dla $y\in E$ i $y\ne x$ iloraz jest zerowy. W konsekwencji
\[
 \limsup_{\substack{y\to x\\y\ne x}}
 \frac{|H(y)-H(x)|}{|y-x|}\le\frac1m.
\]
Dowolność $m$ daje $H'(x)=0$. Poza zbiorem zerowym, na którym $G'$ nie istnieje, mamy więc $F'=G'$. Dla funkcji $\AC G_*$ stosujemy wynik do każdego domkniętego zbioru $E_n$ z definicji. Przeliczalna suma zbiorów wyjątkowych jest zbiorem miary zero, a zbiory $E_n$ pokrywają przedział.
\end{proof}

\begin{twierdzenie}[Charakteryzacja funkcji pierwotnych HK; wynik zewnętrzny]
Dla ciągłej $F:[a,b]\to\R$ równoważne są: (i) $F\in\AC G_*[a,b]$; (ii) istnieje $h\in\HK[a,b]$ takie, że $F(x)=F(a)+\int_a^x h$ dla każdego $x$. Wówczas $F'=h$ prawie wszędzie.
\end{twierdzenie}
Twierdzenie w pełnej postaci pozostaje na tym etapie wynikiem zewnętrznym. Twierdzenie \ref{it:thm:acgstar-derivative-integrable} dowodzi już kierunku od funkcji klasy $\AC G_*$ ze znaną pochodną prawie wszędzie do pierwotnej KH. Istnienie pochodnej prawie wszędzie każdej funkcji tej klasy wynika teraz z twierdzenia \ref{it:thm:acstar-ae-differentiable}. Pozostaje udowodnić kierunek od pierwotnej KH do $\AC G_*$; sama ciągłość i pochodna prawie wszędzie nie wystarczają.
\begin{wniosek}[Denjoy = HK]
Wąska całka Denjoya i całka HK mają te same funkcje całkowalne oraz te same wartości. Jeśli $F\in\AC G_*[a,b]$ i $F'=0$ prawie wszędzie, to $F$ jest stała.
\end{wniosek}
\begin{proof}
Dla funkcji HK-całkowalnej jej pierwotna należy do $\AC G_*$ na mocy poprzedniego twierdzenia i ma pochodną równą funkcji podcałkowej prawie wszędzie. Odwrotnie, jeśli $F\in\AC G_*$ oraz $F'=f$ prawie wszędzie, to $F-F(a)=\int_a^x h$ dla pewnego $h\in\HK$, z $h=f$ prawie wszędzie. Lemat o zmianie na zbiorze zerowym pozwala zastąpić $h$ przez $f$. Dla $f=0$ pierwotna jest stała, co zapewnia jednoznaczność. Kombinacje liniowe takich funkcji pierwotnych pozostają w $\AC G_*$.
\end{proof}
\begin{przyklad}[Całkowalność z definicji Denjoya]
Dla $F(x)=x^2\cos(\pi/x^2)$, $x>0$, $F(0)=0$, bierzemy domknięte pokrycie $\{0\},[1/2,1],[1/3,1],\dots$ przedziału $[0,1]$. Na każdym $[1/n,1]$ funkcja $F$ jest klasy $C^1$, więc absolutnie ciągła, a na zbiorze jednopunktowym warunek $\AC_*$ jest automatyczny. Zatem $F\in\AC G_*$ i $F'$ jest całkowalna w sensie Denjoya z wartością $F(1)-F(0)=-1$. Pochodna nie jest jednak absolutnie całkowalna, co wynika z rozbieżności całki modułu jej oscylacyjnego składnika.
\end{przyklad}
\begin{uwaga}[Wąska i szeroka całka Denjoya]
Wąska całka Denjoya ma tę samą klasę funkcji co HK. Szeroką całkę Denjoya opiszemy przez klasę $\AC G$ i pochodną aproksymatywną; jest to istotnie szersza teoria. Historyczna totalizacja Denjoya nie musi odpowiadać w każdym szczególe poniższej definicji opisowej.
\end{uwaga}

% ============================================================
% UZUPELNIENIE MATEMATYCZNE: ROZDZIAL 8
% ============================================================
\subsection*{Jak w praktyce rozpoznawać pierwotne wąskiej całki Denjoya}
\begin{twierdzenie}[Przeliczalny zbiór punktów nieróżniczkowalności]\label{tw:denjoy-countable}
Niech $F:[a,b]\to\R$ będzie ciągła oraz różniczkowalna ze skończoną pochodną w każdym punkcie $(a,b)\setminus N$, gdzie $N\subset[a,b]$ jest zbiorem co najwyżej przeliczalnym. Wtedy $F\in\AC G_*[a,b]$. Jeśli funkcja $f:[a,b]\to\R$ ma skończone wartości i $f=F'$ na $(a,b)\setminus N$, to $f$ jest całkowalna w sensie HK oraz w wąskim sensie Denjoya, a obie całki są równe $F(b)-F(a)$.
\end{twierdzenie}
\begin{proof}
Dla $n\ge1$ definiujemy
\[
 E_n=\left\{x\in[a,b]: |F(y)-F(x)|\le n|y-x|
 \ \text{dla wszystkich}\ y\in[a,b]\ \text{z}\ |y-x|\le1/n\right\}.
\]
Z ciągłości $F$ wynika, że każdy $E_n$ jest domknięty. Każdy punkt, w którym istnieje skończona pochodna (jednostronna na brzegu), należy do któregoś $E_n$. Dla przedziału $[u,v]$ o końcach w $E_n$ i długości mniejszej niż $1/n$ mamy
$\osc_{[u,v]}F\le2n(v-u)$, jak w dowodzie wcześniejszego twierdzenia. Stąd $F\in\AC_*(E_n)$. Zbiory jednopunktowe $\{x\}$, $x\in N\cup\{a,b\}$, są domknięte, a warunek $\AC_*$ zachodzi na nich trywialnie. Przeliczalna rodzina $E_n$ i tych singletonów pokrywa $[a,b]$, więc $F\in\AC G_*$. Ponieważ $N$ ma miarę Lebesgue'a zero, $F'=f$ prawie wszędzie; stosujemy charakteryzację pierwotnych HK i niezmienniczość całki przy zmianie na zbiorze zerowym.
\end{proof}
\begin{przyklad}[Funkcja o przeliczalnie wielu osobliwościach]
Niech $c_n=2^{-n}$, $r_n=2^{-n-3}$ oraz niech $\psi\in C_c^1((-1,1))$ spełnia $\psi(0)=1$. Definiujemy
$F(0)=0$ oraz $F(x)=\sum_{n\ge1}n^{-1}\psi((x-c_n)/r_n)$ dla $x\in(0,1]$. Nośniki składników są parami rozłączne, a amplitudy $1/n$ maleją do zera, więc $F$ jest ciągła na $[0,1]$ i klasy $C^1$ poza punktem $0$. Jej pochodna jest HK-całkowalna oraz $\int_0^1F'=F(1)-F(0)=0$. Ponieważ $\psi$ jest niezerowa i znika na brzegu nośnika, suma wariacji na rozłącznych nośnikach jest nie mniejsza niż $2\sum_{n\ge1}1/n=\infty$. W istocie $\int_0^1|F'(x)|\,\mathrm d\lambda(x)=\left(\int_{-1}^1|\psi'(t)|\,\mathrm dt\right)\sum_{n\ge1}1/n=\infty$. Stąd $F\notin\BV[0,1]$ i $F\notin\AC[0,1]$. Mimo to całka HK pochodnej istnieje.
\end{przyklad}
\begin{stwierdzenie}[Sklejanie pierwotnych]\label{st:denjoy-sklejanie}
Niech $a<c<b$ i niech $F:[a,b]\to\R$ będzie ciągła. Jeśli $F|_{[a,c]}\in\AC G_*[a,c]$ oraz $F|_{[c,b]}\in\AC G_*[c,b]$, to $F\in\AC G_*[a,b]$. Jeśli dodatkowo $F'=f$ prawie wszędzie na obu podprzedziałach, to $(\mathrm D_*)\int_a^b f=(\mathrm D_*)\int_a^c f+(\mathrm D_*)\int_c^b f$.
\end{stwierdzenie}
\begin{proof}
Domknięte pokrycia obu przedziałów przez zbiory, na których obowiązuje $\AC_*$, są także pokryciami przez zbiory domknięte w $[a,b]$. Łącząc te dwie rodziny, otrzymujemy przeliczalne pokrycie $[a,b]$ bez konieczności kontrolowania przedziałów mających końce po różnych stronach $c$: definicja $\AC G_*$ wymaga kontroli osobno na każdym zbiorze pokrycia. Wzór całkowy jest tożsamością przyrostów $F$.
\end{proof}
\begin{uwaga}[Pochodna prawie wszędzie nie identyfikuje pierwotnej]
Ciągłość i równość $F'=f$ prawie wszędzie nie wystarczają do wzoru Newtona--Leibniza. Przykładem jest funkcja Cantora $C$, dla której $C'=0$ prawie wszędzie, lecz $C(1)-C(0)=1$. Warunek $\AC G_*$ usuwa tę niejednoznaczność. Jest zatem warunkiem na \emph{pierwotną}, a nie tylko na zbiór punktów istnienia jej pochodnej.
\end{uwaga}


\section{Szeroka całka Denjoya (Denjoya--Chinczyna)}
Zmieniamy dwa elementy definicji: oscylacje zastępujemy przyrostami między końcami przedziałów, a zwykłą pochodną --- pochodną aproksymatywną. Zachowujemy ciągłość funkcji pierwotnej.
\begin{definicja}[Pochodna aproksymatywna]
Mierzalny zbiór $E\subset\R$ ma gęstość jeden w punkcie $x$, jeśli $\lambda_1(E\cap(x-r,x+r))/(2r)\to1$ przy $r\downarrow0$. Gęstość zero definiujemy analogicznie, zastępując granicę przez zero. Pochodną aproksymatywną rozpatrujemy w punktach wnętrza przedziału; wartości na końcach nie wpływają na całkę. Mówimy, że $F$ ma pochodną aproksymatywną $F'_{\mathrm{ap}}(x)=L$, jeżeli dla każdego $\varepsilon>0$ zbiór punktów $y\ne x$ spełniających
$\left|\frac{F(y)-F(x)}{y-x}-L\right|\ge\varepsilon$ ma gęstość zero w $x$. Zwykła pochodna, o ile istnieje, jest też pochodną aproksymatywną o tej samej wartości.
\end{definicja}
\begin{lemat}[Jednoznaczność i liniowość pochodnej aproksymatywnej]\label{it:lem:approx-derivative}
Skończona pochodna aproksymatywna w punkcie jest wyznaczona jednoznacznie. Jeśli dwie funkcje mają w tym punkcie skończone pochodne aproksymatywne, to ich kombinacja liniowa ma pochodną aproksymatywną równą odpowiedniej kombinacji.
\end{lemat}
\begin{proof}
Gdyby dwie różne liczby $L_1,L_2$ były pochodnymi aproksymatywnymi tej samej funkcji, bierzemy $\varepsilon=|L_1-L_2|/3$. Zbiory punktów, w których iloraz różnicowy różni się od $L_i$ o mniej niż $\varepsilon$, mają gęstość jeden. Ich przecięcie również ma gęstość jeden, ponieważ miara dopełnienia przecięcia jest ograniczona przez sumę miar dopełnień. Przecięcie nie jest więc puste w dostatecznie małych otoczeniach punktu, ale nierówność trójkąta dawałaby $|L_1-L_2|<2\varepsilon$, sprzeczność.
Dla kombinacji $\alpha F+\beta G$ wybieramy dla obu funkcji zbiory dobrych ilorazów z dokładnością $\varepsilon/(1+|\alpha|+|\beta|)$. Na ich przecięciu błąd ilorazu kombinacji jest mniejszy od $\varepsilon$, a przecięcie ma gęstość jeden. To daje tezę.
\end{proof}

\begin{definicja}[Klasa $\AC G$ i szeroka całka Denjoya]
Dla ciągłej $F:[a,b]\to\R$ i $E\subset[a,b]$ piszemy $F\in\AC(E)$, jeśli dla każdego $\varepsilon>0$ istnieje $\eta>0$ takie, że dla każdej skończonej rodziny przedziałów $[u_j,v_j]$ o rozłącznych wnętrzach i końcach należących do $E$ warunek $\sum_j(v_j-u_j)<\eta$ pociąga $\sum_j|F(v_j)-F(u_j)|<\varepsilon$. Piszemy $F\in\AC G[a,b]$, gdy $[a,b]=\bigcup_{n\ge1}E_n$ dla zbiorów domkniętych $E_n$ z $F\in\AC(E_n)$. Funkcja $f$ jest całkowalna w szerokim sensie Denjoya, jeżeli istnieje ciągła $F\in\AC G[a,b]$ z $F'_{\mathrm{ap}}=f$ prawie wszędzie. Definiujemy $ (\mathrm D_{\mathrm{sz}})\int_a^b f=F(b)-F(a)$.
\end{definicja}
\begin{lemat}[Przedłużenie ze zbioru domkniętego]\label{it:lem:ac-affine-extension}
Jeśli $E\subset[a,b]$ jest domknięty, $F:[a,b]\to\R$ jest ciągła i $F\in\AC(E)$, to $F|_E$ ma absolutnie ciągłe przedłużenie $u:[a,b]\to\R$. Ponadto dla każdego zbioru borelowskiego $B\subset E$ zachodzi
\[
 \lambda_1^*(F(B))\le\int_B|u'(x)|\,\mathrm d\lambda_1(x).
\]
W szczególności dla $B=E$ obraz jest zwarty, więc można użyć miary zamiast miary zewnętrznej.
\end{lemat}
\begin{proof}
Dla pustego lub jednopunktowego $E$ bierzemy przedłużenie stałe. W pozostałych przypadkach oznaczamy luki wewnętrzne przez $(l_k,r_k)$ i przedłużamy $F$ afinicznie na ich domknięciach, a poza $[\min E,\max E]$ stale. Ciągłość tego przedłużenia wynika z argumentu o krótkich i długich lukach z lematu \ref{it:lem:acstar-affine-extension}; ten argument nie wymaga kontroli oscylacji.

Najpierw wykażemy, że $\sum_k|F(r_k)-F(l_k)|<\infty$. Wybieramy $\eta_0$ z warunku $\AC(E)$ dla błędu $1$. Ponieważ suma długości luk nie przekracza $b-a$, po odrzuceniu skończenie wielu luk suma pozostałych długości jest mniejsza od $\eta_0$. Dla każdej skończonej podrodziny ogona warunek $\AC(E)$ ogranicza sumę modułów przyrostów przez $1$. Supremum po podrodzinach daje zbieżność szeregu.
Stąd funkcja
\[
 w(x)=\sum_k\frac{|F(r_k)-F(l_k)|}{r_k-l_k}\chi_{(l_k,r_k)}(x)
\]
należy do $\mathcal L^1$: jest nieujemna, mierzalna, a jej całka na mocy zbieżności monotonicznej jest powyższym szeregiem.

Dla przedziału $[s,t]$ przecinającego $E$ połóżmy
$p=\min(E\cap[s,t])$ i $q=\max(E\cap[s,t])$. Rozpisując przyrost przez $p,q$, otrzymujemy
\[
 |u(t)-u(s)|\le |F(q)-F(p)|+
 \int_{[s,p]\cup[q,t]}w\,\mathrm d\lambda_1.
\]
Gdy przedział nie przecina $E$, moduł przyrostu jest równy całce $w$ po nim (również w luce brzegowej, gdzie oba są zerowe).
Dla skończonej rodziny przedziałów o rozłącznych wnętrzach środkowe odcinki $[p,q]$ mają końce w $E$, rozłączne wnętrza i sumę długości nie większą od długości całej rodziny. Przy dostatecznie małej tej sumie ich przyrosty są ograniczone przez $\varepsilon/2$ z warunku $\AC(E)$. Wszystkie pozostałe odcinki mają rozłączne wnętrza i łączną długość nie większą od długości rodziny. Absolutna ciągłość całki $w$ daje dla dostatecznie małej długości ich łączny wkład mniejszy od $\varepsilon/2$. Dowodzi to $u\in\AC[a,b]$.

Z twierdzenia \ref{it:thm:ac-lebesgue-primitive} wynika $\operatorname{diam}u(J)\le\int_J|u'|\,\mathrm d\lambda_1$ dla każdego przedziału $J\subset[a,b]$ (także względnie otwartego, przez supremum po jego punktach).
Dla $\varepsilon>0$ z regularności miary i absolutnej ciągłości całki wybieramy względnie otwarty $O\supset B$ z $\int_{O\setminus B}|u'|\,\mathrm d\lambda_1<\varepsilon$. Rozkładamy $O$ na co najwyżej przeliczalnie wiele rozłącznych przedziałów $J_k$. Obrazy $u(J_k)$ są przedziałami z ciągłości $u$, a ich długości nie przekraczają całek $|u'|$ po $J_k$. Ponieważ $F(B)=u(B)\subset\bigcup_k u(J_k)$, podaddytywność miary zewnętrznej daje
\[
 \lambda_1^*(F(B))\le\sum_k\int_{J_k}|u'|\,\mathrm d\lambda_1
 =\int_O|u'|\,\mathrm d\lambda_1
 <\int_B|u'|\,\mathrm d\lambda_1+\varepsilon.
\]
Dowolność $\varepsilon$ kończy dowód.
\end{proof}
\begin{twierdzenie}[Pochodna aproksymatywna i jednoznaczność]\label{it:thm:acg-uniqueness}
Każda ciągła funkcja $F\in\AC G[a,b]$ ma skończoną pochodną aproksymatywną prawie wszędzie. Jeśli $F'_{\mathrm{ap}}=0$ prawie wszędzie, to $F$ jest stała. W konsekwencji szeroka całka Denjoya jest dobrze określona i liniowa.
\end{twierdzenie}
\begin{proof}
Niech $[a,b]=\bigcup_n E_n$ będzie domkniętym pokryciem z definicji, a $u_n$ --- absolutnie ciągłym przedłużeniem $F|_{E_n}$ z lematu \ref{it:lem:ac-affine-extension}. Prawie każdy punkt $x\in E_n$ jest punktem gęstości jeden zbioru $E_n$ i punktem różniczkowalności $u_n$. W takim punkcie, dla $y\in E_n$, ilorazy różnicowe $F$ są równe ilorazom dla $u_n$, więc $F'_{\mathrm{ap}}(x)=u_n'(x)$. Wyjątki ze wszystkich $n$ tworzą zbiór miary zero. Pochodna aproksymatywna ma mierzalnego reprezentanta: na zbiorze $E_n\setminus\bigcup_{j<n}E_j$ bierzemy mierzalną wersję $u_n'$ z zerem na zbiorze nieróżniczkowalności. Te zbiory są borelowskie i tworzą rozłączne pokrycie. Otrzymana funkcja zgadza się z pochodną aproksymatywną prawie wszędzie, zatem każda jej wersja o skończonych wartościach na zbiorze wyjątkowym jest mierzalna w sensie Lebesgue'a.

Jeżeli $F'_{\mathrm{ap}}=0$ prawie wszędzie, to $u_n'=0$ prawie wszędzie na $E_n$. Z lematu wynika $\lambda(F(E_n))\le\int_{E_n}|u_n'|\,\mathrm d\lambda=0$, a więc $\lambda(F([a,b]))=0$. Ponieważ ciągły obraz przedziału jest przedziałem, $F([a,b])$ musi być jednym punktem i $F$ jest stała.

Klasa $\AC G$ jest stabilna na kombinacje liniowe: wystarczy przejść do przeliczalnego pokrycia przecięciami zbiorów domkniętych. Liniowość pochodnych aproksymatywnych wynika z lematu \ref{it:lem:approx-derivative}. Dwie pierwotne tej samej funkcji różnią się zatem stałą, a całka jest jednoznaczna i liniowa.
\end{proof}
\begin{wniosek}[Właściwe rozszerzenie klasy HK]
Każda funkcja całkowalna w sensie HK, Perrona lub wąskim sensie Denjoya jest całkowalna w szerokim sensie Denjoya z tą samą wartością całki. Inkluzja jest właściwa.
\end{wniosek}
\begin{proof}
Dla dowolnego $E$ nierówność $|F(v)-F(u)|\le\osc_{[u,v]}F$ daje $\AC G_*\subset\AC G$, a zwykła pochodna jest pochodną aproksymatywną. Właściwość wykaże następny przykład.
\end{proof}
\begin{przyklad}[Oscylacje w lukach zbioru Cantora]
Niech $K\subset[0,1]$ będzie klasycznym zbiorem Cantora. Na etapie $n\ge1$ usuwanych jest $2^{n-1}$ luk $I_{n,k}=(a_{n,k},b_{n,k})$ długości $3^{-n}$. Definiujemy
\[
F(x)=\begin{cases}
2^{-n}\sin^2\!\bigl(\pi3^n(x-a_{n,k})\bigr),&x\in I_{n,k},\\
0,&x\in K.
\end{cases}
\]
Funkcja jest ciągła: na końcach luk wynosi zero, a amplitudy na lukach wysokich poziomów dążą do zera. Jest klasy $\AC$ na domknięciu każdej luki i stała na $K$, zatem należy do $\AC G[0,1]$.

Pokażemy, że $F\notin\AC G_*[0,1]$. Gdyby $K$ było pokryte przeliczalnie wieloma domkniętymi zbiorami $E_m\cap K$ z $F\in\AC_*(E_m)$, twierdzenie Baire'a zastosowane w $K$ dałoby podstawowy przedział Cantora $J$ pewnego poziomu $r$, dla którego $J\cap K\subset E_m$ dla jednego $m$. Wszystkie luki poziomu $n>r$ zawarte w $J$ mają końce w $E_m$. Ich liczba wynosi $2^{n-r-1}$, suma długości wynosi $2^{n-r-1}3^{-n}\to0$, a suma oscylacji jest równa $2^{n-r-1}2^{-n}=2^{-r-1}$, co przeczy definicji $\AC_*(E_m)$.

Niech $f=F'$ na dopełnieniu $K$, a na $K$ określmy $f=0$. W lukach mamy
$f(x)=\pi(3/2)^n\sin(2\pi3^n(x-a_{n,k}))$.
Ponieważ $\lambda(K)=0$, $F'_{\mathrm{ap}}=f$ prawie wszędzie. Stąd $(\mathrm D_{\mathrm{sz}})\int_0^1 f=F(1)-F(0)=0$. Gdyby $f$ była całkowalna HK, jej pierwotna $G$ należałaby do $\AC G_*$. Różnica $G-F$ miałaby zerową pochodną aproksymatywną prawie wszędzie, więc na mocy poprzedniego twierdzenia byłaby stała, wbrew $F\notin\AC G_*$. Ponadto na każdej luce poziomu $n$ całka modułu pochodnej wynosi $2\cdot2^{-n}$: funkcja rośnie od zera do $2^{-n}$ i maleje do zera. Na całym poziomie suma wynosi $2^{n-1}\cdot2^{1-n}=1$. Z addytywności całki nieujemnej po rozłącznych lukach otrzymujemy $\int_0^1|f|\,\mathrm d\lambda_1=\sum_{n\ge1}1=\infty$, więc $f\notin\mathcal L^1$.
\end{przyklad}
\begin{uwaga}[Granica teorii]
Szeroka całka Denjoya odtwarza funkcje pierwotne klasy $\AC G$ z ich pochodnych aproksymatywnych, lecz nie dowolne funkcje ciągłe z pochodnych prawie wszędzie. Funkcja Cantora ma pochodną aproksymatywną zero prawie wszędzie, ale nie jest klasy $\AC G$, gdyż jest niestała.
\end{uwaga}

% ============================================================
% UZUPELNIENIE MATEMATYCZNE: ROZDZIAL 9
% ============================================================
\subsection*{Różnica między $\AC(E)$ i $\AC_*(E)$}
\begin{przyklad}[Brak kontroli oscylacji na lukach]\label{prz:ac-kontra-acstar}
Niech $K$ będzie klasycznym zbiorem Cantora, a $F$ funkcją z przykładu o oscylacjach w lukach $K$ z tego rozdziału. Na $K$ zachodzi $F=0$, więc $F\in\AC(K)$: przyrosty pomiędzy dowolnymi punktami $K$ są równe zeru. Nie zachodzi natomiast $F\in\AC_*(K)$. Istotnie, dla luk $I_{n,k}$ usuwanych na poziomie $n$ mamy końce w $K$ oraz
\[
\sum_{k=1}^{2^{n-1}}|I_{n,k}|=2^{n-1}3^{-n}\longrightarrow0,
\qquad
\sum_{k=1}^{2^{n-1}}\osc_{\overline I_{n,k}}F
=2^{n-1}2^{-n}=\frac12.
\]
Dla każdego $\eta>0$ znajdziemy zatem rozłączne przedziały o końcach w $K$ i łącznej długości mniejszej niż $\eta$, na których suma oscylacji pozostaje równa $1/2$. To pokazuje, dlaczego w szerokiej teorii kontrolujemy jedynie przyrosty, a w wąskiej całe oscylacje.
\end{przyklad}
\begin{twierdzenie}[Wzór Newtona--Leibniza dla szerokiej całki Denjoya]\label{tw:denjoy-szeroka-ftc}
Niech $F\in C[a,b]\cap\AC G[a,b]$. Wówczas $F'_{\mathrm{ap}}$ istnieje i jest skończona prawie wszędzie. Ustalmy dowolną funkcję $f:[a,b]\to\R$ o skończonych wartościach, równą $F'_{\mathrm{ap}}$ wszędzie tam, gdzie ta pochodna istnieje poza pewnym zbiorem miary zero. Wówczas dla dowolnych $a\le u\le v\le b$ funkcja $f$ jest całkowalna w szerokim sensie Denjoya na $[u,v]$ i
\[
 (\mathrm D_{\mathrm{sz}})\int_u^v f(x)\,\mathrm dx=F(v)-F(u).
\]
\end{twierdzenie}
\begin{proof}
Istnienie pochodnej aproksymatywnej prawie wszędzie ustalono we wcześniejszym twierdzeniu tego rozdziału. Jeśli $[a,b]=\bigcup_n E_n$ jest domkniętym pokryciem z $F\in\AC(E_n)$, to zbiory $E_n\cap[u,v]$ są domknięte w $[u,v]$ i pokrywają ten przedział; dlatego $F|_{[u,v]}\in\AC G[u,v]$. W punktach wewnętrznych $[u,v]$ pochodna aproksymatywna ograniczenia zgadza się z pochodną aproksymatywną $F$, o ile istnieje. Definicja szerokiej całki Denjoya daje tożsamość. Jej jednoznaczność wynika z wcześniejszego twierdzenia o zerowej pochodnej aproksymatywnej.
\end{proof}
\begin{stwierdzenie}[Lokalność i liniowość szerokiej całki]\label{st:denjoy-szeroka-lokalnosc}
Jeśli $f,g$ są całkowalne w szerokim sensie Denjoya na $[a,b]$ i $\alpha,\beta\in\R$, to $\alpha f+\beta g$ jest całkowalna i jej całka jest odpowiednią kombinacją liniową. Jeśli $a<c<b$, to $f$ jest całkowalna na $[a,b]$ wtedy i tylko wtedy, gdy jest całkowalna na $[a,c]$ i $[c,b]$; wówczas całki są addytywne.
\end{stwierdzenie}
\begin{proof}
Dla pierwotnych $F,G$ bierzemy $H=\alpha F+\beta G$. Przez przecięcia domkniętych zbiorów pokrycia uzyskujemy $H\in\AC G$, a pochodna aproksymatywna jest liniowa wszędzie, gdzie obie pochodne istnieją. Dla lokalności ograniczamy pierwotną do podprzedziałów. Odwrotnie, dwie pierwotne na podprzedziałach można skleić, dodając jednej z nich stałą, aby wartości w $c$ się zgadzały. Otrzymujemy ciągłą funkcję, a połączone domknięte pokrycia obu przedziałów świadczą, że jest klasy $\AC G$. Przyrost tej funkcji jest sumą przyrostów po obu częściach.
\end{proof}
\begin{wniosek}[Własność Luzina i ograniczone wahanie]\label{it:cor:acg-bv}
Każda funkcja $F\in\AC G[a,b]$ ma własność Luzina (N): obraz zbioru miary Lebesgue'a zero ma miarę zero. Ponadto
\[
 \AC G[a,b]\cap\BV[a,b]=\AC[a,b].
\]
Jeśli $f$ jest całkowalna w szerokim sensie Denjoya, a $F(x)=(\mathrm D_{\mathrm{sz}})\int_a^x f$, to $F\in\BV[a,b]$ wtedy i tylko wtedy, gdy $f\in\mathcal L^1[a,b]$. Wówczas $V_a^b(F)=\int_a^b|f|\,\mathrm d\lambda_1$.
\end{wniosek}
\begin{proof}
Niech $E_n$ będzie domkniętym pokryciem z definicji $\AC G$, a $u_n$ absolutnie ciągłymi przedłużeniami z lematu \ref{it:lem:ac-affine-extension}. Dla zbioru zerowego $N$ wybieramy borelowski zbiór zerowy $B\supset N$. Każdy $B\cap E_n$ jest borelowski, zatem
$\lambda_1^*(F(B\cap E_n))\le\int_{B\cap E_n}|u_n'|\,\mathrm d\lambda_1=0$.
Z przeliczalnej podaddytywności miary zewnętrznej wynika
$\lambda_1^*(F(N))=0$.

Jeśli $F\in\AC G\cap\BV$, z rozkładu Jordana i twierdzenia o różniczkowaniu funkcji niemalejących (argument z lematu \ref{it:lem:bvstar-null-to-acstar}) wynika, że zwykła pochodna $F'$ istnieje prawie wszędzie i jest całkowalna w sensie Lebesgue'a po uzupełnieniu na zbiorze wyjątkowym. Zwykła pochodna jest również aproksymatywna. Funkcja
\[
 G(x)=F(a)+\int_a^x F'(t)\,\mathrm d\lambda_1(t)
\]
jest absolutnie ciągła, a więc należy do $\AC G$, i ma zwykłą pochodną $F'$ prawie wszędzie (twierdzenie \ref{it:thm:ac-lebesgue-primitive}). Stąd $F-G$ należy do $\AC G$ i ma zerową pochodną aproksymatywną prawie wszędzie. Twierdzenie \ref{it:thm:acg-uniqueness} daje $F-G$ stałą; wartość w $a$ jest zerowa, więc $F=G\in\AC$. Inkluzja odwrotna wynika z definicji oraz ograniczonego wahania funkcji absolutnie ciągłej.

Dla pierwotnej szerokiej całki Denjoya mamy $F'_{\mathrm{ap}}=f$ prawie wszędzie. Jeśli $F\in\BV$, pierwsza część daje $F\in\AC$, a zwykła pochodna $F'$ jest równa $f$ prawie wszędzie przez jednoznaczność pochodnej aproksymatywnej (lemat \ref{it:lem:approx-derivative}). Zatem $f\in\mathcal L^1$. Odwrotnie, jeśli $f\in\mathcal L^1$, jego pierwotna Lebesgue'a jest absolutnie ciągła i ma pochodną aproksymatywną równą $f$ prawie wszędzie. Jednoznaczność pierwotnych szerokiej całki Denjoya daje jej równość z $F$. Wzór na wahanie wynika teraz z twierdzenia \ref{it:thm:denjoy-bv-ac}, ponieważ $F$ jest absolutnie ciągła, a więc jest także pierwotną wąskiej całki Denjoya.
\end{proof}

\begin{uwaga}[Nazewnictwo i oddzielne poziomy teorii]
W terminologii historycznej szeroka całka Denjoya jest także nazywana całką Denjoya--Chinczyna. Jej opis wymaga $\AC G$ i pochodnej aproksymatywnej. Wąska całka Denjoya, równoważna HK i Perronowi, wymaga $\AC G_*$ oraz zwykłej pochodnej prawie wszędzie. Przyjęcie tylko $\AC G$ i zwykłej pochodnej bez dodatkowych warunków prowadzi do innego schematu definicyjnego; obu pojęć nie należy utożsamiać bez dowodu.
\end{uwaga}


\section{Miary związane z całkowaniem}
Określenie „miara Lebesgue'a--Riemanna” nie jest standardową nazwą jednej konstrukcji. Rozróżniamy zawartość Jordana, miarę Lebesgue'a i miary Lebesgue'a--Stieltjesa.
\subsection*{Miara Jordana („riemannowska”)}
Dla ograniczonego $A\subset\R$ kładziemy $c_*(A)=\sup\{|E|:E\subset A,\ E\ \text{elementarny}\}$ oraz $c^*(A)=\inf\{|E|:A\subset E,\ E\ \text{elementarny}\}$, gdzie zbiory elementarne są skończonymi sumami zbiorów ograniczonych przedziałów, a $|E|$ to suma ich długości po sprowadzeniu do przedziałów o rozłącznych wnętrzach. Zbiór jest mierzalny w sensie Jordana, jeżeli $c_*(A)=c^*(A)$.
\begin{stwierdzenie}
Dla $A\subset[a,b]$ następujące warunki są równoważne: (i) $A$ jest mierzalny w sensie Jordana; (ii) $\ind A\in\mathcal R[a,b]$; (iii) $\lambda(\partial A)=0$. Wówczas $c(A)=\int_a^b\ind A(x)\,\mathrm dx=\lambda(A)$.
\end{stwierdzenie}
\begin{proof}
Niech $P$ będzie podziałem $[a,b]$. Suma domkniętych komórek zawartych w $A$ jest zbiorem elementarnym zawartym w $A$, a jej długość jest sumą dolną $L(\chi_A,P)$. Suma domkniętych komórek przecinających $A$ zawiera $A$, a jej długość jest sumą górną $U(\chi_A,P)$. Zatem
\[
 L(\chi_A,P)\le c_*(A)\le c^*(A)\le U(\chi_A,P).
\]
Jeżeli $\chi_A$ jest całkowalna Riemanna, różnica skrajnych wyrazów może być dowolnie mała, więc $c_*(A)=c^*(A)$, a wspólna wartość jest całką Riemanna.

Odwrotnie, jeżeli $c_*(A)=c^*(A)$, dla $\varepsilon>0$ wybieramy zbiory elementarne $E\subset A\subset H$ z $|H|-|E|<\varepsilon$. Każdy punkt brzegu $A$ leży w $H\setminus E$ albo w brzegu $E$ lub $H$: poza tymi zbiorami jest bowiem albo we wnętrzu $E$, albo na zewnątrz domknięcia $H$. Brzegi zbiorów elementarnych są skończone, a $H\setminus E$ jest elementarny i ma długość $|H|-|E|$. Stąd $\lambda_1(\partial A)\le |H|-|E|<\varepsilon$. Dowolność $\varepsilon$ daje (iii).
Punkty nieciągłości $\chi_A$ tworzą względny brzeg $A$ w $[a,b]$, różniący się od brzegu w $\R$ co najwyżej końcami $a,b$. Z kryterium Lebesgue'a całkowalności Riemanna wynika (iii) $\Rightarrow$ (ii).
Jeżeli brzeg jest zerowy, zbiór $A$ różni się od swojego wnętrza o podzbiór brzegu, więc jest mierzalny w sensie Lebesgue'a. Dla dowolnych $E\subset A\subset H$ jak wyżej zachodzi $|E|\le\lambda_1(A)\le |H|$. Przechodząc do kresów, otrzymujemy $c(A)=\lambda_1(A)$.
\end{proof}
Zbiory mierzalne w sensie Jordana tworzą algebrę, nie $\sigma$-algebrę: $\Q\cap[0,1]$ jest przeliczalną sumą zbiorów jednopunktowych mierzalnych w sensie Jordana, ale sam nie jest mierzalny w tym sensie. Przejście od skończonych do przeliczalnych pokryć prowadzi do miary zewnętrznej Lebesgue'a $\lambda^*(A)=\inf\{\sum_k|I_k|:A\subset\bigcup_kI_k\}$ oraz konstrukcji Carathéodory'ego. Dla ograniczonego $A$ mamy $c_*(A)\le\lambda_*(A)\le\lambda^*(A)\le c^*(A)$.

\subsection*{Miary Lebesgue'a--Stieltjesa}
\begin{twierdzenie}[Konstrukcja miary Lebesgue'a--Stieltjesa]\label{it:thm:stieltjes-construction}
Niech $\alpha:\R\to\R$ będzie niemalejąca i prawostronnie ciągła. Istnieje dokładnie jedna lokalnie skończona miara borelowska $\mu_\alpha$ na $\R$, dla której
\[
 \mu_\alpha((u,v])=\alpha(v)-\alpha(u)\qquad(u<v).
\]
Jej uzupełnienie nazywamy miarą Lebesgue'a--Stieltjesa. Odwrotnie, każda lokalnie skończona miara borelowska $\nu$ na $\R$ ma tę postać. Funkcja $\alpha$ jest wyznaczona z dokładnością do stałej.
\end{twierdzenie}
\begin{proof}
Na pierścieniu $\mathcal A$ skończonych sum rozłącznych ograniczonych przedziałów $(u,v]$ definiujemy $m$ jako sumę przyrostów $\alpha$. Wspólne rozdrobnienie skończonych rodzin w ich końcach oraz teleskopowanie przyrostów dowodzą niezależności od przedstawienia i skończonej addytywności. W szczególności $m$ jest nieujemna, monotoniczna i skończenie podaddytywna.

Udowodnimy ciągłość z góry w zbiorze pustym. Niech $A_n\in\mathcal A$, $A_n\downarrow\varnothing$ i niech $\varepsilon>0$. Każdy przedział $(u,v]$ w skończonym przedstawieniu $A_n$ zastępujemy zwartym podprzedziałem $[u+h,v]$, $0<h<v-u$. Na mocy prawostronnej ciągłości $\alpha$ wybieramy $h$ tak, by suma wartości $m((u,u+h])$ dla wszystkich komórek $A_n$ była mniejsza od $\varepsilon 2^{-n}$. Otrzymujemy zbiór zwarty $K_n\subset A_n$ oraz $D_n\in\mathcal A$ z $A_n\setminus K_n\subset D_n$ i $m(D_n)<\varepsilon 2^{-n}$. Dla pustego $A_n$ bierzemy $K_n=D_n=\varnothing$.
Ponieważ $\bigcap_n K_n\subset\bigcap_n A_n=\varnothing$, własność skończonych przecięć zbiorów zwartych daje $N$ takie, że $\bigcap_{n=1}^N K_n=\varnothing$. Dla $x\in A_N$ mamy $x\in A_n$ dla każdego $n\le N$, a $x\notin K_n$ dla pewnego takiego $n$. Zatem $A_N\subset\bigcup_{n=1}^N D_n$ i
\[
 m(A_N)\le \sum_{n=1}^Nm(D_n)<\varepsilon.
\]
Monotoniczność daje $m(A_n)\to0$.

Jeśli $A=\bigsqcup_{n\ge1}B_n\in\mathcal A$ oraz $B_n\in\mathcal A$, reszty $R_N=A\setminus\bigcup_{n=1}^N B_n$ należą do $\mathcal A$ i maleją do zbioru pustego. Skończona addytywność daje $m(A)=\sum_{n=1}^N m(B_n)+m(R_N)$; przejście do granicy dowodzi, że $m$ jest przedmiarą. Twierdzenie Carathéodory'ego o rozszerzaniu przedmiary daje miarę na $\sigma(\mathcal A)=\mathcal B(\R)$. Jednoznaczność wynika z $\sigma$-skończoności: przedziały $(-n,n]$ pokrywają prostą i mają skończone masy. Każdy zbiór zwarty leży w jednym z tych przedziałów, więc miara jest lokalnie skończona.

Odwrotnie, dla lokalnie skończonej $\nu$ definiujemy
\[
 \alpha(x)=
 \begin{cases}
 \nu((0,x]),&x\ge0,\\
 -\nu((x,0]),&x<0.
 \end{cases}
\]
W obu przypadkach $\alpha(v)-\alpha(u)=\nu((u,v])$, także dla $u<0\le v$ przez addytywność. Z nieujemności miary wynika monotoniczność. Dla $y\downarrow x$ mamy $\nu((x,y])\to0$ z ciągłości miary z góry na ograniczonym przedziale, a zatem $\alpha$ jest prawostronnie ciągła. Jednoznaczność miary daje $\nu=\mu_\alpha$. Jeżeli dwie funkcje mają te same przyrosty, ich różnica jest stała.
\end{proof}
Dla $\alpha(x)=x$ otrzymujemy $\lambda_1$, a dla $\alpha=\chi_{[c,\infty)}$ miarę Diraca $\delta_c$, ponieważ obie spełniają właściwy wzór na przedziałach $(u,v]$. Dla integratora prawostronnie ciągłego w porównaniu z całką Riemanna--Stieltjesa na $[a,b]$ ograniczamy miarę do $(a,b]$ (zob. twierdzenie \ref{it:thm:rs-oscillation} i stwierdzenie \ref{st:ks-miara-konwencja}).
\subsection*{Rozkład miary i całkowanie względem gęstości}
% Import TM-02, TM-17, TM-18; zachowane etykiety źródłowe.
\begin{definicja}[Bezwzględna ciągłość miary]
\label{def:ac-measure-17}
Niech $\mu$ i $\nu$ będą miarami nieujemnymi na $(X,\mathcal F)$.
Mówimy, że $\nu$ jest \emph{bezwzględnie ciągła względem $\mu$},
i piszemy $\nu\ll\mu$, jeżeli dla każdego $E\in\mathcal F$
z warunku $\mu(E)=0$ wynika $\nu(E)=0$.
Dla skończonej miary zespolonej $\eta$ przyjmujemy tę samą definicję,
równoważnie $|\eta|\ll\mu$, gdzie $|\eta|$ jest wariacją całkowitą.
\end{definicja}
Miary nieujemne $\nu,\mu$ na wspólnej przestrzeni mierzalnej są osobliwe ($\nu\perp\mu$), gdy istnieje mierzalny zbiór $S$ taki, że $\mu(S)=0$ i $\nu(X\setminus S)=0$. Korzystamy z twierdzenia o rozkładzie Lebesgue'a oraz twierdzenia Radona--Nikodyma (Teoria miary, rozdziały 17--18).
\begin{definicja}[Atom i miara bezatomowa]
\label{def:atom}
Niech $(X,\mathcal F,\mu)$ będzie przestrzenią z miarą.
Zbiór $A\in\mathcal F$ nazywamy \emph{atomem miary $\mu$},
jeżeli $\mu(A)>0$ oraz dla każdego $B\in\mathcal F$
zawartego w $A$ zachodzi $\mu(B)=0$ lub
$\mu(A\setminus B)=0$.
Miarę $\mu$ nazywamy \emph{bezatomową}, jeśli nie ma atomów.
\end{definicja}
\begin{twierdzenie}[Dalszy rozkład miary osobliwej na prostej]
\label{thm:atomic-cont-singular-17}
Niech $\nu$ będzie skończoną miarą borelowską na
$\mathbb R$. Zbiór
$A=\{x\in\mathbb R:\nu(\{x\})>0\}$ jest co najwyżej
przeliczalny, a miary
\[
 \nu_{\mathrm{at}}(E):=\nu(E\cap A),\qquad
 \nu_{\mathrm{na}}(E):=\nu(E\setminus A)
 \quad(E\in\mathcal B(\mathbb R))
\]
są jednoznacznie określone. Pierwsza ma postać
$\nu_{\mathrm{at}}=\sum_{a\in A}\nu(\{a\})\delta_a$,
a druga jest bezatomowa. Jeśli $\nu\perp\lambda_1$,
to również $\nu_{\mathrm{at}}\perp\lambda_1$ i
$\nu_{\mathrm{na}}\perp\lambda_1$.
\end{twierdzenie}
\begin{proof}
Dla każdego $n\ge1$ zbiór
$A_n=\{x:\nu(\{x\})\ge1/n\}$ jest skończony:
gdyby zawierał $m$ różnych punktów, mielibyśmy
$\nu(\mathbb R)\ge m/n$. Ponieważ
$A=\bigcup_{n\ge1}A_n$, jest przeliczalny.
Przeliczalna addytywność daje podany wzór na
$\nu_{\mathrm{at}}$ oraz rozkład $\nu$.
Z definicji $\nu_{\mathrm{na}}(\{x\})=0$ dla
każdego $x$. Wykażemy, że miara ta nie ma także
żadnych atomów w sensie teorii miary.

Przypuśćmy, że mierzalny zbiór $B$ jest atomem
$\nu_{\mathrm{na}}$ o dodatniej masie $b$:
każdy mierzalny podzbiór $B$ ma masę $0$ lub $b$.
Ponieważ $\mathbb R=\bigcup_{m\ge1}[-m,m]$,
pewien domknięty ograniczony przedział $I_0$
spełnia $\nu_{\mathrm{na}}(B\cap I_0)=b$.
Dzielimy $I_0$ na dwa domknięte przedziały równej
długości. Co najmniej jeden z nich, nazwany $I_1$,
ma $\nu_{\mathrm{na}}(B\cap I_1)>0$, a więc masę $b$.
Powtarzając, otrzymujemy zstępujący ciąg przedziałów
zamkniętych $I_n$ o długościach malejących do zera,
dla których $\nu_{\mathrm{na}}(B\cap I_n)=b$.
Ich przecięcie jest singletonem $\{x\}$.
Z ciągłości skończonej miary z góry mamy
$\nu_{\mathrm{na}}(B\cap\{x\})=b>0$,
sprzecznie z zerowością wszystkich singletonów.
Zatem $\nu_{\mathrm{na}}$ jest bezatomowa.

Jeśli $\nu$ skupia się na zbiorze borelowskim $S$
o $\lambda_1(S)=0$, to obie części jako ograniczenia
$\nu$ również skupiają się na $S$, więc są osobliwe.
\end{proof}
\begin{lemat}[Całkowanie względem miary o gęstości]
\label{lem:integration-density-18}
Niech $\lambda$ będzie miarą nieujemną na $(X,\mathcal F)$,
$w:X\to[0,\infty]$ funkcją mierzalną oraz
$\mu(E)=\int_Ew\,\mathrm d\lambda$ dla $E\in\mathcal F$.
Wówczas dla każdej mierzalnej $\varphi:X\to[0,\infty]$
zachodzi
\begin{equation}\label{eq:integration-density-18}
\int_X\varphi(x)\,\mathrm d\mu(x)
=\int_X\varphi(x)w(x)\,\mathrm d\lambda(x),
\end{equation}
przy konwencji $0\cdot\infty=0$ dla iloczynu
nieujemnych funkcji mierzalnych. Jeśli $\varphi$ jest
zespolona i całkowalna względem $\mu$, to analogiczna
równość zachodzi dla całek zespolonych.
\end{lemat}
\begin{proof}
Dla $\varphi=\chi_A$ jest to definicja $\mu$.
Z liniowości dla funkcji prostych nieujemnych otrzymujemy
równość dla wszystkich takich funkcji. Każdą funkcję
$\varphi\ge0$ aproksymujemy z dołu ciągiem mierzalnych
funkcji prostych $\varphi_n\uparrow\varphi$.
Dwukrotne zastosowanie twierdzenia Beppo Leviego
daje \eqref{eq:integration-density-18}. W przypadku
funkcji zespolonej stosujemy tę tożsamość do części
dodatnich i ujemnych jej części rzeczywistej i urojonej,
które są całkowalne dzięki założeniu.
\end{proof}

\begin{twierdzenie}[Kryterium oparte na rozkładzie miary]
Niech $\mu_\alpha$ będzie skończoną miarą Lebesgue'a--Stieltjesa ograniczoną do $(a,b]$. Istnieje rozkład
\[
\mu_\alpha=w\lambda_1+\mu_{\mathrm{sc}}+\sum_{k\in J}c_k\delta_{t_k},
\]
gdzie wszystkie miary rozpatrujemy na $(a,b]$, $w\ge0$ jest mierzalna, $\mu_{\mathrm{sc}}\perp\lambda_1$ jest bezatomowa, $J$ jest zbiorem co najwyżej przeliczalnym, $t_k$ są różne, $c_k>0$ i $\sum_{k\in J}c_k<\infty$. Dla każdej funkcji borelowskiej $f:(a,b]\to\R$ zachodzi
\[
\int_{(a,b]}|f|\,\mathrm d\mu_\alpha
=\int_{(a,b]}|f|w\,\mathrm d\lambda_1+
\int_{(a,b]}|f|\,\mathrm d\mu_{\mathrm{sc}}+
\sum_{k\in J}c_k|f(t_k)|.
\]
Funkcja jest bezwzględnie całkowalna względem $\mu_\alpha$ wtedy i tylko wtedy, gdy prawa strona jest skończona.
\end{twierdzenie}
\begin{proof}
Obie miary $\mu_\alpha$ i $\lambda_1|_{(a,b]}$ są skończone, więc twierdzenie o rozkładzie Lebesgue'a daje część bezwzględnie ciągłą i osobliwą. Twierdzenie Radona--Nikodyma daje nieujemną gęstość $w$ pierwszej części względem $\lambda_1$. Drugą część przedłużamy przez zero na $\R$ i stosujemy twierdzenie \ref{thm:atomic-cont-singular-17}. Masa atomów sumuje się do masy części atomowej, nie większej od $\mu_\alpha((a,b])$.
Dla nieujemnej borelowskiej funkcji prostej równość całek względem sumy miar wynika z definicji całki i przeliczalnej addytywności sum nieujemnych. Dla dowolnej nieujemnej funkcji mierzalnej przybliżamy ją od dołu rosnącym ciągiem funkcji prostych i stosujemy twierdzenie Beppo Leviego do każdej miary oraz do sumy po $J$ (miara licząca). Lemat \ref{lem:integration-density-18} identyfikuje całkę względem $w\lambda_1$, a całka względem $\delta_{t_k}$ jest wartością w $t_k$. Stosujemy otrzymaną równość do $|f|$.
\end{proof}

\subsection*{Miara Cantora: obliczenia i porównanie całek}
\begin{przyklad}[Miara Cantora — osobliwość bez atomów]
\label{ex:cantor-measure-17}
Podamy konstrukcję niewymagającą wcześniejszego twierdzenia
o miarach Lebesgue'a--Stieltjesa. Dla $x\in[0,1)$ niech
\[
b_n(x)=\lfloor2^nx\rfloor-2\lfloor2^{n-1}x\rfloor\in\{0,1\},
\qquad n\ge1.
\]
Są to cyfry rozwinięcia dwójkowego, przy czym dla liczb
dwójkowo wymiernych przyjęliśmy rozwinięcie zakończone
samymi zerami. Określamy
$T(x)=\sum_{n\ge1}2b_n(x)3^{-n}$ dla $x\in[0,1)$
oraz $T(1)=1$. Każda funkcja $b_n$ jest borelowska,
a szereg definiujący $T$ jest jednostajnie zbieżny,
więc $T:[0,1]\to[0,1]$ jest borelowska. Jej wartości
należą do środkowego zbioru Cantora $C$, ponieważ ich
rozwinięcia trójkowe używają wyłącznie cyfr $0$ i $2$.

Niech $\kappa=T_*(\lambda_1|_{[0,1]})$. Jest to miara
probabilistyczna, $\kappa(C)=1$. Na poziomie $n$ konstrukcji zbioru Cantora
każdy z $2^n$ rozłącznych przedziałów podstawowych
$I_{\varepsilon_1,\ldots,\varepsilon_n}$ odpowiada
słowu $(\varepsilon_1,\ldots,\varepsilon_n)\in\{0,1\}^n$:
jego punkty mają rozwinięcie trójkowe rozpoczynające się
od cyfr $2\varepsilon_1,\ldots,2\varepsilon_n$.
Z wyjątkiem przeliczalnego zbioru końców przedziałów dwójkowych
przeciwobraz $I_{\varepsilon_1,\ldots,\varepsilon_n}$
przez $T$ jest zbiorem liczb, których pierwsze $n$
cyfr dwójkowych to $\varepsilon_1,\ldots,\varepsilon_n$.
Jest to przedział dwójkowy długości $2^{-n}$, więc
$\kappa(I_{\varepsilon_1,\ldots,\varepsilon_n})=2^{-n}$.
Każdy $y\in C$ należy na każdym poziomie do jednego
takiego przedziału, zatem
$\kappa(\{y\})\le2^{-n}$ dla wszystkich $n$ i
$\kappa(\{y\})=0$. Dla $y\notin C$ mamy
$\kappa(\{y\})=0$, gdyż $\kappa$ jest skupiona na $C$.
Miara $\kappa$ jest więc bezatomowa. Ponieważ
$\lambda_1(C)=0$, otrzymujemy $\kappa\perp\lambda_1$.
Pokazuje to, że \emph{osobliwość} nie jest tym samym co
\emph{dyskretność}. Dystrybuanta miary $\kappa$ jest
klasyczną funkcją Cantora.
\end{przyklad}
W dalszych przykładach oznaczamy funkcję Cantora przez $C$, środkowotrzeci zbiór Cantora przez $K$, a skonstruowaną miarę przez $\mu_C=\kappa$.
\begin{przyklad}[Momenty miary osobliwej]
Dla losowej zmiennej $X$ o rozkładzie $\mu_C$ mamy $X\stackrel d= (B+X')/3$, gdzie $B$ przyjmuje wartości $0,2$ z prawdopodobieństwami $1/2$, a $X'$ jest niezależną kopią $X$. Wynika to z tego, że dla miary Lebesgue'a każdy skończony ciąg cyfr dwójkowych długości $n$ ma prawdopodobieństwo $2^{-n}$. Cyfry są więc niezależne i każda ma rozkład Bernoulliego z parametrem $1/2$; po oddzieleniu pierwszej cyfry szereg pozostałych cyfr ma po przeskalowaniu ten sam rozkład i jest niezależny od pierwszej cyfry. Stąd
$\mathbb E X=\frac13\mathbb E X+\frac13$, czyli $\mathbb E X=1/2$, a także $\mathbb E X^2=\frac19\mathbb E X^2+\frac29\mathbb E X+\frac29$, czyli $\mathbb E X^2=3/8$. Zatem
\[
\int_0^1 1\,\mathrm dC=1,\qquad \int_0^1 x\,\mathrm dC=\frac12,\qquad \int_0^1 x^2\,\mathrm dC=\frac38.
\]
Ponadto z całkowania przez części, ponieważ $C$ jest ciągła i ma wahanie ograniczone, wynika $2\int_0^1 C\,\mathrm dC=C(1)^2-C(0)^2=1$. Choć $C'=0$ prawie wszędzie, nie wolno zastępować $\mathrm dC$ przez $C'(x)\,\mathrm dx$: zgubilibyśmy część osobliwą miary. Ponieważ $x\mapsto x^2$ jest ciągła, równości obowiązują zarówno dla całek Riemanna--Stieltjesa, jak i Lebesgue'a--Stieltjesa.
\end{przyklad}
\begin{przyklad}[Ta sama funkcja, różne miary]
Niech $K$ będzie zbiorem Cantora. Funkcja $\ind K$ jest całkowalna Riemanna względem $\lambda$ z całką $0$, ponieważ $K$ ma miarę Lebesgue'a zero. Jest również całkowalna Lebesgue'a--Stieltjesa względem $\mu_C$ z całką $1$. Natomiast zwykła całka Riemanna--Stieltjesa $\int_0^1\ind K\,\mathrm dC$ nie istnieje: punkty nieciągłości funkcji $\ind K$ tworzą $K$, a $\mu_C(K)=1$, więc nie jest spełnione kryterium oscylacyjne.
\end{przyklad}

% ============================================================
% UZUPELNIENIE MATEMATYCZNE: ROZDZIAL 10
% ============================================================
\subsection*{Atomy i wartości na końcach: porównanie z KS}
\begin{stwierdzenie}[Reprezentacja całki KS przez miarę przy prawostronnej ciągłości]\label{st:ks-miara-konwencja}
Niech $\alpha:[a,b]\to\R$ będzie niemalejąca, prawostronnie ciągła na $[a,b)$ i niech $\mu_\alpha$ będzie miarą na $(a,b]$ zdefiniowaną przez $\mu_\alpha((s,t])=\alpha(t)-\alpha(s)$. Jeśli $f:[a,b]\to\R$ jest regulowana, to całka Kurzweila--Stieltjesa $\int_a^b f\,\mathrm d\alpha$ istnieje i
\[
\int_a^b f\,\mathrm d\alpha=\int_{(a,b]}f\,\mathrm d\mu_\alpha.
\]
W szczególności $\mu_\alpha(\{c\})=\alpha(c)-\alpha(c-)$ dla $a<c\le b$, a punkt $a$ nie wnosi atomu w tej konwencji.
\end{stwierdzenie}
\begin{proof}
Z twierdzenia \ref{tw:ks-measure} wiemy, że dla integrandu regulowanego i integratora BV całka istnieje oraz że zgadza się z całką Lebesgue'a--Stieltjesa w przyjętej konwencji. Tutaj $\alpha$ jest monotoniczna, więc ma wahanie ograniczone. Przechodząc z $s\uparrow c$ w równości $\mu_\alpha((s,c])=\alpha(c)-\alpha(s)$, otrzymujemy wskazaną masę atomu.
\end{proof}
\begin{przyklad}[Skok prawostronny i wartość w punkcie]
Dla $a<c<b$ zdefiniujmy $g(x)=\ind{(c,b]}(x)$. Funkcja $g$ nie jest prawostronnie ciągła w $c$, a jej przyrost na przedziale $[c,t]$, $t>c$, wynosi $1$. Dla dowolnej regulowanej funkcji $f$ całka KS istnieje i $\int_a^b f\,\mathrm dg=f(c)$ (jest to wkład prawego skoku). Zastąpienie $g$ przez prawostronnie ciągłą funkcję $g^+(x)=\ind{[c,b]}(x)$ nie zmienia tu wartości całki KS dla regulowanej $f$, lecz przenosi skok z prawej strony $c$ na lewą. Zwykła miara Stieltjesa jest związana z wybranym prawostronnie ciągłym reprezentantem.
\end{przyklad}
\begin{przyklad}[Dwa skoki przeciwnego znaku w jednym punkcie]
Niech $g(x)=\ind{\{c\}}(x)$, gdzie $a<c<b$. Funkcja ma skok lewy $\Delta^-g(c)=1$ i prawy $\Delta^+g(c)=-1$. Dla regulowanej $f$ wzór skokowy z rozdziału 6 daje $\int_a^b f\,\mathrm dg=f(c)-f(c)=0$. Nie należy utożsamiać takiej funkcji $g$ z dystrybuantą dodatniej miary: $g$ nie jest monotoniczna ani prawostronnie ciągła.
\end{przyklad}
\begin{przyklad}[Dwa różne końce przedziału]
Na $[a,b]$ niech $g_a=\ind{(a,b]}$ oraz $g_b=\ind{\{b\}}$. Dla dowolnej regulowanej $f$ zachodzi
$\int_a^b f\,\mathrm dg_a=f(a)$ oraz $\int_a^b f\,\mathrm dg_b=f(b)$ w sensie KS. Są to odpowiednio skok prawostronny w $a$ i skok lewostronny w $b$. W miarze $\mu_\alpha$ na $(a,b]$ pierwszy z nich nie występuje, ponieważ funkcja $g_a$ nie jest prawostronnie ciągła w $a$; drugi odpowiada atomowi w $b$.
\end{przyklad}

\providecommand{\atlasref}[2]{#2}
\section{Całka Lebesgue'a}\label{ch:calkowanie-lebesgue}

Ograniczamy się do podstawowych twierdzeń o zbieżności całek.
Konstrukcję całki, mierzalność, całkowalność i rachunek całek
omawia skrypt \emph{Teoria miary}, w szczególności
\atlasref{ME-06::thm:simple-up}{ME-06: aproksymacja funkcji mierzalnych},
\atlasref{ME-07::ch:calkowanie-lebesgue}{ME-07: konstrukcja całki Lebesgue'a}
oraz \atlasref{ME-08::thm:mct}{ME-08: twierdzenia graniczne}.
Dalszą teorię rozwinięto w rozdziałach ME-09--ME-14
(rodzaje zbieżności, regularność, całki z parametrem,
miary produktowe, zamiana zmiennych i przestrzenie $L^p$)
oraz ME-25 (różniczkowanie całek i miar).
W szczególności zob.
\atlasref{ME-11::ch:parametry-calki}{ME-11: całki zależne od parametru},
\atlasref{ME-12::thm:fubini-12}{ME-12: twierdzenia Tonellego i Fubiniego},
\atlasref{ME-13::thm:change-variables-13}{ME-13: zamiana zmiennych},
\atlasref{ME-14::thm:riesz-fischer-14}{ME-14: przestrzenie $L^p$}
i \atlasref{ME-25::thm:lebesgue-diff-25}{ME-25: różniczkowanie całek}.
Nie powtarzamy konstrukcji ani dowodów tych wyników.

Niech $(X,\mathcal F,\mu)$ będzie dowolną przestrzenią z miarą.
Rozważamy funkcje mierzalne; $L^1(\mu)$ oznacza przestrzeń
klas funkcji całkowalnych modulo równość prawie wszędzie.
Nie zakładamy zupełności miary ani $\sigma$-skończoności.

\subsection{Dlaczego nie wolno bezwarunkowo przestawiać granicy i całki?}\label{ch:twierdzenia-graniczne}
Na $(0,1)$ z miarą Lebesgue'a $\lambda_1$ funkcje
$f_n=n\chi_{(0,1/n)}$ zbiegają punktowo do zera:
dla ustalonego $x>0$ i $n>1/x$ mamy $f_n(x)=0$.
Jednak $\int_{(0,1)} f_n\,\mathrm d\lambda_1
=n\lambda_1((0,1/n))=1$ dla każdego $n$, podczas gdy
całka funkcji granicznej wynosi zero.
Pokazuje to zarazem, że w lemacie Fatou nierówność
może być ostra.
Poniższe twierdzenia określają warunki zamiany granicy
i całkowania; pełna teoria znajduje się w
\atlasref{ME-08::thm:mct}{ME-08}. Symbol $f_n\uparrow f$
oznacza niemalejącą zbieżność punktową.

\subsection{Twierdzenie Beppo Leviego}


\begin{twierdzenie}[O zbieżności monotonicznej; Beppo Levi]
\label{thm:mct}
Niech $(f_n)_{n\ge1}$ będzie ciągiem mierzalnych funkcji
$f_n:X\to[0,\infty]$, takim że dla każdego $x\in X$
$f_n(x)\le f_{n+1}(x)$ dla każdego $n\ge1$. Niech
$f(x)=\lim_{n\to\infty}f_n(x)=\sup_nf_n(x)$.
Wtedy $f$ jest mierzalna i
\begin{equation}\label{eq:mct}
\int_Xf(x)\,\mathrm d\mu(x)
=\lim_{n\to\infty}\int_Xf_n(x)\,\mathrm d\mu(x)
=\sup_{n\ge1}\int_Xf_n(x)\,\mathrm d\mu(x).
\end{equation}
Dopuszczamy wartość $+\infty$ i nie wymagamy
żadnego założenia o $\mu(X)$.
\end{twierdzenie}
\noindent\emph{Dowód:} zob. \atlasref{ME-08::thm:mct}{Teoria miary, rozdział 8 (dowód tego twierdzenia)}.

\begin{wniosek}[Konstruktywne obliczanie całki]
\label{cor:simple-approx-integral}
Niech $(X,\mathcal F,\mu)$ będzie przestrzenią z miarą.
Jeżeli $f:X\to[0,\infty]$ jest mierzalna, a
$(s_n)$ jest \emph{dowolnym} ciągiem nieujemnych
funkcji prostych mierzalnych takim, że $s_n\uparrow f$
wszędzie, to
$\int_Xf\,\mathrm d\mu(x)=
\lim_{n\to\infty}\int_Xs_n\,\mathrm d\mu(x)$.
W szczególności można użyć funkcji z
\atlasref{ME-06::thm:simple-up}{ME-06: twierdzeniu o aproksymacji funkcjami prostymi}.
\end{wniosek}
\noindent\emph{Dowód:} zob. \atlasref{ME-08::cor:simple-approx-integral}{Teoria miary, rozdział 8 (dowód tego twierdzenia)}.

\begin{twierdzenie}[Zbieżność monotoniczna prawie wszędzie]
\label{thm:mct-ae}
Niech $f_n,f:X\to[0,\infty]$ będą mierzalne.
Jeśli dla każdego $n\ge1$ zachodzi $f_n\le f_{n+1}$ prawie wszędzie oraz $f_n\to f$
$\mu$-prawie wszędzie, to zachodzi równość
\eqref{eq:mct}.
\end{twierdzenie}
\noindent\emph{Dowód:} zob. \atlasref{ME-08::thm:mct-ae}{Teoria miary, rozdział 8 (dowód tego twierdzenia)}.

\subsection{Lemat Fatou}


\begin{twierdzenie}[Lemat Fatou]
\label{thm:fatou}
Niech $f_n:X\to[0,\infty]$ będą funkcjami mierzalnymi.
Wtedy
\begin{equation}\label{eq:fatou}
\int_X\liminf_{n\to\infty}f_n(x)\,\mathrm d\mu(x)
\le\liminf_{n\to\infty}\int_Xf_n(x)\,\mathrm d\mu(x).
\end{equation}
Nie wymagamy monotoniczności ciągu ani skończoności miary.
\end{twierdzenie}
\noindent\emph{Dowód:} zob. \atlasref{ME-08::thm:fatou}{Teoria miary, rozdział 8 (dowód tego twierdzenia)}.

\subsection{Twierdzenie Lebesgue'a o zbieżności zmajoryzowanej}


\begin{twierdzenie}[O zbieżności zmajoryzowanej; Lebesgue]
\label{thm:dct}
Niech $f_n:X\to\mathbb R$ będą mierzalne i
$f_n(x)\to f(x)$ dla $\mu$-prawie wszystkich $x$, gdzie
$f:X\to\mathbb R$ jest mierzalna. Załóżmy, że
istnieje nieujemna całkowalna funkcja mierzalna
$g:X\to[0,\infty]$ taka, że
$|f_n(x)|\le g(x)$ dla wszystkich $n$ i
$\mu$-prawie wszystkich $x$, przy czym zbiór
wyjątkowy można wybrać wspólny dla całego ciągu.
Wówczas $f_n$ i $f$ są całkowalne oraz
\begin{equation}\label{eq:dct}
\lim_{n\to\infty}\int_Xf_n(x)\,\mathrm d\mu(x)
=\int_Xf(x)\,\mathrm d\mu(x).
\end{equation}
\end{twierdzenie}
\noindent\emph{Dowód:} zob. \atlasref{ME-08::thm:dct}{Teoria miary, rozdział 8 (dowód tego twierdzenia)}.

\begin{wniosek}[Zbieżność zmajoryzowana w $L^1$]
\label{cor:dct-l1}
Przy założeniach twierdzenia \ref{thm:dct} zachodzi
\[
\lim_{n\to\infty}\int_X|f_n(x)-f(x)|\,\mathrm d\mu(x)=0.
\]
Dla funkcji zespolonych wniosek i twierdzenie
\ref{thm:dct} także zachodzą, jeśli zastąpić
wartość bezwzględną modułem.
\end{wniosek}
\noindent\emph{Dowód:} zob. \atlasref{ME-08::cor:dct-l1}{Teoria miary, rozdział 8 (dowód tego twierdzenia)}.

\begin{wniosek}[Twierdzenie o zbieżności ograniczonej]
\label{cor:bct}
Niech $\mu(X)<\infty$ oraz $f_n,f:X\to\mathbb R$
będą mierzalne. Jeśli $f_n\to f$ prawie wszędzie
i istnieje $M\in[0,\infty)$ takie, że $|f_n|\le M$
prawie wszędzie dla wszystkich $n$, to wszystkie $f_n$ oraz $f$ są całkowalne i
$\int_X f_n\,\mathrm d\mu\to\int_X f\,\mathrm d\mu$ oraz $\int_X|f_n-f|\,\mathrm d\mu\to0$.
\end{wniosek}
\noindent\emph{Dowód:} zob. \atlasref{ME-08::cor:bct}{Teoria miary, rozdział 8 (dowód tego twierdzenia)}.

\subsection{Szeregi funkcji i całek}


\begin{twierdzenie}[Szeregi nieujemnych funkcji mierzalnych]
\label{thm:sum-nonnegative-integrals}
Niech $f_n:X\to[0,\infty]$ będą mierzalne. Wtedy
\begin{equation}\label{eq:sum-nonnegative-integrals}
\int_X\sum_{n=1}^{\infty}f_n(x)\,\mathrm d\mu(x)
=\sum_{n=1}^{\infty}\int_Xf_n(x)\,\mathrm d\mu(x),
\end{equation}
gdzie obie strony mogą być nieskończone.
\end{twierdzenie}
\noindent\emph{Dowód:} zob. \atlasref{ME-08::thm:sum-nonnegative-integrals}{Teoria miary, rozdział 8 (dowód tego twierdzenia)}.

\begin{twierdzenie}[Całkowanie szeregów bezwzględnie sumowalnych]
\label{thm:sum-absolutely-integrable}
Niech $f_n:X\to\mathbb C$ będą mierzalne i
całkowalne oraz niech
$\sum_{n\ge1}\int_X|f_n(x)|\,\mathrm d\mu(x)<\infty$.
Wówczas $\sum_nf_n(x)$ jest bezwzględnie zbieżny
$\mu$-prawie wszędzie. Jeśli $F$ oznacza jego sumę
na zbiorze zbieżności i zero na pozostałym zbiorze,
to $F$ jest całkowalna, a dla sum częściowych $S_N$
zachodzi oszacowanie ogona
\[
 \int_X|F-S_N|\,\mathrm d\mu\le
 \sum_{n>N}\int_X|f_n|\,\mathrm d\mu\longrightarrow0.
\]
Ponadto szereg liczbowy $\sum_n\int_X f_n\,\mathrm d\mu$
jest bezwzględnie zbieżny.
W szczególności
\begin{equation}\label{eq:interchange-series}
\int_XF(x)\,\mathrm d\mu(x)
=\sum_{n=1}^{\infty}\int_Xf_n(x)\,\mathrm d\mu(x).
\end{equation}
\end{twierdzenie}
\begin{proof}
Zob. \atlasref{ME-08::thm:sum-absolutely-integrable}{Teoria miary, rozdział 8},
gdzie wykazano również oszacowanie ogona
i bezwzględną zbieżność szeregu całek.
\end{proof}


\subsection{Którego twierdzenia użyć?}
Zbieżność monotoniczna pozwala obliczać całkę funkcji
nieujemnej za pomocą rosnących obcięć i dopuszcza
nieskończoną wartość całki. Lemat Fatou daje jednostronne
oszacowanie bez monotoniczności. Zbieżność zmajoryzowana
pozwala zamienić granicę z całką i uzyskać zbieżność w $L^1$,
jeżeli istnieje wspólna całkowalna majoranta.
Na przestrzeni o skończonej mierze wystarczy wspólne
ograniczenie stałą. W przypadku szeregu stosujemy
zbieżność monotoniczną do nieujemnych sum częściowych,
a przy funkcjach ze znakiem sprawdzamy bezwzględną
sumowalność całek modułów.

\begin{przyklad}
Na $(0,1)$ z miarą Lebesgue'a $\lambda_1$ niech
$f_n(x)=\min(n,x^{-1/2})$. Funkcje są mierzalne,
nieujemne i $f_n\uparrow x^{-1/2}$. Rozdzielając
przedział w punkcie $n^{-2}$, otrzymujemy
\[
 \int_{(0,1)} f_n\,\mathrm d\lambda_1
 = n\cdot n^{-2}+\int_{n^{-2}}^1 x^{-1/2}\,\mathrm dx
 =\frac1n+2(1-\tfrac1n)=2-\frac1n.
\]
Drugi składnik obliczamy w sensie Riemanna na zwartym
przedziale, korzystając ze
\atlasref{ME-08::thm:riemann-to-lebesgue}{zgodności całek Riemanna i Lebesgue'a (ME-08)}.
Dla $n=1$ drugi składnik wynosi zero.
Twierdzenie \ref{thm:mct} daje więc
$\int_{(0,1)}x^{-1/2}\,\mathrm d\lambda_1(x)=2$.
Jeżeli natomiast $g_n(x)=x^n$ na $(0,1)$, to $g_n\to0$
punktowo i $0\le g_n\le1$. Stała majoranta $1$ jest
całkowalna, ponieważ $\lambda_1((0,1))=1$.
Z twierdzenia \ref{thm:dct} całki dążą zatem do zera.
Sprawdzamy to bezpośrednio: po ciągłym przedłużeniu
$g_n$ na $[0,1]$ korzystamy ze zgodności całek
Riemanna i Lebesgue'a (końce mają miarę zero), skąd
\[
 \int_{(0,1)}x^n\,\mathrm d\lambda_1(x)
 =\left[\frac{x^{n+1}}{n+1}\right]_0^1
 =\frac1{n+1}\longrightarrow0.
\]
\end{przyklad}

Dalsze warianty twierdzeń granicznych, w tym jednostajną
całkowalność i twierdzenie Vitaliego, przedstawiono w
\atlasref{ME-08::thm:vitali-l1}{ME-08: twierdzenie Vitaliego}.

% CAL-12: materiał kanoniczny znajduje się w ME-32.
% Wszystkie 23 etykiety oryginalnego rozdziału zachowano w ME-32.
% Do czasu wycofania całej książki pozostawiamy jedynie odsyłacz.
\providecommand{\atlasref}[2]{#2}
\section{Całka Daniella--Stone'a}

Pełną konstrukcję całki Daniella--Stone'a przeniesiono do
\atlasref{ME-32::ch:daniell-stone-32}{Teorii miary, rozdział 32}.
W odróżnieniu od całek uogólniających całkę Riemanna
konstrukcja ta wychodzi od dodatniego funkcjonału liniowego
na kracie funkcji i prowadzi do miary oraz całki Lebesgue'a.
Nie powtarzamy tu definicji, twierdzeń ani dowodów.

W rozdziale docelowym znajdują się:
\begin{itemize}
 \item \atlasref{ME-32::it:def:daniell-functional}{funkcjonał Daniella i ciągłość porządkowa};
 \item \atlasref{ME-32::it:thm:daniell-stone}{konstrukcja i jednoznaczność miary skończonej};
 \item \atlasref{ME-32::it:thm:daniell-unbounded-unit}{wersja dla funkcji nieograniczonych z jedynką};
 \item \atlasref{ME-32::it:thm:daniell-stone-general}{wersja ogólna z warunkiem Stone'a i lokalna jednoznaczność};
 \item \atlasref{ME-32::it:cor:daniell-general-density}{gęstość kraty i uzupełnienie miary};
 \item \atlasref{ME-32::it:thm:daniell-radon-comparison}{porównanie z miarą Radona}.
\end{itemize}

% CAL-13: materiał kanoniczny w ME-20, ME-21, ME-32 i FA-12.
% Wszystkie 44 etykiety oryginalnego rozdziału zachowano w tych źródłach.
% Do wycofania książki pozostawiamy jedynie odsyłacze.
\providecommand{\atlasref}[2]{#2}
\section{Całka Radona}

Topologiczne własności miar i reprezentację dodatnich
funkcjonałów przez całki omawiamy w teorii miary.
Nie powtarzamy tutaj definicji, dowodów ani przykładów.

\begin{itemize}
 \item \atlasref{ME-20::def:radon-measure-20}{Miary Radona i regularność};
 \item \atlasref{ME-20::lem:cutoff-radon-20}{Funkcje odcinające o zwartym nośniku};
 \item \atlasref{ME-20::thm:uniqueness-cc-20}{Jednoznaczność miary Radona};
 \item \atlasref{ME-21::thm:riesz-positive-21}{Twierdzenie Riesza--Markowa--Kakutaniego dla lokalnie zwartych przestrzeni Hausdorffa};
 \item \atlasref{ME-21::thm:compact-riesz-21}{Alternatywna konstrukcja na zwartych przestrzeniach metrycznych};
 \item \atlasref{ME-21::eq:gluing-measures-21}{Sklejanie miar na wyczerpaniu};
 \item \atlasref{ME-21::thm:riesz-c0-positive-21}{Dodatnie funkcjonały na przestrzeni $C_0(X)$};
 \item \atlasref{ME-21::thm:riesz-complex-21}{Reprezentacja funkcjonałów zespolonych};
 \item \atlasref{FA-12::prop:ck-banach12}{Zupełność przestrzeni $C(K)$};
 \item \atlasref{ME-32::it:thm:daniell-radon-comparison}{Porównanie miar Daniella i Radona};
 \item \atlasref{ME-32::it:cor:radon-cc-l1-general}{Gęstość funkcji o zwartym nośniku w $L^1$}.
\end{itemize}
Przykłady miar atomowych, miar skupionych na okręgu
oraz sum miar atomowych i miary Lebesgue'a znajdują się w
\atlasref{ME-21::ch:riesz-markov-21}{rozdziale 21 teorii miary}.

% CAL-14: materiał kanoniczny w ME-23.
% Wszystkie 37 etykiet oryginalnego rozdziału zachowano w ME-23.
% Do wycofania książki pozostawiamy jedynie odsyłacze.
\providecommand{\atlasref}[2]{#2}
\section{Całka na grupie topologicznej: miara Haara}

Pełny wykład miary Haara, wraz z dowodami i przykładami,
znajduje się w \atlasref{ME-23::ch:haar-23}{rozdziale 23 teorii miary}.
Nie powtarzamy tu definicji ani konstrukcji.

\begin{itemize}
 \item \atlasref{ME-23::def:haar-23}{Lewe i prawe miary Haara};
 \item \atlasref{ME-23::thm:haar-main-23}{Istnienie i jednoznaczność bez założeń przeliczalności};
 \item \atlasref{ME-23::cor:haar-finiteness-23}{Skończoność i $\sigma$-skończoność};
 \item \atlasref{ME-23::thm:haar-modular-23}{Funkcja modularna};
 \item \atlasref{ME-23::cor:haar-integral-inversion-23}{Wzory całkowe i inwersja};
 \item \atlasref{ME-23::it:lem:haar-lie-volume}{Miara Haara na grupach Liego};
 \item \atlasref{ME-23::it:thm:haar-l1-convolution}{Algebra splotu $L^1(G)$};
 \item \atlasref{ME-23::it:ex:haar-ell2}{Znaczenie lokalnej zwartości};
 \item \atlasref{ME-23::tw:haar-automorfizm}{Moduł automorfizmu}.
\end{itemize}
Przykłady grup zwartych, dyskretnych, torusa,
grupy afinicznej i całek względem miary Haara
są omówione w tym samym rozdziale.

\section{Całka Mawhina}
Wielowymiarowe całki cechowane wymagają uwzględnienia
kształtu komórek podziału: dla bardzo cienkich
prostokątów stosunek miary brzegu do objętości
może być dowolnie duży.
W niniejszym skrypcie nie opracowano definicji ani twierdzeń
dotyczących całki Mawhina. Jej przyszłe opracowanie
wstrzymano do osobnej decyzji o zakresie i miejscu.

\providecommand{\atlasref}[2]{#2}
\section*{Literatura i status dowodów}
\addcontentsline{toc}{section}{Literatura i status dowodów}
\begin{enumerate}[label={[\arabic*]},leftmargin=*]
\item S. Saks, \emph{Theory of the Integral}, 2. wyd., Warszawa 1937; rozdziały o całkach Denjoya i Perrona.
\item R. A. Gordon, \emph{The Integrals of Lebesgue, Denjoy, Perron, and Henstock}, Graduate Studies in Mathematics 4, American Mathematical Society, 1994.
\item A. Fonda, \emph{The Kurzweil--Henstock Integral for Undergraduates}, Birkh\"{a}user, 2018.
\item M. Tvrd\'{y}, \emph{Kurzweil--Stieltjes Integration}, \url{https://users.math.cas.cz/~tvrdy/ks-dekk.pdf}.
\item Y. Kubota, \emph{A constructive definition of the approximately continuous Denjoy integral}, Canad. Math. Bull. 15 (1972), 103--109.
\item Materiały Biblioteki matematycznej: Analiza matematyczna I (rozdziały 19, 29, 36), Teoria miary (rozdziały 06--08, 14, 17--18, 20--21, 23), Analiza funkcjonalna (rozdział 12). Wykorzystano opracowane tam dowody w odpowiednich częściach niniejszego skryptu.
\end{enumerate}

\paragraph*{Zakres samodzielności dowodów.}
Bezpośrednio wykazano twierdzenia o sumach całkowych, wybrane kryteria, rachunki z funkcjami skokowymi, lemat uzupełniania podziałów, proste twierdzenia Perrona oraz własności wynikające z funkcji pierwotnych. Charakteryzację wszystkich pierwotnych KH przez $\AC G_*$ udowodniono w twierdzeniu \atlasref{AM1-38::it:thm:kh-acgstar-characterization}{charakteryzacja funkcji pierwotnych KH}, z wykorzystaniem wskazanych wyników teorii miary i całki Lebesgue'a. Twierdzenie Riesza--Markowa--Kakutaniego (twierdzenie \atlasref{ME-21::thm:riesz-positive-21}{ME-21: twierdzenie Riesza--Markowa--Kakutaniego}) udowodniono dla dowolnych lokalnie zwartych przestrzeni Hausdorffa przez zawartość zbiorów otwartych, miarę zewnętrzną i kryterium Carathéodory'ego; zachowano również wcześniejszy dowód metryczny. Dowody istnienia i jednoznaczności miary Haara bez założeń przeliczalności znajdują się w \atlasref{ME-23::thm:haar-main-23}{ME-23: twierdzenie Haara}. Wykorzystują wcześniejsze wyniki topologii i teorii miary, w szczególności twierdzenie Carathéodory'ego, przywoływane w tekście. Równoważność całek Perrona i HK udowodniono w twierdzeniu \atlasref{AM1-39::it:thm:perron-hk}{równoważność całek Perrona i KH}, z wykorzystaniem małej ciągłej korekty na zbiorze wyjątkowym (lemat \atlasref{AM1-39::it:lem:perron-null-correction}{korekta na zbiorze wyjątkowym}). Istnienie i jednoznaczność miary w skończonej wersji twierdzenia Daniella--Stone'a udowodniono w \atlasref{ME-32::it:thm:daniell-stone}{ME-32: twierdzenie Daniella--Stone'a} przez konstrukcję miary zewnętrznej, mierzalność zbiorów poziomicowych i aproksymację kratową, wykorzystując \atlasref{ME-03::thm:caratheodory}{ME-03: twierdzenie Carathéodory'ego}. Tam też znajduje się \atlasref{ME-32::it:thm:daniell-stone-general}{ME-32: wersja bez jedynki i bez założenia skończoności miary}. Wykaz nie zastępuje szczegółowej kontroli zależności dowodów; migracja nie oznacza zakończenia audytów matematycznych.


\end{document}
